\chapter*{圆锥曲线}

\section*{圆的标准方程和圆的一般方程}


A 组：

\begin{enumerate}

    \item 当\blkda[4]{$D^2+E^2-4F>0$}时,二元二次方程 ${x}^{2} + {y}^{2} + {Dx} + {Ey} + F = 0$ 叫做圆的一般方程. 其中圆心为\blkda[4]{$(-\dfrac{D}{2},-\dfrac{E}{2})$},圆的半径为 $r =$ \blkda[4]{$\dfrac{\sqrt{D^2+E^2-4F}}{2}$}.

    \item 某圆的圆心为(2, - 1),且过点(4, - 2),则圆的标准方程是\blkda[4]{$(x-2)^2+(y+1)^2=5$}.
    \begin{jieda}
        由题意半径 $r=\sqrt{(4-2)^2+(-2+1)^2}=\sqrt{5}$ ，所以圆的标准方程为 $(x-2)^2+(y+1)^2=5$.
    \end{jieda}

    \item 已知点 $P(1,-5)$ ,则该点与圆 ${x}^{2} + {y}^{2} = {25}$ 的位置关系是\blkda[3]{在圆的外部}
    \begin{jieda}
        由圆 $x^2+y^2=25$ 的圆心 $(0,0)$ 到点 $P(1,-5)$ 的距离为 $d=\sqrt{1^2+(-5)^2}=\sqrt{26}>5=r$ ，知点 $P(1,-5)$ 在圆的外部．
    \end{jieda}

    \item 若圆 ${\left( x - 1\right) }^{2} + {\left( y - 2\right) }^{2} = 5$ 的圆心到直线 $x - y + a = 0$ 的距离为 $\sqrt{2}$ ,则实数 $a$ 的值为\blkda[3]{-1或3}

    \item 若圆 $C : {x}^{2} + {y}^{2} + {2mx} - {4y} - 1 = 0$ 关于直线 $y = {3x} + 1$ 对称,则实数 $m =$ \blkda[3]{$-\dfrac{1}{3}$}
    \begin{jieda}
        由题知，直线 $y=3 x+1$ 过圆 $C$ 的圆心 $(-m, 2)$ ，则 $-3 m+1=2$ ，解得 $m=-\dfrac{1}{3}$ ．
    \end{jieda}
    
    \item “ $A = C \neq  0,B = 0$ ” 是 “方程 $A{x}^{2} + {Bxy} + C{y}^{2} + {Dx} + {Ey} + F = 0$ 表示圆” 的\blkda[3]{必要非充分}条件.

    \item 已知一个圆的直径端点是 $A\left( {{x}_{1},{y}_{1}}\right)$ 、 $B\left( {{x}_{2},{y}_{2}}\right)$ ,求这个圆的标准方程.
    \begin{jieda}
        易知圆心为$\left(\dfrac{x_1+x_2}{2}, \dfrac{y_1+y_2}{2}\right)$，半径为$\sqrt{\dfrac{\left(x_1-x_2\right)^2+\left(y_1-y_2\right)^2}{4}}$
    
        $\left(x-\dfrac{x_1+x_2}{2}\right)^2+\left(y-\dfrac{y_1+y_2}{2}\right)^2=\dfrac{\left(x_1-x_2\right)^2+\left(y_1-y_2\right)^2}{4}$
    \end{jieda}

    \item 求经过 $A\left( {3,2}\right)$ 、 $B\left( {1,1}\right)$ 、 $C\left( {2, - 1}\right)$ 三点的圆的方程.
    \begin{jieda}
        设过 $A(3,2), B(1,1), C(2,-1)$ 三点的圆的方程为：$x^2+y^2+D x+E y+F=0, D^2+E^2-4 F>0$，则 $\left\{\begin{array}{l}13+3 D+2 E+F=0 \\ 2+D+E+F=0 \\ 5+2 D-E+F=0\end{array}\right.$ ，解得 $\left\{\begin{array}{l}D=-5 \\ E=-1, \\ F=4\end{array}\right.$
        
        $\therefore$ 所求圆的方程为 $x^2+y^2-5 x-y+4=0$ ．
    \end{jieda}
    
\end{enumerate}


B 组
\begin{enumerate}

    \item 以点(-3,2)为圆心,且与 $x$ 轴相切的圆的标准方程是\blkda[4]{$(x+3)^2+(y-2)^2=4$}.

    \item 过 $A\left( {2, - 3}\right)$ 、 $B\left( {-2, - 5}\right)$ 两点,且圆心在直线 $x - {2y} - 3 = 0$ 上的圆的标准方程是\blkda[4]{$(x+1)^2+(y+2)^2=10$}.

    \item 若 $\bigtriangleup {ABO}$ 的顶点坐标分别为 $A\left( {4,0}\right) \text{、}B\left( {0,3}\right) \text{、}O\left( {0,0}\right)$ ,则 $\bigtriangleup {ABO}$ 的内切圆的标准方程为\blkda[4]{$(x-1)^2+(y-1)^2=1$}.
    \begin{jieda}
        $\bigtriangleup {ABO}$ 是直角三角形,斜边长为5, 内切圆半径为 $r = \dfrac{3 + 4 - 5}{2} = 1$
        
        故内切圆圆心坐标为$(1,1)$,其方程为 ${\left( x - 1\right) }^{2} + {\left( y - 1\right) }^{2} = 1$ .
    \end{jieda}

    \item 若圆 $C$ 与圆 ${\left( x + 2\right) }^{2} + {\left( y - 1\right) }^{2} = 1$ 关于原点对称,则圆 $C$ 的标准方程为\blkda[4]{$(x-2)^2+(y+1)^2=1$}.
    \begin{jieda}
        两圆圆心关于原点对称，半径相同.
    \end{jieda}

    \item 过圆 ${\left( x + 2\right) }^{2} + {y}^{2} = 4$ 的圆心且与直线 $x + y = 0$ 垂直的直线方程为\blkda[4]{ $x-y+2=0$ }.
    \begin{jieda}
        圆 $(x+2)^2+y^2=4$ 的圆心为 $(-2,0)$ ，与直线 $x+y=0$ 垂直的直线的斜率为 1 ，所以所求直线为 $y-0=1 \cdot(x+2)$ ，即 $x-y+2=0$ ．
    \end{jieda}
    
    \item 圆 $C : {x}^{2} + {y}^{2} - {2x} - {2y} + 1 = 0$ 上的动点 $Q$ 到直线 $l : {3x} + {4y} + 8 = 0$ 的距离的最小值是\blkda[2]{2}.
    \begin{jieda}
        圆心$(1,1)$到直线 $l : {3x} + {4y} + 8 = 0$ 的距离为$\dfrac{|3+4+8|}{\sqrt{9+16}} = 3$，因此圆上动点$Q$到直线 $l : {3x} + {4y} + 8 = 0$ 的距离的最小值是2.
    \end{jieda}

    \item 圆 $C : {x}^{2} + {y}^{2} - {2ax} + {2ay} + 3{a}^{2} - {2a} - 1 = 0$ 的面积的最大值是\blkda[2]{$2\pi$}.
    \begin{jieda}
        ${x}^{2} + {y}^{2} - {2ax} + {2ay} + 3{a}^{2} - {2a} - 1 = 0$即$(x-a)^2+(y-a)^2 = 2-(a-1)^2$

        因此半径的最大值为$\sqrt{2}$，所以面积最大值为$2\pi$
    \end{jieda}

    \item 已知坐标原点在圆 ${x}^{2} + {y}^{2} + x + {3y} - a = 0$ 的外部,则实数 $a$ 的取值范围为\blkda[3]{$\left(-\dfrac{5}{2},0\right)$}.
    \begin{jieda}
        将圆的一般方程化为标准方程得到$\left(x + \dfrac{1}{2}\right)^2 + \left(y + \dfrac{3}{2}\right)^2 = a + \dfrac{5}{2}$，故当原点在圆外时，应有$\left(0 + \dfrac{1}{2}\right)^2 + \left(0 + \dfrac{3}{2}\right)^2 > a + \dfrac{5}{2} > 0$，因此$a \in \left(-\dfrac{5}{2},0\right)$
    \end{jieda}

    \item 已知 $A\left( {-\sqrt{3},0}\right) ,B\left( {\sqrt{3},0}\right) ,C\left( {0,3}\right)$ ,则 $\bigtriangleup  {{ABC}}$ 外接圆的方程为\blkda[4]{$x^2+(y-1)^2=4$}
    \begin{jieda}
        边$AB$的中垂线为$y$轴，设圆心为$(0, t), t > 0$，则$(3-t)^2 = t^2 + 3$，解得$t = 1$，故外接圆的圆心为$(0, 1)$，半径为$2$，外接圆方程为$x^2+(y-1)^2=4$.
    \end{jieda}

    \item 曲线 ${x}^{2} + {y}^{2} = 4\left| x\right|  + 4\left| y\right|$ 的周长为\blkda[3]{$8\sqrt{2}\pi$}
    \begin{jieda}
        在第一象限及$x$轴,$y$轴正方向上，方程可化为${x}^{2} + {y}^{2} = 4x+4y$，即$(x-2)^2+(y-2)^2 = 8$，计算可知该段周长为$2\sqrt{2}\pi$

        若$(x_0, y_0)$在曲线上，则$(x_0, -y_0), (-x_0, y_0), (-x_0, -y_0)$均在曲线上，因此曲线关于$x$轴，$y$轴对称，且关于原点中心对称，故周长为$8\sqrt{2}\pi$
    \end{jieda}

    \item 若圆 ${x}^{2} + {y}^{2} + {ax} + {by} + c = 0$ 关于直线 $x - {2y} = 0$ 对称,则
    \dotfill{(\kh{B})}
    \xx{$a = {2c}$}{$a = {2b}$}{$a + {2b} = 0$}{ $a = b$}
    \begin{jieda}
        圆心$\left(-\dfrac{a}{2}, -\dfrac{b}{2}\right)$在直线上.
    \end{jieda}

    \item 当 $a \neq  0$ 时,方程 ${x}^{2} + {y}^{2} + {2ax} - {2ay} = 0$ 表示的图形
    \dotfill{(\kh{C})}
    \xx{关于 $x$ 轴对称}{关于 $y$ 轴对称}{关于直线 $x + y = 0$ 对称}{关于直线 $x - y = 0$ 对称}
    \begin{jieda}
        ${x}^{2} + {y}^{2} + {2ax} - {2ay} = 0 \Longleftrightarrow (x+a)^2+(y-a)^2 = 2a^2$，圆心为$(-a, a)$，在直线$x + y = 0$上.
    \end{jieda}


    \item 已知两点 ${A}\left( {-1,0}\right) ,{B}\left( {1,0}\right)$ ,且点 ${P}$ 在直线 $x\cos \theta  + y\sin \theta  + 1 = 0$ 上,当 $\angle {APB} = {90}^{ \circ  }$ 时,则点 ${P}$ 有
    \dotfill{(\kh{D})}
    \xx{0个}{1个}{2个}{个数与 $\theta$ 有关}
    \begin{jieda}
        问题转化为 以$AB$为直径的圆$x^2+y^2 = 1$与直线$x\cos \theta  + y\sin \theta  + 1 = 0$的公共点个数

        圆心$(0, 0)$到直线$x\cos \theta  + y\sin \theta  + 1 = 0$的距离$d = \dfrac{1}{\sqrt{sin^2\theta + cos^2\theta}} = 1$，故$x^2+y^2 = 1$与直线$x\cos \theta  + y\sin \theta  + 1 = 0$相切，只有一个公共点.

        本题条件为$\angle {APB} = {90}^{ \circ  }$，因此轨迹会挖掉$A、B$两点，故为D.
    \end{jieda}


    \item 已知圆 $C : {x}^{2} + {y}^{2} + {2ax} + {2y} + {\left( a - 1\right) }^{2} = 0$ . 若 $0 < a < 1$ ,则坐标系原点 $O$
    \dotfill{(\kh{B})}
    \xx{在圆上}{在圆外}{在圆内}{以上关系都可能}
    \begin{jieda}
        圆的标准方程为$(x+a)^2+(y+1)^2 = 2a$，因为当$a \in (0, 1)$时，$(0+a)^2+(0+1)^2 > 2a$，所以原点在圆外.
    \end{jieda}

    \item 求圆 $C : {x}^{2} + {y}^{2} + {4x} - {12y} + {39} = 0$ 关于直线 $l : {3x} - {4y} + 5 = 0$ 对称的圆 ${C}^{\prime }$ 的标准方程.
    \begin{jieda}
        $C(-2, 6), r = 1$
        
        设$C'(a, b)$，则有$\left\{\begin{array}{l}
             3 \times \dfrac{a-2}{2} - 4 \times \dfrac{b+6}{2} + 5 = 0 \\
             (a+2) \times (-4) = (b-6) \times 3
        \end{array} \right.$，解得$C' (4, -2)$
        
        故圆$C' : (x-4)^2+(y+2)^2=1$
    \end{jieda}

    \item 讨论方程 ${x}^{2} + {y}^{2} + {ax} + {2y} + 2 = 0$ 所表示的曲线.
    \begin{jieda}
        当 $a \in(-\infty,-2) \cup(2,+\infty)$ ，曲线为圆心为 $\left(-\dfrac{a}{2},-1\right)$ ，半径为 $\sqrt{\dfrac{a^2}{4}-1}$ 的圆；当 $a=2$ 时，点（ $-1,-1$ ）；当 $a=-2$ 时，点 $(1,-1)$ ；当 $a \in(-2,2)$ 时，曲线不存在。
    \end{jieda}

    \item 已知点 $A\left( {0,1}\right)$ 、 $B\left( {4,t}\right)$ . 是否存在实数 $t$ ,满足过 $A$ 、 $B$ 两点作与 $x$ 轴相切的圆有且只有一个? 若存在满足条件的圆, 求出这个圆的方程; 若不存在满足条件的圆, 说明理由.
    \begin{jieda}
        设圆心为$(a, b)$，根据题意可知半径为$b$，故圆的标准方程为$(x-a)^2+(y-b)^2 = b^2$
        
        代入点$A, B$的坐标可知$\left\{\begin{array}{l}
             a^2+(1-b)^2 = b^2 \\
             (4-a)^2+(t-b)^2 = b^2
        \end{array}\right.$

        上式化简可得$a^2+1=2b$，两式做差可得$4(4-2a)+(t+1-2b)(t-1) = 0$，消元得$4(4-2a)+(t-a^2)(t-1) = 0$，整理成$a$为主元的形式：$(1-t)a^2 - 8a + 16+t^2-t = 0$

        存在实数 $t$ ,满足过 $A$ 、 $B$ 两点作与 $x$ 轴相切的圆有且只有一个，即存在实数 $t$使得关于$a$的方程$(1-t)a^2 - 8a + 16+t^2-t = 0$只有一个实根，即$t = 1$或$\Delta = 64 - 4(1-t)(16+t^2-t) = 0 (t \neq 1)$

        $t \neq 1$时，$64 - 4(1-t)(16+t^2-t) = 0$即$\dfrac{16}{1-t} = t^2-t+16$，结合图像可知$y = \dfrac{16}{1-t}$与$y = t^2-t+16$只有一个交点$(0, 16)$.

        故$t \in \{0, 1\}$时满足要求.

        \begin{dingautolist}{172}
            \item $t=1,(x-2)^2+\left(y-\dfrac{5}{2}\right)^2=\dfrac{25}{4}$
            \item $t=0,(x-4)^2+\left(y-\dfrac{17}{2}\right)^2=\dfrac{289}{4}$
        \end{dingautolist}
    \end{jieda}
    

    \item 已知直线 ${l}_{1} : {2x} - y + a = 0$ ,直线 ${l}_{2} : {4x} - {2y} - 1 = 0$ 和直线 ${l}_{3} : x + y - 1 = 0$ .
    \begin{enumerate}
        \item 若直线 ${l}_{1}$ 与 ${l}_{2}$ 的距离为 $\dfrac{\sqrt{5}}{2}$ ,求实数 $a$ 的值;
        \item 曲线 ${x}^{2} + {y}^{2} = {10}\left( {x > 0,y > 0}\right)$ 上是否存在点 $P$ ,使得 $P$ 到直线 ${l}_{2}$ 的距离与 $P$ 到直线 ${l}_{3}$ 距离之比是 $3 : \sqrt{10}$ ,若存在,求出所有满足条件的 $P$ 点坐标,若不存在, 说明理由.
    \end{enumerate}
    \begin{jieda}
        \begin{enumerate}
            \item 将 $l_2$ 的方程化为： $2 x-y-\dfrac{1}{2}=0$ ，$\therefore$ 直线 $l_1$ 与 $l_2$ 的距离 $d=\dfrac{\left|a+\dfrac{1}{2}\right|}{\sqrt{2^2+(-1)^2}}=\dfrac{\sqrt{5}}{2}$，解得 $a=-3$ 或 $a=2$ ；

            \item 设 $P\left(x_0, y_0\right)$ ，则 $P$ 点到 $l_2$ 距离为 $d_1=\dfrac{\left|4 x_0-2 y_0-1\right|}{\sqrt{4^2+(-2)^2}}$ ，则 $P$ 点到 $l_3$ 距离为 $d_2=\dfrac{\left|x_0+y_0-1\right|}{\sqrt{1^2+(-1)^2}}$ ，则 $\dfrac{\left|4 x_0-2 y_0-1\right|}{\sqrt{4^2+(-2)^2}} \times \sqrt{10}=\dfrac{\left|x_0+y_0-1\right|}{\sqrt{1^2+(-1)^2}} \times 3$ ，即 $\left|4 x_0-2 y_0-1\right|=\left|3 x_0+3 y_0-3\right|$ ，又因为 $x_0^2+y_0^2=10$ ，解得：$\left\{\begin{array}{l}x_0=3 \\ y_0=1\end{array}\right.$ 或 $\left\{\begin{array}{l}x_0=\dfrac{3}{25} \\ y_0=\dfrac{79}{25}\end{array}\right.$ ；
        \end{enumerate}

    \end{jieda}
    
\end{enumerate}

\clearpage
\section*{直线和圆的位置关系}

A 组：

\begin{enumerate}

    \item 已知圆 ${x}^{2} + {y}^{2} - {2x} + {6y} = 6$ ,则直线 ${3x} - {4y} + 1 = 0$ 和圆的位置关系为\blkda[2]{相交}.
    \begin{jieda}
        由圆 $x^2+y^2-2 x+6 y=6$ 得 $(x-1)^2+(y+3)^2=16$ ，圆心 $(1,-3)$ ，半径 $r=4$ ，圆心 $(1,-3)$ 到直线 $3 x-4 y+1=0$ 的距离 $d=\dfrac{|3 \times 1-4 \times(-3)+1|}{\sqrt{3^2+4^2}}=\dfrac{16}{5}<4$ ，所以直线 $3 x-4 y+1=0$ 和圆的位置关系为相交，
    \end{jieda}

    \item 若直线 $x - y + 1 = 0$ 与圆 ${x}^{2} + {y}^{2} - {2x} + 1 - a = 0$ 相切,则 $a$ 等于\blkda[2]{2}.
    \begin{jieda}
$\left\{\begin{array}{l}
x-y+1=0 \\
x^2+y^2-2 x+1-a=0
\end{array}\right.
$解得 $2 x^2+2-a=0$ ，因为直线与圆相切，所以 $-4 \times 2(2-a)=0$ ，解得 $a=2$ ．
    \end{jieda}

    \item 直线 ${x} + {y} = 0$ 被圆 ${x}^2 + {y}^2 - 6{x} + 2{y} + 4 = 0$ 截得的弦长等于\blkda[2]{4}.
    \begin{jieda}
        因为 $x^2+y^2-6 x+2 y+4=0$ ，所以 $(x-3)^2+(y+1)^2=6$ ，圆心 $(3,-1)$ 到直线 $x+y=0$ 的距离 $d=\dfrac{|3-1|}{\sqrt{2}}=\sqrt{2}$ ，直线 $x+y=0$ 被圆 $x^2+y^2-6 x+2 y+4=0$ 截得的弦长 $l =2 \sqrt{(\sqrt{6})^2-(\sqrt{2})^2}=4$ .
    \end{jieda}

    \item 过点 $A\left( {1, - 3}\right)$ 的圆 ${x}^{2} + {y}^{2} = {10}$ 的切线方程是 \blkda[4]{$x-3y-10=0$}.
    \begin{jieda}
        $A$在圆上，为切点，故切线方程为$x\times1+y\times3 = 10$，即$x-3y-10=0$.
    \end{jieda}

    \item 圆 ${x}^{2} + {y}^{2} = 9$ 的过点 $M\left( {3,4}\right)$ 的切线方程为\blkda[6]{$7x-24y+75=0$或$x=3$}.
    \begin{jieda}
        $M$在圆外
        \begin{dingautolist}{172}
            \item 当斜率不存在时，过$M$的直线为$x = 3$，满足要求
            \item 当斜率存在时，设切线为$y = k(x-3)+4$，故原点到切线的距离$d = \dfrac{|4-3k|}{\sqrt{1+k^2}} = 3$，解得$k = \dfrac{7}{24}$，故切线方程为$7x-24y+75=0$
        \end{dingautolist}
    \end{jieda}

\end{enumerate}

B 组:

\begin{enumerate}
    \item 若过点$(0,2)$且与圆 ${x}^{2} - {2x} + {y}^{2} - {2y} + m = 0$ 相切的直线只有一条,则 $m =$ \blkda[2]{0}.
    \begin{jieda}
        $(0,2)$在圆上.
    \end{jieda}

    \item 已知过点 $M\left( {2,1}\right)$ 的直线 $l$ 与圆 $C : {x}^{2} + {y}^{2} - {2x} - {6y} + 5 = 0$ 相切,则直线 $l$ 的方程为\blkda[3]{$x-2y=0$}.
    \begin{jieda}
        圆$C: (x-1)^2+(y-3)^2 = 5$，因此$M$在圆上，为切点，故该处切线方程为$(x-1)(2-1)+(y-3)(1-3) = 5$，即$x-2y=0$.
    \end{jieda}

    \item 直线 ${ax} - {by} = 0$ 与圆 ${x}^{2} + {y}^{2} - {ax} + {by} = 0\left( {{a}^{2} + {b}^{2} \neq  0}\right)$ 的位置关系是\blkda[3]{相切}.

    \item 若动圆 $O$ 与直线 ${4x} + {3y} - 7 = 0$ 相切,且半径为 4,则圆心 $O$ 的轨迹方程为\blkda[6]{$4x+3y-27=0$或$4x+3y+13=0$}.
    \begin{jieda}
        圆心轨迹应为两平行线，且与已知直线的距离等于圆的半径4.

        设圆心轨迹为$4x+3y+c=0$，$d = \dfrac{|c+7|}{5} = 4$，故$c = 13$或$-27$.
    \end{jieda}

    \item 直线 $l : {ax} + y - 2 = 0$ 与圆 $C : {\left( x - 1\right) }^{2} + {\left( y - 2\right) }^{2} = 1$ 的公共点的个数为\blkda[2]{2}.
    \begin{jieda}
        直线过定点$(0, 2)$，在圆上，又因为直线的斜率存在，故两者不相切，因此必然有两个公共点.
    \end{jieda}

    \item 由曲线 $y = \left| x\right|$ 和圆 ${x}^{2} + {y}^{2} = 4$ 可围成两个面积不等的封闭图形,其中较小的一个面积是\blkda[2]{$\pi$}.
    \begin{jieda}
        较小的一个为圆心角为直角的扇形，面积为圆的面积的$\dfrac{1}{4}$，即$\pi$
    \end{jieda}

    \item 直线 $\sqrt{3}x + y - 2\sqrt{3} = 0$ 截圆 ${x}^{2} + {y}^{2} = 4$ 所得的劣弧所对的圆心角为\blkda[3]{$\dfrac{\pi}{3}$}.

    \item 与圆 ${x}^{2} + {y}^{2} + {8x} + {6y} = 0$ 相切且在两坐标轴上的截距相等的直线方程为\blkda[11]{$x+y+7+5\sqrt{2}=0$或$x+y+7-5\sqrt{2}=0$或$4x+3y=0$}.
    \begin{jieda}
        ${x}^{2} + {y}^{2} + {8x} + {6y} = 0 \Longleftrightarrow (x+4)^2+(y+3)^2 = 25$.
        
        \begin{dingautolist}{172}
            \item 截距为$0$时，设直线方程为$y = kx$，则$d = \dfrac{|3k-4|}{\sqrt{1+k^2}} = 5$，解得$k = -\dfrac{4}{3}$，故直线为$4x+3y = 0$
            \item 截距不为$0$时，设直线方程为$\dfrac{x}{a} + \dfrac{y}{a} = 1$，即$x+y = a (a \neq 0)$，则$d = \dfrac{|-3-4-a|}{\sqrt{2}} = 5$，解得$a = -7 \pm 5\sqrt{2}$，故直线方程为$x+y+7 \pm 5\sqrt{2} = 0$
        \end{dingautolist}
    \end{jieda}

    \item 已知直线 ${l}_{1} : x - {8y} + {30} = 0$ 与直线 ${l}_{2} : x + {5y} - {22} = 0$ 的交点为点 $A$ ,直线 ${l}_{3}$ 过点 $A$ ,且直线 ${l}_{3}$ 被圆 ${x}^{2} + {y}^{2} - {2x} + {2y} - {14} = 0$ 所截得的线段的长为 $2\sqrt{3}$ ,直线 ${l}_{3}$ 的方程为\blkda[8]{$2 x-3 y+8=0$ 和 $3 x+2 y-14=0$ }.
    \begin{jieda}
        由 $\left\{\begin{array}{l}x-8 y+30=0, \\ x+5 y-22=0,\end{array}\right.$ 得 $\left\{\begin{array}{l}x=2, \\ y=4 .\end{array}\right.$ 
        
        易知斜率不存在时不满足要求，故设直线 $l_3$ 的方程为 $y-4=k(x-2)$．因为圆 $(x-1)^2+(y+1)^2 =16$ 的圆心为 $(1,-1)$ ，半径 $r=4$ ，所以圆心到直线 $l$ 的距离 $d=\dfrac{|5-k|}{\sqrt{1+k^2}}$ ．由题设，得 $\dfrac{k^2-10 k+25}{1+k^2}+ (\sqrt{3})^2=16$ ．解得 $k_1=\dfrac{2}{3}, k_2=-\dfrac{3}{2}$ ．因此，直线 $l_3$ 的方程为 $2 x-3 y+8=0$ 和 $3 x+2 y-14=0$ ．
    \end{jieda}

    
    \item 设圆 $O$ 上的点 $A\left( {2,3}\right)$ 关于直线 $x + {2y} = 0$ 的对称点仍在圆 $O$ 上,且直线$x - y + 1 = 0$ 被圆 $O$ 所截得的线段的长为 $2\sqrt{2}$ ,圆 $O$ 的方程为\blkda[9]{$(x-6)^2+(y+3)^2=52$ 或 $(x-14)^2+(y+ 7)^2=244$}.
    
    \begin{jieda}
        设圆的方程为 $(x-a)^2+(y-b)^2=r^2$ ，由题意，知直线 $x+2 y=0$ 过圆心，
$
\therefore a+2 b=0 .
$
又点 $A$ 在圆上，$\therefore(2-a)^2+(3-b)^2=r^2$ ．
$\because$ 直线 $x-y+1=0$ 被圆截得的弦长为 $2 \sqrt{2}$ ，

$
\therefore(\sqrt{2})^2+\left[\dfrac{|a-b+1|}{\sqrt{1^2+(-1)^2}}\right]^2=r^2 .
$


联立可得 $\left\{\begin{array}{l}a=6 \\ b=-3 \text { 或 } \\ r^2=52\end{array}\left\{\begin{array}{l}a=14 \\ b=-7 \\ r^2=244\end{array}\right.\right.$ ，
故所求方程为 $(x-6)^2+(y+3)^2=52$ 或 $(x-14)^2+(y+ 7)^2=244$ ．

    \end{jieda}

    \item 已知 $\odot  C : {\left( x - 2\right) }^{2} + {y}^{2} = 1$ ,过坐标原点 $O$ 作 $\odot  C$ 的两条切线,切点为 $A\text{、}B$ ,则四边形 ${OACB}$ 的面积为\blkda[2]{$\sqrt{3}$}.
    \begin{jieda}
        $S_{OACB} = 2S_{\triangle OAB} = \sqrt{3}$
    \end{jieda}

    \item 若关于 $x$ 的方程 ${kx} = \sqrt{-{x}^{2} + {4x} - 3} + 1$ 有两个不同的实数解,则实数 $k$ 的取值范围是\blkda[3]{$\left[1,\dfrac{4}{3}\right)$}.
    \begin{jieda}
        方程 $k x=\sqrt{-x^2+4 x-3}+1 \Longleftrightarrow k x-1=\sqrt{1-(x-2)^2}$ ，令 $y=k x-1, y=\sqrt{1-(x-2)^2}$ ，由 $y=\sqrt{1-(x-2)^2}$ ，得 $(x-2)^2+y^2=1$ ，因此 $y=\sqrt{1-(x-2)^2}$ 表示以点 $C(2,0)$ 为圆心， 1 为半径的上半圆（含端点）， $y=k x-1$ 表示纵截距为 -1 ，斜率为 $k$ 的直线 $l$ ，依题意，直线 $l$ 与半圆 $C$ 有两个交点，
        在同一坐标系内作出直线 $l$ 与半圆 $C$ 
        \begin{center}
            \includegraphics[width=0.25\linewidth]{bodyxb//src/zhixianheyuan1.png}
        \end{center}
        
        如图,当直线 $l$ 与半圆 $C$ 相切时，$\dfrac{|2 k-1|}{\sqrt{k^2+1}}=1$ ，显然 $k>0$ ，解得 $k=\dfrac{4}{3}$ ，当直线 $l$ 过点 $(1,0)$ 时，$k=1$ ，此时直线 $l$ 与半圆 $C$ 有两个交点，当 $1 \leq k<\dfrac{4}{3}$ 时，直线 $l$ 与半圆 $C$ 有两个交点，即方程 $k x=\sqrt{-x^2+4 x-3}+1$ 有两个不同的实数解，所以实数 $k$ 的取值范围是 $1 \leq k<\dfrac{4}{3}$ ．
    \end{jieda}
    

    \item 直线 $x\cos a + y\sin a = 2$ 与圆 ${x}^{2} + {y}^{2} = 4$ 的位置关系是
    \dotfill{(\kh{C})}
    \xx{相交}{相离}{相切}{以上各种关系都可能}

    \begin{jieda}
        圆心$(0, 0)$到直线的距离$d = \dfrac{2}{\sqrt{sin^2a + cos^2a}} = 2$
    \end{jieda}


    \item 已知 ${BC}$ 是圆 ${x}^{2} + {y}^{2} = {25}$ 的动弦. 若 $\left| {BC}\right|  = 6$ ,则 ${BC}$ 的中点的轨迹方程是
    \dotfill{(\kh{B})}
    \xx{${x}^{2} + {y}^{2} = 9$}
    {${x}^{2} + {y}^{2} = {16}$ }
    {${x}^{2} + {y}^{2} = 4$ }
    {$x + y = 4$}
    \begin{jieda}
        弦$BC$中点到圆心$(0, 0)$的距离为4.
    \end{jieda}

     

    \item  “ $a = 1$ ” 是直线 $\left( {a + 1}\right) x + \left( {a - 1}\right) y + {2a} = 0\left( {a \in  \mathbf{R}}\right)$ 与圆 ${x}^{2} + {y}^{2} = 4$ 相交的\blkda[4]{充分非必要} 条件
    \begin{jieda}
        当 $a=1$ 时，直线为 $2 x+2=0$ ，即 $x=-1$ ，此时直线和圆相交，充分性成立；
    当直线 $(a+1) x+(a-1) y+2 a=0(a \in R)$ 与圆 $x^2+y^2=4$ 相交时，
    圆心到直线的距离 $d=\dfrac{|2 a|}{\sqrt{(a+1)^2+(a-1)^2}}<2$ ，
    化简得 $a^2+2>0$ ，显然恒成立，故不能推出 $a=1$ ，必要性不成立．
    \end{jieda}

    \item 已知圆 $O$ 的方程为 ${x}^{2} + {y}^{2} = {r}^{2}$ ,点 $P\left( {a,b}\right) ,\left( {{ab} \neq  0}\right)$ 是圆 $O$ 内一点,设以 $P$ 为中点的弦所在的直线为 $m$ ,方程为 ${ax} + {by} = {r}^{2}$ 的直线为 $n$ ,则
    \dotfill{(\kh{B})}
    \xx{$m//n$ ,且 $n$ 与圆 $O$ 相交} 
    {$m//n$ ,且 $n$ 与圆 $O$ 相离}
    {$m\bot n$ ,且 $n$ 与圆 $O$ 相交}
    {$m\bot n$ ,且 $n$ 与圆 $O$ 相离}
    \begin{jieda}
        直线 $O P$ 的斜率为 $\dfrac{b}{a}$ ，由垂径定理可知，$m \perp O P$ ，所以，直线 $m$ 的方程为 $y-b=-\dfrac{a}{b}(x-a)$ ，即 $a x+b y=a^2+b^2$ ，
由于点 $P$ 是圆 $O$ 内一点，则 $a^2+b^2<r^2$ ，
又直线 $n$ 的方程为：$a x+b y=r^2$ ，
所以，$m / / n$ ．
圆心 $O$ 到直线 $n$ 的距离为 $\dfrac{r^2}{\sqrt{a^2+b^2}}>\dfrac{r^2}{r}=r$ ，因此，直线 $n$ 与圆 $O$ 相离．
    \end{jieda}

    \item 已知圆 $C : {x}^{2} + {y}^{2} - {6x} - {8y} + m = 0$ ,其中 $m \in  \mathbf{R}$ .
    \begin{enumerate}
        \item 如果圆 $C$ 与圆 ${x}^{2} + {y}^{2} = 4$ 外切,求 $m$ 的值;
        \item 如果直线 $l : x + y - 3 = 0$ 与圆 $C$ 相交所得的弦长为 $4\sqrt{5}$ ,求 $m$ 的值.
    \end{enumerate}
    \begin{jieda}
        \begin{enumerate}
            \item 圆 $C$ 的圆心为 $(3,4)$ ，半径 $r=\sqrt{25-m}$ （ $m<25$ ）。因为圆 $C$ 与圆 $x^2+y^2=1$ 外切，所以两圆的圆心距等于其半径和，即 $\sqrt{(3-0)^2+(4-0)^2}=2+\sqrt{25-m}$ ，解得 $m=16$ ；
            \item 圆 $C$ 的圆心到直线 $x+y-3=0$ 的距离 $d=\dfrac{|3+4-3|}{\sqrt{2}}=2 \sqrt{2}$ ．因为直线 $x+y-3=0$与圆 $C$ 相交所得的弦长为 $4 \sqrt{5}$ ，所以 $r^2=25- m=(2 \sqrt{2})^2+(2 \sqrt{5})^2$ ，解得 $m=-3$ 。
        \end{enumerate}
    \end{jieda}
    
    \item 已知圆 $C : {x}^{2} + {y}^{2} = 4$ ,求圆 $C$ 的一组斜率为 1 的平行弦的中点的轨迹方程.
    \begin{jieda}
        设弦的两端点为 $A\left(x_1, y_1\right) 、 B\left(x_2, y_2\right)$ ，中点为 $M(x, y)$ ，则 $x_1^2+y_1^2=4$（1）,$x_2^2+y_2^2=4$（2）．
由（1）－（2），可得
$\left(x_1-x_2\right)\left(x_1+x_2\right)+\left(y_1-y_2\right)\left(y_1+y_2\right)=0$ ．因为
$x_1 \neq x_2$ ，所以 $x_1+x_2+\dfrac{y_1-y_2}{x_1-x_2}\left(y_1+y_2\right)=0$（3）．又
因为 $x_1+x_2=2 x, y_1+y_2=2 y, \dfrac{y_1-y_2}{x_1-x_2}=1$ ，
所以由（3）式，可得 $y=-x$ ．由 $\left\{\begin{array}{l}y=-x, \\ x^2+y^2=4,\end{array}\right.$ 可得 $x= \pm \sqrt{2}$ ．因此，所求方程为 $y=-x$ ， $x \in(-\sqrt{2}, \sqrt{2})$.
    \end{jieda}

    
    \item 在平面直角坐标系 ${xOy}$ 中,已知两点 $S\left( {4,0}\right) ,T\left( {1,0}\right)$ ,动点 $P$ 满足 $\left| {PS}\right|  = 2\left| {PT}\right|$ , 设点 $P$ 的轨迹为 $C$ . 如图,动直线 $l$ 与曲线 $C$ 交于不同的两点 $A,B$ ( $A,B$ 均在 $x$ 轴上方),且 $\angle {ATO} + \angle {BTO} = {180}^{ \circ  }$ .
    \begin{center}
        \includegraphics[width=0.25\linewidth]{bodyxb//src/bx01ch2.2-20.png}
    \end{center}
    
    \begin{enumerate}
        \item 求曲线 $C$ 的方程；
        \item 当 ${A}$ 为曲线 $C$ 与 $y$ 轴正半轴的交点时,求直线 $l$ 的方程;
        \item 是否存在一个定点,使得直线 $l$ 始终经过此定点？若存在,求出定点的坐标；若不存在,请说明理由. 
    \end{enumerate}
    \begin{jieda}
        \begin{enumerate}
            \item 设 ${P}({x}, {y})$ ，由 $|P S|=2|P T|$ 得 $\sqrt{(x-4)^2+y^2}=2 \sqrt{(x-1)^2+y^2}$ ，化简得 $x^2+y^2=4$ ，则曲线 $C$ 的方程为 $x^2+y^2=4$ ；
            \item 由题意知 $A(0,2)$ ，设 ${B}\left({x}_2, {y}_2\right)$ ，
            依题意可知直线 $l$ 的斜率存在，设直线 $l$ 的方程为 $y=k x+2$ ，
            由 $\angle A T O+\angle B T O=180^{\circ}$ ，得 $k_{A T}+k_{B T}=0$ ，
            则 $\left\{\begin{array}{l}-2+\dfrac{y_2}{x_2-1}=0 \\ x_2^2+y_2^2=4\end{array}\right.$ ，所以 $\left\{\begin{array}{l}x_2=0 \\ y_2=-2\end{array}\right.$（舍去）或 $\left\{\begin{array}{l}x_2=\dfrac{8}{5} \\ y_2=\dfrac{6}{5}\end{array}\right.$ ，即 $B\left(\dfrac{8}{5}, \dfrac{6}{5}\right)$ ，
            则 $k=\dfrac{\dfrac{6}{5}-2}{\dfrac{8}{5}-0}=-\dfrac{1}{2}$ ，
            则直线 $l$ 方程为 $y=-\dfrac{1}{2} x+2$ ；
            
            
            \item 设直线 $l$ 方程为 $y=k x+b$ ，设 ${A}\left({x}_1, {y}_1\right), {B}\left({x}_2, {y}_2\right)$ ，
            
            $\because \angle A T O+\angle B T O=180^{\circ}, \therefore k_{A T}+k_{B T}=0$
            
            $\therefore \dfrac{y_1}{x_1-1}+\dfrac{y_2}{x_2-1}=0 \Longrightarrow y_1\left(x_2-1\right)+y_2\left(x_1-1\right)=0$
             
            $\therefore \left(k x_1+b\right)\left(x_2-1\right)+\left(k x_2+b\right)\left(x_1-1\right)=0$
            
            $\therefore 2 k x_1 x_2+(b-k)\left(x_1+x_2\right)-2 b=0$
            
            
            联立方程 $\left\{\begin{array}{l}x^2+y^2=4 \\ y=k x+b\end{array}\right.$ ，得 $\left(k^2+1\right) x^2+2 k b x+b^2-4=0$ ，因此$\left\{\begin{array}{l}
                 x_1+x_2=\dfrac{-2 k b}{k^2+1} \\
                 x_1 x_2=\dfrac{b^2-4}{k^2+1}
            \end{array}\right.$

            $\therefore 2 k \dfrac{b^2-4}{k^2+1}+(b-k)\left(\dfrac{-2 k b}{k^2+1}\right)-2 b=0$，即$2 k\left(b^2-4\right)-2 k b(b-k)-2 b\left(k^2+1\right)=0$

            故$b = -4k$，则直线 $l$ 始终经过此定点 $(4,0)$ ．
            
        \end{enumerate}
    \end{jieda}
    
\end{enumerate}



C 组:

\begin{enumerate}
    \item 在平面直角坐标系 ${xOy}$ 中,圆 $C$ 的方程为 ${\left( x - 2\right) }^{2} + {y}^{2} = 1$ ,点 $P$ 在直线 $l : x + y + 1 = 0$ 上,若过点 $P$ 存在直线 $m$ 与圆 $C$ 交于 ${A},{B}$ 两点,且点 ${A}$ 为 ${PB}$ 中点,则点 $P$ 的横坐标的取值范围是\blkda[3]{$[-1,2]$}.
    \begin{jieda}
        设点 $P\left(x_0,-x_0-1\right), B(2+\cos \theta, \sin \theta)$ ，则
由条件得 $A$ 点坐标为 $x=\dfrac{1}{2}\left(x_0+2+\cos \theta\right), y=\dfrac{1}{2}\left(\sin \theta-x_0-1\right)$ ，
从而 $\left[\dfrac{1}{2}\left(x_0+2+\cos \theta\right)-2\right]^2+\left[\dfrac{1}{2}\left(\sin \theta-x_0-1\right)\right]^2=1$ ，
整理得 $\left(x_0-2\right) \cos \theta-\left(x_0+1\right) \sin \theta+x_0^2-x_0+1=0$ ，
从而 $\sqrt{2 {x}_0{ }^2-2 {x}_0+5} \sin (\theta+\alpha)=-x_0{ }^2+x_0-1$ ，
于是由 $\sqrt{2 {x}_0{ }^2-2 {x}_0+5} \geqslant\left|-x_0{ }^2+x_0-1\right|$ ，解得 $-1 \leqslant x_0 \leqslant 2$ ．
故答案为：$[-1,2]$ ．
    \end{jieda}
    
    \item 在平面直角坐标系 ${xOy}$ 中,已知 $A\left( {-1, - 1}\right)$ ,圆 $O : {x}^{2} + {y}^{2} = 1$ ,在直线 ${AO}$ 上存在异于 ${A}$ 的定点 $Q$ ,使得对圆 $O$ 上任意一点 $P$ ,都有 $\dfrac{\left| PA\right| }{\left| PQ\right| } = \lambda \left( {\lambda \text{ 为常数 }}\right)$ ,则 $Q$ 的坐标为\blkda[3]{$(-\dfrac{1}{2},-\dfrac{1}{2})$}.
    \begin{jieda}
    设 $Q\left(x_0, y_0\right), P(x, y)$ ，
则 $|P A|=\sqrt{(x+1)^2+(y+1)^2},|P Q|=\sqrt{\left(x-x_0\right)^2+\left(y-y_0\right)^2}$ ．
若在直线 $A O$ 上存在异于 A 的定点 $Q$ ，使得对圆 $O$ 上任意一点 $P$ ，都有 $\dfrac{|P A|}{|P Q|}=\lambda$（ $\lambda$ 为常数）．
等价于 $\dfrac{\sqrt{(x+1)^2+(y+1)^2}}{\sqrt{\left(x-x_0\right)^2+\left(y-y_0\right)^2}}=\lambda$ 对圆 $O$ 上任意点 $P(x, y)$ 恒成立．
即 $(x+1)^2+(y+1)^2=\lambda^2\left(x-x_0\right)^2+\lambda^2\left(y-y_0\right)^2$ ，
整理得 $\left(1-\lambda^2\right)\left(x^2+y^2\right)+\left(2+2 \lambda^2 x_0\right) x+\left(2+2 \lambda^2 y_0\right) y+2-\lambda^2\left(x_0^2+y_0^2\right)=0$ ，
因为点 $Q$ 在直线 $A O$ 上，所以 $x_0=y_0$ ，
由于 $P$ 在圆 $O$ 上，所以 $x^2+y^2=1$ ，
故 $\left(2+2 \lambda^2 x_0\right)(x+y)+3-\lambda^2-2 \lambda^2 x_0^2=0$ 恒成立，
其中点 $P(x, y)$ 在圆 $O: x^2+y^2=1$ 上，令 $x+y=m$ ，则 $x+y-m=0$ ，
所以直线 $x+y-m=0$ 与圆有交点，
所以圆心到直线的脑离小于等于半径，即 $d=\dfrac{|m|}{\sqrt{1^2+1^2}} \leq 1$ ，船得 $-\sqrt{2} \leq m \leq \sqrt{2}$ ，即 $x+y \in[-\sqrt{2}, \sqrt{2}]$ ，
所以 $\left\{\begin{array}{l}2+2 \lambda^2 x_0=0 \\ 3-\lambda^2-2 \lambda^2 x_0^2=0\end{array}\right.$ ，显然 $\lambda \neq 0$ ，所以 $x_0=-\dfrac{1}{\lambda^2}$ ，
故 $3-\lambda^2-\dfrac{2}{\lambda^2}=0$ ．
因为 $\lambda>0$ ，解得 $\lambda=\sqrt{2}$ 或 $\lambda=1$ ．
当 $\lambda=1$ 时，$Q(-1,-1)$ ，此时 $Q, A$ 重合，舍去．
当 $\lambda=\sqrt{2}$ 时，$Q\left(-\dfrac{1}{2},-\dfrac{1}{2}\right)$ ，
综上，存在满足条件的定点 $Q\left(-\dfrac{1}{2},-\dfrac{1}{2}\right)$ ，此时 $\lambda=\sqrt{2}$ ．
    \end{jieda}

    
    \item 已知二次函数 $y = {x}^{2} + \left( {{2m} - 3}\right) x - 4 - {11m}\left( {m \in  \mathbf{R}}\right)$ 与 $x$ 轴交于 $A,B$ 两点,点 $C\left( {1,3}\right)$ , 圆 $G$ 过 $A,B,C$ 三点,存在一条定直线 $l$ 被圆 $G$ 截得的弦长为定值,则该定值为\blkda[3]{$\sqrt{13}$}.
    \begin{jieda}
        设圆 $G$ 的方程为 $x^2+y^2+D x+E y+F=0$ ，因为圆 $G$过 $A, B$ 两点，
令 $y=0$ ，得 $x^2+D x+F=0$ ，
又 $A, B$ 两点的横坐标满足方程 $x^2+(2 m-3) x-4- 11 {~m}=0$ ，
所以 $D=2 m-3, F=-4-11 m$ ，
所以圆G的方程为 $x^2+y^2+(2 m-3) x+E y-4-$
$11 {~m}=0$ ，
又 $C(1,3)$ 在圆 $G$ 上，
所以 $10+(2 m-3)+3 E-4-11 m=0$ ，解得 $E=3 m-$
1，
所以圆G的方程为 $x^2+y^2+(2 m-3) x+(3 m-1) y- 4-11 {~m}=0$ ，
即 $x^2+y^2-3 x-y-4+m(2 x+3 y-11)=0$ ，
令 $\left\{\begin{array}{l}x^2+y^2-3 x-y-4=0 \\ 2 x+3 y-11=0\end{array}\right.$ ，解得 $\left\{\begin{array}{l}x=1 \\ y=3\end{array}\right.$ 或 $\left\{\begin{array}{l}x=4 \\ y=1\end{array}\right.$ ，
即圆 $G$ 恒过点 $(1,3)$ 和 $(4,1)$ ，
又 $\sqrt{(1-4)^2+(3-1)^2}=\sqrt{13}$ ，所以该定值为 $\sqrt{13}$ ．
    \end{jieda}
    
\end{enumerate}


\section*{两圆的位置关系}

A 组：

\begin{enumerate}
    \item 圆 ${C}_{1} : {\left( x - 1\right) }^{2} + {y}^{2} = 1$ 与圆 ${C}_{2} : {\left( x - 4\right) }^{2} + {\left( y - 4\right) }^{2} = {18}$ 的位置关系为\blkda[2]{相交}.

    \item 圆 ${x}^{2} + {y}^{2} - {4y} = 0$ 与圆 ${x}^{2} + {y}^{2} - {16} = 0$ 的位置关系是\blkda[2]{内切}.

    \item 若圆 ${x}^{2} + {y}^{2} + {4x} - {4y} = 0$ 与圆 ${x}^{2} + {y}^{2} + {2x} - {12} = 0$ 相交于 $P\text{、}Q$ 两点,则 ${PQ}$ 所在的直线方程为\blkda[3]{$x-2y+6=0$}.

    \item 在满足 ${\left( x - 3\right) }^{2} + {\left( y - 3\right) }^{2} = 6$ 的所有实数对$(x,y)$中, $\dfrac{y}{x}$ 的最大值是\blkda[3]{$3+2\sqrt{2}$}.

    \item 由点 $P\left( {3,2}\right)$ 引圆 ${x}^{2} + {y}^{2} = 4$ 的两条切线 ${PA}\text{、}{PB},A\text{、}B$ 为切点,则直线 ${AB}$ 的方程是\blkda[4]{$3x+2y-4=0$}.
    
\end{enumerate}

B 组:

\begin{enumerate}
    \item 已知圆 ${C}_{1} : {x}^{2} + {y}^{2} = 4$ 与圆 ${C}_{2} : {\left( x - 3\right) }^{2} + {\left( y - 4\right) }^{2} = {r}^{2}\left( {r > 0}\right)$ 相交,则 $r$ 的取值范围为\blkda[3]{$(3,7)$}.
    \begin{jieda}
        因为圆 $C_1: x^2+y^2=4$ 的圆心 $(0,0)$ ，半径为 2 ，圆 $C_2:(x-3)^2+(y-4)^2=r^2(r>0)$ 的圆心为 $(3,4)$ ，半径为 $r$ ，则圆心距为 5 ，且两圆相交，则 $|r-2|<5<|r+2|$ ，解得 $3<r<7$ ．
    \end{jieda}
    
    \item 与圆 ${x}^{2} + {y}^{2} - {6x} - {8y} + {21} = 0$ 外切且圆心在原点的圆的标准方程为\blkda[4]{$x^2+y^2=9$}.
    \begin{jieda}
        因为 $x^2+y^2-6 x-8 y+21=0$ ，即 $(x-3)^2+(y-4)^2=4$ ，可知圆心 $C(3,4)$ ，半径 $r=2$ ，则 $|O C|=\sqrt{3^2+4^2}=5$ ，由题意可得圆 $O$ 的半径 $R=|O C|-r=3$ ，所以圆 $O$ 的标准方程为 $x^2+y^2=9$ ．
    \end{jieda}

    \item 已知圆 ${O}_{1} : {x}^{2} + {y}^{2} = 4$ ,圆 ${O}_{2} : {x}^{2} + {y}^{2} - {2mx} - {2my} - 4 = 0\left( {m \neq  0}\right)$ ,则同时与圆 ${O}_{1}$ 和圆 ${O}_{2}$ 相切的直线有\blkda[2]{2}条.
    \begin{jieda}
        由 $O_1: x^2+y^2=4$ ，得圆 $O_1(0,0)$ ，半径为 $r_1=2$ ，由 $O_2: x^2+y^2-2 m x-2 m y-4=0(m \neq 0)$ ，得 $O_2(m, m)$ ，半径为$r_2=\dfrac{1}{2} \sqrt{(-2 m)^2+(-2 m)^2-4 \times(-4)}=\sqrt{2 m^2+4}$，所以 $\left|O_1 O_2\right|=\sqrt{(m-0)^2+(m-0)^2}=\sqrt{2 m^2}>0$ ，$\left|r_2-r_1\right|=\sqrt{2 m^2+4}-2>0, \quad r_1+r_2=2+\sqrt{2 m^2+4}$，所以 $\left|r_2-r_1\right|<\left|O_1 O_2\right|<r_1+r_2$ ，所以圆 $O_1$ 与圆 $O_2$ 相交，所以圆 ${O}_1$ 与圆 ${O}_2$ 有两条公共的切线．

    \end{jieda}

    
    \item 写出一个与圆 ${x}^{2} + {y}^{2} = 1$ 外切,并与直线 $y = \dfrac{\sqrt{3}}{3}x$ 及 $y$ 轴都相切的圆的方程.
    \begin{jieda}
        设所求圆的方程为：$(x-a)^2+(y-b)^2=r^2$ ，因为与圆 $x^2+y^2=1$ 外切，所以 $\sqrt{a^2+b^2}=1+r$ ，又因为与直线 $y=\dfrac{\sqrt{3}}{3} x$ 及 $y$ 轴都相切，所以圆心在角分线 $y=\sqrt{3} x$ 上或 $y=-\dfrac{\sqrt{3}}{3} x$ 上
        \begin{dingautolist}{172}
            \item 当圆心在 $y=\sqrt{3} x$ 上，所以 $b=\sqrt{3} a, r=|a|$ 
            
            联立得： $3 a^2=2|a|+1$ ，解得：$\left\{\begin{array}{l}a=1 \\ b=\sqrt{3}\end{array}\right.$ 或 $\left\{\begin{array}{l}a=-1 \\ b=-\sqrt{3}\end{array}, \quad r=1\right.$
            
            所以求得圆的方程为：$(x-1)^2+(y-\sqrt{3})^2=1$ 或 $(x+1)^2+(y+\sqrt{3})^2=1$

            \item 当圆心在 $y=-\dfrac{\sqrt{3}}{3} x$ 上，所以 $b=-\dfrac{\sqrt{3}}{3} a, r=|a|$ 
            
            联立得：$\dfrac{1}{3} a^2=2|a|+1$ ，解得：$\left\{\begin{array}{l}a=3+2 \sqrt{3} \\ b=-\sqrt{3}-2\end{array}\right.$ 或 $\left\{\begin{array}{l}a=-3-2 \sqrt{3} \\ b=\sqrt{3}+2\end{array}, r=3+2 \sqrt{3}\right.$
            
            所以求得圆的方程为：$(x-2 \sqrt{3}-3)^2+(y+2+\sqrt{3})^2=21+12 \sqrt{3}$ 或 $(x+2 \sqrt{3}+3)^2+(y-2-\sqrt{3})^2=21+12 \sqrt{3}$

        \end{dingautolist}
        故答案为：$(x-1)^2+(y-\sqrt{3})^2=1$ 或 $(x+1)^2+(y+\sqrt{3})^2=1$ 或$(x-2 \sqrt{3}-3)^2+(y+2+\sqrt{3})^2=21+12 \sqrt{3}$ 或 $(x+2 \sqrt{3}+3)^2+(y-2-\sqrt{3})^2=21+12 \sqrt{3}($ 写出其中一个即可）
    \end{jieda}

    \item  设 $P\left( {x,y}\right)$ 为圆 $C : {x}^{2} + {\left( y - 1\right) }^{2} = 1$ 上任意一点,欲使不等式 $x + y + c \geq  0$ 恒成立,则实数 $c$ 的取值范围为\blkda[3]{$[\sqrt{2}-1,+\infty$)}.

    \item 圆 ${\left( x - 1\right) }^{2} + {\left( y - 1\right) }^{2} = 8$ 与圆 ${x}^{2} + {y}^{2} + {2x} + {2y} = 0$ 的公共弦之长为\blkda[3]{$\dfrac{\sqrt{30}}{2}$}.
    \begin{jieda}
        ${\left( x - 1\right) }^{2} + {\left( y - 1\right) }^{2} = 8 \Longleftrightarrow x^2+y^2-2x-2y-6=0$
        
        公共弦$2x+2y+3=0$
        
        ${x}^{2} + {y}^{2} + {2x} + {2y} = 0$的圆心$(-1, -1)$到弦的距离$d = \dfrac{\sqrt{2}}{4}$，半径为$\sqrt{2}$，根据勾股定理有$\dfrac{l}{2} = \sqrt{r^2 - d^2} = \dfrac{\sqrt{30}}{4}$，因此弦长$l = \dfrac{\sqrt{30}}{2}$
    \end{jieda}

    \item 已知圆 ${C}_{1} : {\left( x + 1\right) }^{2} + {y}^{2} = {r}^{2}$ 过圆 ${C}_{2} : {\left( x - 4\right) }^{2} + {\left( y - 1\right) }^{2} = 4$ 的圆心,则两圆相交弦的方程为\blkda[6]{$5x+y-19=0,$$x\in[\dfrac{42}{13},4]$}.
    \begin{jieda}
        圆 $C_2:(x-4)^2+(y-1)^2=4$ 的圆心坐标为 $(4,1)$ ，因为圆 $C_1$ 过圆 $C_2$ 的圆心，所以 $(4+1)^2+1^2=r^2$ ，所以 $r^2=26$ ，所以 $C_1:(x+1)^2+y^2=26$ ，两圆的方程相减可得相交弦方程为 $5 x+y-19=0$ 
        因为是弦的方程，因此要有定义域，联立方程可得$x\in[\dfrac{42}{13},4]$．
    \end{jieda}

    \item 已知 $O$ 为坐标原点, $A\left( {3,0}\right) ,B$ 点在第一象限, $\bigtriangleup {OAB}$ 的内切圆 $Q$ 的方程为 ${\left( x - 2\right) }^{2} + {\left( y - 1\right) }^{2} = 1$ ,分别以 $O,A,B$ 为圆心作圆,且 $\odot  O, \odot  A, \odot  B$ 两两相外切,则 $\odot  B$ 的标准方程为\blkda[4]{$(x-3)^2+(y-4)^2=9$}.
    \begin{jieda}
        由已知，过点 A 与圆 $Q$ 相切的直线有两条，一条为直线 $O A$ ，另一条为直线 $A B$ ，如图，即 $x$ 轴与直线 $x=3$ ．
        \begin{center}
            \includegraphics[width=0.25\linewidth]{xx01ch02tupian2.png}
        \end{center}

        设 $\odot Q$ 分别切 $O A, A B$ 于 $M, N$ 两点，四边形 $Q M A N$ 为正方形．
            \end{jieda}
        则 $|Q N|=|Q M|=|M A|=1,|O M|=2,|Q A|=\sqrt{2},|A N|=1$ ，故 $\tan \angle Q O M=\dfrac{1}{2}$ ，
        设 $B(3, y)$ ，则 $\dfrac{y}{3}=\tan \angle B O A=\tan 2 \angle Q O M=\dfrac{2 \times \dfrac{1}{2}}{1-\left(\dfrac{1}{2}\right)^2}=\dfrac{4}{3}$ ，
        解得 $y=4$ ，故 $B(3,4)$ ，则 $|O B|=5,|A B|=4,|O A|=3$ ，
        
        设 $\odot O, \odot A, \odot B$ 的半径分别为 $r_O, r_A, r_B$ ，
        
        则 $\left\{\begin{array}{l}r_O+r_B=5 \\ r_O+r_A=3, \text { 解得 } r_B=3, \\ r_A+r_B=4\end{array}\right.$
        
        故所求 $\odot B$ 的标准方程为 $(x-3)^2+(y-4)^2=9$ ．
    

    \item 已知圆 $M : {x}^{2} + {y}^{2} - {2ay} + {a}^{2} - 4 = 0$ 与圆 $N : {x}^{2} + {y}^{2} - {4x} + 3 = 0$ 有 4 条公切线, 则 $a$ 的取值范围为\blkda[5]{$(-\infty,-\sqrt{5})\cup（\sqrt{5},+\infty)$}.
    \begin{jieda}
        由题可得，圆 $M: x^2+(y-a)^2=4$ ，圆心为 $(0, a)$ ，半径为 2 ，圆 $N:(x-2)^2+y^2=1$ ，圆心为 $(2,0)$ ，半径为 1 ，因为两圆有 4 条公切线，所以两圆外离，故圆心距 $|M N|=\sqrt{2^2+a^2}>3$ ，解得 $a \in(-\infty,-\sqrt{5}) \cup(\sqrt{5},+\infty)$.
    \end{jieda}

    
    \item 若直线 $y = {kx} - 1$ 与曲线 $y =  - \sqrt{1 - {\left( x - 2\right) }^{2}}$ 有公共点,则 $k$ 的取值范围是\blkda[4]{$[0,1]$}.

    
    \item 函数 $y = \dfrac{\sqrt{1 - {x}^{2}}}{2 + x}$ 的值域是\blkda[4]{$\left[0,\dfrac{\sqrt{3}}{3}\right]$}.

    \item 圆系: ${C}_{k} : {x}^{2} + {y}^{2} + {2kx} + \left( {{4k} + {10}}\right) y + {10k} + {20} = 0\left( {k \neq   - 1}\right)$ 中任两圆之间的关系是
    \dotfill{(\kh{D})}
    \xx{相交}{相离}{外切}{内切或外切}


    \item 从原点向圆 $C : {x}^{2} + {y}^{2} - {12y} + {27} = 0$ 作两条切线,则该圆夹在两条切线间的劣弧长是
    \dotfill{(\kh{B})}
    \xx{$\pi$}{${2\pi }$}{${4\pi }$}{${6\pi }$}


    \item 已知点 $P\left( {x,y}\right)$ 是直线 ${kx} + y + 4 = 0\left( {k > 0}\right)$ 上一动点,PA、PB 是圆 ${C} : {x}^{2} + {y}^{2} - {2y} = 0$ 的两条切线, ${A}\text{、}{B}$ 是切点,若四边形 ${PACB}$ 的最小面积是 2,则 ${k}$ 的值为
    \dotfill{(\kh{D})}
    \xx{ $\sqrt{2}$}{$\dfrac{\sqrt{21}}{2}$ }{$2\sqrt{2}$}{2}    

    
    \item 画出方程 $\left| x\right|  - 1 = \sqrt{1 - {\left( y - 1\right) }^{2}}$ 所表示的曲线.
    \begin{figure}
        \centering
        \includegraphics[width=0.5\linewidth]{xx01ch2tupian.png}
        \caption{12题}
        \label{fig:placeholder}
    \end{figure}
    
    \item 已知直线 $x + {2y} - 3 = 0$ 和圆 ${x}^{2} + {y}^{2} + x - {2ay} + a = 0$ 相交于 $A$ 、 $B$ 两点,且以 ${AB}$ 为直径的圆经过坐标原点,求实数 $a$ 的值.
    \begin{jieda}
        过$A,B$两点的圆系方程为$x^2+y^2+(1-\lambda)x-2(a+\lambda)y+(a+3\lambda) = 0$
        \begin{enumerate}
            \item 方法1
            
                半径$r = \dfrac{1}{2}\sqrt{(1-\lambda)^2+4(a+\lambda)^2-4(a+3\lambda)} = \dfrac{1}{2}\sqrt{5\lambda^2+(8a-14)\lambda + (4a^2-4a+1)}$，当$\lambda = \dfrac{7-4a}{5}$时，半径取得最小值，即以$A, B$为直径的圆，又因为该圆过原点，所以有$a + 3\lambda = 0$，两式联立$\left\{\begin{array}{l}
                 \lambda = \dfrac{7-4a}{5}  \\
                 a + 3\lambda = 0
            \end{array}\right.$，解得$\left\{\begin{array}{l}
                 \lambda = -1  \\
                 a = 3
            \end{array}\right.$.
            \item 方法2
            
            $x^2+y^2+(1-\lambda)x-2(a+\lambda)y+(a+3\lambda) = 0$应满足两个约束：过原点，圆心在$x + {2y} - 3 = 0$上

            满足过原点时方程变为$x^2+y^2+(1-\lambda)x+4\lambda y = 0$，圆心为$\left(\dfrac{\lambda-1}{2}, -2\lambda\right)$，代入直线方程可得$\lambda = -1$，故$a=3$.
        \end{enumerate}
        
    \end{jieda}

    \item 已知点 $P\left( {5,4}\right)$ ,圆 $C$ 的方程是 ${x}^{2} + {y}^{2} - {6x} - {8y} - {11} = 0$ . 判断是否存在满足下列三个条件的圆 $D$ : 
    \begin{enumerate*}
        \item 圆 $D$ 过点 $P$
        \item 圆 $D$ 与圆 $C$ 相切
        \item 圆 $D$ 的圆心的横、纵坐标均为正整数
    \end{enumerate*}
    若存在,求出圆 $D$ 的方程；若不存在,说明理由.
    \begin{jieda}
        存在满足条件的圆 $D$ ．圆 $C$ 即 $(x-3)^2+(y-4)^2=36$ ，易知点 $P(5,4)$ 在圆 $C$ 内部，圆 $D$ 在圆 $C$ 内部，所以两圆只能内切，设圆 $D$ 方程为 $(x-a)^2+(y-b)^2=r^2$ ．因为 $P$ 在圆 $D$ 上，所以 $(a-5)^2+(b-4)^2= r^2 $
        
        又因为两圆内切，所以 $\sqrt{(a-3)^2+(b-4)^2}=6-r$，故可得 $a=13-3 r$ ，易得 $1 \leqslant r \leqslant 4$
        
        \begin{dingautolist}{172}
            \item 当 $r=1$ 时，$a=10$ ，代入$(a-5)^2+(b-4)^2= r^2 $，得 $25+(b-4)^2=1$ 无解
            \item 当 $r=2$ 时，$a=7$ ，代入$(a-5)^2+(b-4)^2= r^2 $，得 $4+(b-4)^2=4$ ，所以 $b=4$
            \item 当 $r=3$ 时，$a=4$ ，代入$(a-5)^2+(b-4)^2= r^2 $，得 $1+(b-4)^2=9$ ，无整数解
            \item 当 $r=4$ 时，$a=1$ ，代入$(a-5)^2+(b-4)^2= r^2 $，得 $16+(b-4)^2=16$ ，所以 $b=4$
        \end{dingautolist}
        
        于是 $\left\{\begin{array}{l}a=7, \\ b=4, \\ r=2,\end{array}\right.$ 或 $\left\{\begin{array}{l}a=1, \\ b=4, \\ r=4,\end{array}\right.$ 因此，在在满足条件的圆 $D$, 它的方程是 $(x-7)^2+(y-4)^2=4$ 或 $(x-1)^2+(y-4)^2=16$ ．
    \end{jieda}
    
\end{enumerate}

\section*{椭圆的标准方程 }
\subsection*{椭圆的标准方程 (1)}

A 组:
\begin{enumerate}

    \item 集合 $P = \left\{  {M\left| {M{F}_{1}}\right|  + \left| {M{F}_{2}}\right|  = {2a}}\right\}  ,\left| {{F}_{1}{F}_{2}}\right|  = {2c}$ ,其中 $a > 0,c > 0$ ,且 $a\text{、}c$ 为常数：
    \begin{enumerate}
        \item 若\blkda[3]{$a>c$},则集合 $P$ 表示椭圆；
        \item 若\blkda[3]{$a=c$},则集合 $P$ 表示线段；
        \item 若\blkda[3]{$a<c$},则集合 $P$ 为空集.
    \end{enumerate}

    \item 椭圆 $\dfrac{{x}^{2}}{49} + \dfrac{{y}^{2}}{100} = 1$ 上一点到两个焦点的距离之和为\blkda[2]{20}.

    \item 过 $\left( {0,2}\right) .\left( {\dfrac{1}{2},\sqrt{3}}\right)$ 两点的椭圆的标准方程为\blkda[4]{$\dfrac{y^2}{4}+x^2=1$}.


    \item “ $1 < m < 7$ ” 是 “方程 $\dfrac{{x}^{2}}{7 - m} + \dfrac{{y}^{2}}{m - 1} = 1$ 表示的曲线为椭圆” 的\blkda[3]{必要非充分}条件.

    \item 已知两点 $A\left( {0, - 3}\right)$ 、 $B\left( {0,3}\right)$ ,若动点 $P$ 满足 $\left| {PA}\right|  + \left| {PB}\right|  = 6$ ,则点 $P$ 的轨迹为
    \dotfill{(\kh{D})}
    \xx{焦点在 $x$ 轴上的椭圆}{焦点在 $y$ 轴上的椭圆}{$x$ 轴上的线段 }{$y$ 轴上的线段}

    \item “平面内存在两个定点, 满足一动点到两定点的距离之和为一定值” 是 “这个动点的轨迹是椭圆”的\blkda[3]{必要非充分}条件

    \item 写出适合下列条件的椭圆的标准方程:
        \begin{enumerate}
            \item $a = 5,c = 2$ ,焦点在 $y$ 轴上； 
            \item 焦点是 $\left( {-3,0}\right) ,\left( {3,0}\right)$ ,且离心率为 $\dfrac{3}{5}$ ；
            \item 焦点在 $x$ 轴上,且经过点(2,0)和点(0,1)；
            \item 与椭圆 $4{x}^{2} + 9{y}^{2} = {36}$ 有相同焦点,且过点(3, - 2).
        \end{enumerate}
    \begin{jieda}
        \begin{enumerate}
            \item {$\dfrac{y^2}{25}+\dfrac{x^2}{21}=1$}
            \item {$\dfrac{x^2}{25}+\dfrac{y^2}{16}=1$}
            \item {$\dfrac{x^2}{4}+y^2=1$}
            \item {$\dfrac{x^2}{15}+\dfrac{y^2}{10}=1$}
        \end{enumerate}
    \end{jieda}
\end{enumerate}


B 组：
\begin{enumerate}
    \item 已知椭圆 $\dfrac{{x}^{2}}{16} + \dfrac{{y}^{2}}{4} = 1$ 上一点 $P$ 到其一个焦点的距离为 3,则点 $P$ 到另一个焦点的距离为\blkda[2]{5}.

    \item 若椭圆 $\dfrac{{x}^{2}}{{m}^{2}} + \dfrac{{y}^{2}}{4} = 1$ 经过点 $M\left( {-2,\sqrt{3}}\right)$ ,则这个椭圆的焦距长等于\blkda[2]{$4\sqrt{3}$}, 离心率等于\blkda[2]{$\dfrac{\sqrt{3}}{2}$}.

    \item 若椭圆 $\dfrac{{x}^{2}}{9} + \dfrac{{y}^{2}}{4} = 1$ 上点 $P$ 到右焦点的距离为 4 ,则点 $P$ 的横坐标为\blkda[3]{$-\dfrac{3\sqrt{5}}{5}$}.

    \item 已知 $F$ 为椭圆 $C : \dfrac{{x}^{2}}{4} + \dfrac{{y}^{2}}{3} = 1$ 的右焦点, $P$ 为 $C$ 上的一点,若 $\left| {PF}\right|  = 1$ ,则点 $P$ 的坐标为\blkda[3]{$(2,0)$}.

    \item 若方程 $\dfrac{{x}^{2}}{{a}^{2} - 4} + \dfrac{{y}^{2}}{3a} = 1$ 表示焦点在 $y$ 轴上的椭圆,则 $a$ 的取值范围是\blkda[3]{$a\in(2,4)$}.

    \item 已知椭圆的两个焦点为 ${F}_{1}\left( {-4,0}\right) ,{F}_{2}\left( {4,0}\right) ,{AB}$ 是椭圆过焦点 ${F}_{1}$ 的弦, 且 $\bigtriangleup  {AB}{F}_{2}$ 的周长为 20,那么该椭圆的方程为\blkda[4]{$\dfrac{x^2}{25}+\dfrac{y^2}{9}=1$}.

    \item 已知 ${F}_{1}\left( {-1,0}\right)$ 、 ${F}_{2}\left( {1,0}\right)$ 两点. 若动点 $P$ 满足 $\left| {P{F}_{1}}\right|  + \left| {P{F}_{2}}\right|  = 2\left| {{F}_{1}{F}_{2}}\right|$ ,则点 $P$ 的轨迹方程为\blkda[4]{$\dfrac{x^2}{4}+\dfrac{y^2}{3}=1$}.

    \item 若 $\bigtriangleup {ABC}$ 底边的两个端点分别是 ${B}\left( {-6,0}\right) ,{C}\left( {6,0}\right)$ ,周长为 32 ,则顶点 $A$ 的轨迹方程为\blkda[5]{$\dfrac{x^2}{100}+\dfrac{y^2}{64}=1(y\neq 0)$}.

    \item “ $m > 2$ ” 是 “方程 $\dfrac{{x}^{2}}{m - 2} + \dfrac{{y}^{2}}{5 - m} = 1$ 表示的曲线是椭圆” 的\blkda[2]{必要非充分}条件


    \item 在 $\bigtriangleup  {ABC}$ 中,已知 $A\left( {-1,0}\right) \text{、}C\left( {1,0}\right)$ ,且三边长满足 $\left| {BC}\right|  > \left| {AC}\right|  > \left| {AB}\right|$ . 若 $\sin A + \sin C = 2\sin B$ ,则动点 $B$ 的轨迹方程是\blkda[4]{$\dfrac{x^2}{4}+\dfrac{y^2}{3}=1(x<0)$}.


    \item 已知 ${F}_{1}\left( {-4,0}\right) \text{、}{F}_{2}\left( {4,0}\right)$ 两定点,若动点 $M$ 满足 $\left| {M{F}_{1}}\right|  + \left| {M{F}_{2}}\right|  = {2a}\left( {a > 0}\right)$ , 求动点 $M$ 的轨迹.
    \begin{jieda}
         当 $a > 4$ 时,动点M  的轨迹是以 $ {F}_{1}\text{、}{F}_{2} $ 为焦点的椭圆 $ \dfrac{{x}^{2}}{{a}^{2}} + \dfrac{{y}^{2}}{{a}^{2} - {16}} = 1 $ ; 当 $ a = 4 $ 时,它的轨迹是一条
线段 $ y = 0\left( {-4 \leq  x \leq  4}\right) $ ; 当 $ a < 4 $ 时,轨迹不存在.
    \end{jieda}

    \item 已知点 $A\left( {\sqrt{3},0}\right) ,Q$ 是圆 ${x}^{2} + {y}^{2} = 4$ 上的动点,线段 ${AQ}$ 的垂直平分线交半径 ${OQ}$ 于点 $P$ . 当点 $Q$ 在圆周上运动时,求点 $P$ 的轨迹方程.
    \begin{jieda}
        因为 $\left| {PA}\right|  = \left| {PQ}\right|$ ,所以 $\left| {PO}\right|  + \left| {PA}\right|  = \left| {PO}\right|  + \left| {PQ}\right|  = 2 > \left| {OA}\right|$ .
由椭圆定义,知点 $P$ 的轨迹是以 $O\left( {0,0}\right) \text{、}A\left( {\sqrt{3},0}\right)$ 为焦点,以 2 为长轴长的椭圆,其中
心为 $\left( {\dfrac{\sqrt{3}}{2},0}\right) ,a = 1,c = \dfrac{\sqrt{3}}{2},{b}^{2} = \dfrac{1}{4}$ . 因此,所求轨迹方程为 ${\left( x - \dfrac{\sqrt{3}}{2}\right) }^{2} + 4{y}^{2} = 1$ .
    \end{jieda}
    
    \item 已知圆 ${x}^{2} + {y}^{2} - {6x} - {55} = 0$ ,动圆 $M$ 经过定点 $A\left( {-3,0}\right)$ ,且与已知圆内切,求圆心 $M$ 的轨迹方程.
    \begin{jieda}
        已知圆 ${\left( x - 3\right) }^{2} + {y}^{2} = {64}$ 的圆心为 $C\left( {3,0}\right)$ ,半径 $r = 8$ . 
        
        设圆 $M$与圆 $C$ 内切于 $B$ ,则 $C$ 、 $M$ 、 $B$ 在一直线上,即 $\left| {MB}\right|  + \left| {MC}\right|  = \left| {CB}\right|  = 8$ . 又因为$\left| {MB}\right|  = \left| {MA}\right|$ ,所以 $\left| {MA}\right|  + \left| {MC}\right|  = 8$ . 因为 $8 > 6 = \left| {AC}\right|$ ,所以由椭圆定
        义,知点 $M$ 的轨迹是以 $A\left( {-3,0}\right) \text{、}C\left( {3,0}\right)$ 为焦点,长轴长为 8 的椭圆,且 $a = 4$ ,$c = 3,{b}^{2} = {a}^{2} - {c}^{2} = 7$ . 故点 $M$ 的轨迹方程是 $\dfrac{{x}^{2}}{16} + \dfrac{{y}^{2}}{7} = 1$ .
    \end{jieda}
\end{enumerate}


C 组:
\begin{enumerate}

    \item 已知点 $M\left( {0,2}\right) ,N\left( {-3,6}\right)$ 到直线 $l$ 的距离分别为 1,4,则满足条件的直线 $l$ 的条数是 \blkda[2]{3}.
    \begin{jieda}
        易得 ${MN} = 5$ ,以点 $M$ 为圆心、半径为1的圆与以点$N$ 为圆心、半径为 4 的圆外切,故满足条件的直线 $l$有3条.
    \end{jieda}
    
    \item 在平面直角坐标系 ${xOy}$ 中,圆 $O : {x}^{2} + {y}^{2} = 1$ ,圆 ${O}_{1} : {\left( x - 3\right) }^{2} + {y}^{2} = 4$ . 过 $x$ 轴的左半轴上一点 $M$ 作圆 $O$ 的切线,与圆 $O$ 相切于点 $A$ ,与圆 ${O}_{1}$ 分别相交于点 $B,C$ ,若 ${AB} = {BC}$ ,则点 $M$ 的坐标为\blkda[3]{(-4,0)}.

    \item 满足 $\sqrt{1 - {x}^{2}} \geq  x$ 的实数 $x$ 的取值范围是\blkda[3]{$[-1,\dfrac{\sqrt{2}}{2}]$}
    
\end{enumerate}

\clearpage
\subsection*{椭圆的标准方程 (2)}

\begin{enumerate}


    \item 已知椭圆的两个焦点坐标分别是 $\left( {0,5}\right) ,\left( {0, - 5}\right)$ ,椭圆上一点 $P$ 到两个焦点的 距离之和为 26,则该椭圆方程为\blkda[3]{$\dfrac{y^2}{169}+\dfrac{x^2}{144}=1$}.

    \item 已知椭圆 $9{x}^{2} + {25}{y}^{2} - {225} = 0$ 的两焦点为 ${F}_{1}\text{、}{F}_{2}$ ,若 ${AB}$ 为过 ${F}_{1}$ 的弦, 则 $\bigtriangleup  {AB}{F}_{2}$ 的周长为\blkda[3]{20}.

    \item 在平面直角坐标系 ${xoy}$ 中,已知 $\bigtriangleup {ABC}$ 的顶点 $A\left( {-1,0}\right)$ 和 $C\left( {1,0}\right)$ ,且顶点 $B$ 在椭圆 $\dfrac{{x}^{2}}{4} + \dfrac{{y}^{2}}{3} = 1$ 上,则 $\dfrac{\sin A + \sin C}{\sin B}$ 的值等于\blkda[2]{2}. 
    
    \item 若方程 $\dfrac{{x}^{2}}{4 - k} + \dfrac{{y}^{2}}{6 + k} = 1$ 表示椭圆,则实数 $k$ 的范围是\blkda[3]{$(-6,-1)\cup (-1,4)$}. 
    
    \item 若方程 ${x}^{2}\cos \theta  + {y}^{2}\sin \theta  = 1\left( {\theta  \in  \left( {0,{2\pi }}\right) }\right)$ 表示椭圆,则 $\theta$ 的取值范围是\blkda[3]{$(0,\dfrac{\pi}{4})\cup (\dfrac{\pi}{4},\dfrac{\pi}{2})$}. 
    
    \item 如果方程 ${x}^{2} + k{y}^{2} = 2$ 表示焦点在 $y$ 轴的椭圆,那么实数 $k$ 的取值范围是\blkda[3]{$k\in (0,1)$}. 
    
    \item 已知椭圆 $\dfrac{{x}^{2}}{8} + \dfrac{{y}^{2}}{{m}^{2}} = 1$ 的焦点在 ${x}$ 轴上,则其焦距等于\blkda[3]{$2\sqrt{8-m^2}$}. 
    
    
    \item 椭圆 $5{x}^{2} - k{y}^{2} = 5$ 的一个焦点是(0,2),那么 $k =$ \blkda[2]{-1}. 
    
    \item 若方程 $\left( {m - 2}\right) {x}^{2} + \left( {4 - m}\right) {y}^{2} = \left( {m - 2}\right) \left( {4 - m}\right)$ 表示焦点在 $y$ 轴上的椭圆,则实数 $m$ 的取值范围是\blkda[2]{(3,4)}.

    \item 椭圆 $a{x}^{2} + b{y}^{2} + {ab} = 0\left( {a < b < 0}\right)$ 的焦点坐标为\blkda[5]{$(0,-\sqrt{b-a})$和$(0,\sqrt{b-a})$}.

    \item 已知椭圆 $C : \dfrac{{x}^{2}}{36} + \dfrac{{y}^{2}}{9} = 1$ 的左右焦点分别为 ${F}_{1},{F}_{2}$ ,点 $M$ 是椭圆 $C$ 上一点, 点 $N$ 是线段 $M{F}_{1}$ 的中点, $O$ 为坐标原点,若 $\left| {MN}\right|  = 4$ ,则 $\left| {ON}\right|  =$ \blkda[2]{2}.

    \item 设点 $A\left( {-2,\sqrt{3}}\right) \text{、}B\left( {2,0}\right)$ ,点 $M$ 在椭圆 $\dfrac{{x}^{2}}{16} + \dfrac{{y}^{2}}{12} = 1$ 上运动,当 $\left| {MA}\right|  + \left| {MB}\right|$ 最大时,点 $M$ 的坐标为\blkda[3]{$(-2, -3)$}.
    \begin{jieda}
        $|MA|+|MB| = |MF_1|+|MB|+|MA|-|MF_1| = 8+|MA|-|MF_1| \leq 8 + |AF_1| = 8+\sqrt{3}$

        当且仅当$M$坐标为$(-2,-3)$，$A,F_1,M$三点共线时，取得最大值.
    \end{jieda}

    \item 设集合 $A = \{ 1,2,3,4,5\} ,a,b \in  A$ ,则方程 $\dfrac{{x}^{2}}{a} + \dfrac{{y}^{2}}{b} = 1$ 表示焦点位于 $y$ 轴上的椭圆有\blkda[2]{10}个.


    \item 已知边长为 4 的菱形 ${ABCD}$ 的两条对角线在坐标轴上,且顶角 $\angle {DAB} = {60}^{ \circ  }$ . 以菱形 ${ABCD}$ 的不相邻的两个顶点 $A\text{、}C$ 为焦点,且过菱形的另两个顶点的椭圆方程是\blkda[8]{$\dfrac{x^2}{16}+\dfrac{y^2}{4}=1$或$\dfrac{y^2}{16}+\dfrac{x^2}{4}=1$}.

    \item 已知椭圆 $C : \dfrac{{x}^{2}}{4} + {y}^{2} = 1$ 的左、右焦点分别为 ${F}_{1},{F}_{2}$ ,点 $P$ 为椭圆上一点,且点 $P$ 位于第一象限, $P{F}_{1} \bot  P{F}_{2}$ ,则 $\left| {P{F}_{1}}\right|  =$ \blkda[3]{$2+\sqrt{2}$}, $\left| {P{F}_{2}}\right|  =$ \blkda[3]{$2-\sqrt{2}$}.

    \item 已知点 ${M}\left( {x,y}\right)$ 与定点 ${F}\left( {c,0}\right)$ 的距离和它到定直线 $l : x = \dfrac{{a}^{2}}{c}$ 的比是离心率 $\dfrac{c}{a}\left( {a > c > 0}\right)$ ,求点 ${M}$ 的轨迹方程.
    \begin{jieda}
        $\dfrac{x^2}{a^2}+\dfrac{y^2}{a^2-c^2}=1$
    \end{jieda}
    

    \item 求与两圆 ${\left( x + 2\right) }^{2} + {y}^{2} = 1,{x}^{2} - {4x} + {y}^{2} - {45} = 0$ 都相切的动圆圆心的轨迹方程.
    \begin{jieda}
        \begin{dingautolist}{172}
            \item 与圆1外切，与圆2内切，此时轨迹方程为$\dfrac{{x}^{2}}{16} + \dfrac{{y}^{2}}{12} = 1$
            \item 与两圆均内切，此时轨迹方程为$\dfrac{{x}^{2}}{9} + \dfrac{{y}^{2}}{5} = 1$
        \end{dingautolist}
    \end{jieda}

    \item 已知椭圆的中心在坐标原点 ${O}$ ,焦点在 $x$ 轴上,直线 $y = x + 1$ 与椭圆相交于点 $P$ 和点 $Q$ ,且 ${OP} \bot  {OQ},\left| {PQ}\right|  = \dfrac{\sqrt{10}}{2}$ ,求椭圆方程.

    \begin{jieda}
        设椭圆的标准方程是$\dfrac{x^2}{a^2}+\dfrac{y^2}{b^2} = 1(a > b > 0)$，联立直线和椭圆的方程，得到方程组$\left\{\begin{array}{l}
             y = x + 1 \\
             \dfrac{x^2}{a^2}+\dfrac{y^2}{b^2} = 1
        \end{array}\right.$，消去$y$，得到$b^2x^2+a^2(x+1)^2 = a^2b^2$，整理得$(a^2+b^2)x^2+2a^2x+a^2(1-b^2) = 0$，设$P(x_1, y_2), Q(x_2, y_2)$，根据韦达定理有$\left\{\begin{array}{l}
             x_1+x_2 = -\dfrac{2a^2}{a^2+b^2} \\
             x_1x_2 = \dfrac{a^2(1-b^2)}{a^2+b^2}
        \end{array}\right.$\\
        根据${OP} \bot  {OQ}$可知$x_1x_2+y_1y_2 = 0$，即$2x_1x_2+(x_1+x_2)+1 = 0$，根据$\left| {PQ}\right|  = \dfrac{\sqrt{10}}{2}$可知$|x_1-x_2| = \dfrac{\sqrt{5}}{2}$，即$(x_1+x_2)^2-4x_1x_2 = \dfrac{5}{4}$，联立两式可得$\left\{\begin{array}{l}
             x_1+x_2 = -\dfrac{3}{2} \\
             x_1x_2 = \dfrac{1}{4}
        \end{array}\right.$，因此有$\left\{\begin{array}{l}
             -\dfrac{3}{2} = -\dfrac{2a^2}{a^2+b^2} \\
             \dfrac{1}{4} = \dfrac{a^2(1-b^2)}{a^2+b^2}
        \end{array}\right.$，解得$a^2 = 2, b^2 = \dfrac{2}{3}$，因此$\dfrac{x^2}{2}+\dfrac{3y^2}{2} = 1$
    \end{jieda}
    
    \item 以椭圆 $\dfrac{{x}^{2}}{{a}^{2}} + {y}^{2} = 1\left( {a > 1}\right)$ 上的一个端点 $B\left( {0,1}\right)$ 为直角顶点,作椭圆内接等腰直角三角形 ${ABC}$ ,这样的三角形存在吗? 若存在,最多能作几个?
    \begin{jieda}
        由题设,知 ${AB}\text{、}{BC}$ 不可能垂直或平行于 $x$ 轴,故不妨设 ${BC}$ 所在直线方程为 $y = {kx} + 1\left( {k < 0}\right)$ ,代入椭圆方程,整理得 $\left( {1 + {a}^{2}{k}^{2}}\right) {x}^{2} + {2k}{a}^{2}x = 0$ ,解得${x}_{1} = 0,{x}_{2} =  - \dfrac{2{a}^{2}k}{1 + {a}^{2}{k}^{2}}$ . 因此 $\left| {BC}\right|  = \sqrt{1 + {k}^{2}}\left| {{x}_{1} - {x}_{2}}\right|  = \dfrac{2{a}^{2}\left| k\right| \sqrt{1 + {k}^{2}}}{1 + {a}^{2}{k}^{2}}$ . 
        
        同理,用 $- \dfrac{1}{k}$ 代上式中的 $k$ ,得 $\left| {AB}\right|  = \dfrac{2{a}^{2}\sqrt{1 + {k}^{2}}}{{a}^{2} + {k}^{2}}$ .
        
        由 $\left| {BC}\right|  = \left| {AB}\right|$ ,得 $\left| k\right| \left( {{a}^{2} + {k}^{2}}\right)  = 1 + {a}^{2}{k}^{2}$ . 又 $k < 0$ ,得 $\left( {k + 1}\right) \left\lbrack  {{k}^{2} + \left( {{a}^{2} - 1}\right) k}+1\right\rbrack  = 0$ . 
        
        当 ${\left( {a}^{2} - 1\right) }^{2} - 4 < 0$ ,即 $1 < a < \sqrt{3}$ 时,方程$\left( {k + 1}\right) \left\lbrack  {{k}^{2} + \left( {{a}^{2} - 1}\right) k}+1\right\rbrack  = 0$有唯一解 $k =  - 1$ ; 
        
        当 $a = \sqrt{3}$ 时,方程 ${k}^{2} + \left( {{a}^{2} - 1}\right) k + 1 = 0$ 的解为 $k =  - 1$ ,此时方程$\left( {k + 1}\right) \left\lbrack  {{k}^{2} + \left( {{a}^{2} - 1}\right) k}+1\right\rbrack  = 0$有唯一解 $k =  - 1$ ; 
        
        当 $a > \sqrt{3}$ 时,方程$\left( {k + 1}\right) \left\lbrack  {{k}^{2} + \left( {{a}^{2} - 1}\right) k}+1\right\rbrack  = 0$有三个不同解. 
        
        综上所述,当 $1 < a \leq  \sqrt{3}$ 时,最多作出一个满足条件的三角形; 当 $a > \sqrt{3}$ 时,最多可作出三个满足条件的三角形.
    \end{jieda}

\end{enumerate}


\section*{椭圆的性质}
\subsection*{椭圆的性质(1)}

A 组：
\begin{enumerate}
    
    \item 椭圆 $4{x}^{2} + 2{y}^{2} = {16}$ 的长轴长是\blkda[2]{$4\sqrt{2}$},短轴长是\blkda[2]{4},离心率是\blkda[2]{$\dfrac{\sqrt{2}}{2}$}, 焦点坐标是\blkda[4]{$(0,2) (0,-2)$},上顶点坐标是\blkda[3]{$(0,2\sqrt{2})$}.

    \item 椭圆的一个焦点到长轴两端点的距离之比为 $1 : 4$ ,短轴长为 8 ,则椭圆的标准方程为\blkda[8]{$\dfrac{{x}^{2}}{25} + \dfrac{{y}^{2}}{16} = 1$ 或 $\dfrac{{y}^{2}}{25} + \dfrac{{x}^{2}}{16} = 1$}.

    \item 已知椭圆 $\dfrac{{x}^{2}}{{a}^{2}} + \dfrac{{y}^{2}}{{b}^{2}} = 1\left( {a > b > 0}\right)$ 的离心率为 $\dfrac{\sqrt{2}}{2}$ ,则长轴与短轴的比值为\blkda[3]{$\sqrt{2}$}.

    \item 若椭圆 $C : \dfrac{{x}^{2}}{5} + \dfrac{{y}^{2}}{m} = 1$ 的一个焦点坐标为(-1,0),则实数 $m$ 的值为\blkda[2]{4}.

    \item 若椭圆 $C : \dfrac{{x}^{2}}{m} + \dfrac{{y}^{2}}{9} = 1\left( {m > 9}\right)$ 比椭圆 $D : \dfrac{{x}^{2}}{6} + \dfrac{{y}^{2}}{3} = 1$ 更扁,则椭圆 $C$ 的长轴长的取值范围是\blkda[4]{$(6\sqrt{2},+\infty)$}.

    \item 已知椭圆 $C : \dfrac{{y}^{2}}{5} + {x}^{2} = 1$ 的右顶点为 ${A},P$ 为 $C$ 上一点,则 ${PA}$ 的最大值为\blkda[3]{$\dfrac{5}{2}$}.
\end{enumerate}

B 组:
\begin{enumerate}


    \item 已知椭圆 $\dfrac{{x}^{2}}{16} + \dfrac{{y}^{2}}{7} = 1$ 均两个焦点为 ${F}_{1}\text{、}{F}_{2}$ . 若点 $P$ 是此椭圆上的点,且 $\vv{P{F}_{1}}\bot \vv{P{F}_{2}}$ , 则 $\left| \vv{OP}\right|  =$ \blkda[2]{3}.

    \item 已知椭圆 $\dfrac{{x}^{2}}{9} + \dfrac{{y}^{2}}{4} = 1$ 的焦点为 ${F}_{1}$ 、 ${F}_{2}$ ,点 $P$ 为椭圆上的动点,当 $\angle {F}_{1}P{F}_{2}$ 是钝角时,则 $P$ 的横坐标的范围是\blkda[4]{$\left(-\dfrac{3\sqrt{5}}{5},\dfrac{3\sqrt{5}}{5}\right)$}.

    \item 若椭圆 ${x}^{2} + m{y}^{2} = 1$ 的焦点在 $y$ 轴上,长轴长是短轴长的两倍,则 $m$ 的值是\blkda[2]{$\dfrac{1}{4}$}.

    \item 已知 $\bigtriangleup  {ABC}$ 以椭圆上一点和椭圆两焦点为顶点,若三角形 ${ABC}$ 面积的最大值为1 , 则椭圆长轴长的最小值为\blkda[2]{$2\sqrt{2}$}.

    \item 若椭圆的长轴长、焦距、短轴长依次成等差数列,则满足条件的椭圆的标准方程可以为\blkda[4]{$\dfrac{{x}^{2}}{25} + \dfrac{{y}^{2}}{9} = 1$}.


    \item 已知椭圆 $\dfrac{{x}^{2}}{98} + \dfrac{{y}^{2}}{49} = 1$ 内有一点 $A\left( {0,5}\right)$ . 若在此椭圆上求一点 $B$ ,使 $\left| {AB}\right|$ 最大,则点 $B$ 的坐标是\blkda[4]{$(4\sqrt{3},-5)$或($4\sqrt{3},-5)$}.

    \item 已知方程 $y = {ax} + b$ 和 ${a}^{2}{x}^{2} + {y}^{2} = {b}^{2}\left( {a > b > 1}\right)$ ,它们在同一坐标系中的图像可能是下图中的\blkda[2]{C}
    \begin{center}
    \includegraphics[width=0.3\textwidth]{tuoyuandebiaozhunfangchengtupian1.png}
\end{center}


    \item 方程 ${x}^{2}\sin {2\theta } + {y}^{2}\cos \theta  = 1\left( {\theta  \in  \lbrack 0,{2\pi }}\right) )$ 表示椭圆,则 $\theta$ 的取值范围是
    \dotfill{(\kh{C})}
    \xx{$\left( {0,\dfrac{\pi }{2}}\right)$ }
    {$\left\lbrack  {0,\dfrac{\pi }{2}}\right\rbrack$ }
    {$\left( {0,\dfrac{\pi }{6}}\right)  \cup  \left( {\dfrac{\pi }{6},\dfrac{\pi }{2}}\right)$ }
    {$\left\lbrack  {0,\dfrac{\pi }{6}}\right)  \cup  \left( {\dfrac{\pi }{6},\dfrac{\pi }{2}}\right\rbrack$}


    \item 已知 $P$ 是椭圆 $\dfrac{{x}^{2}}{4} + \dfrac{{y}^{2}}{3} = 1$ 上的点, ${F}_{1},{F}_{2}$ 是两个焦点,则 $\left| {P{F}_{1}}\right|  \cdot  \left| {P{F}_{2}}\right|$ 的最大值与最小值之差是 \blkda[2]{1}


    \item 已知 $P$ 为椭圆 $\dfrac{{x}^{2}}{25} + \dfrac{{y}^{2}}{9} = 1$ 上的点, ${F}_{1}\text{、}{F}_{2}$ 为椭圆的焦点. 若 $\angle {F}_{1}P{F}_{2} = \dfrac{\pi }{6}$ ,则 $\bigtriangleup  P{F}_{1}{F}_{2}$ 的面积等于\blkda[3]{$9\left( {2 - \sqrt{3}}\right)$}


    \item 已知椭圆 $\dfrac{{x}^{2}}{16} + \dfrac{{y}^{2}}{{b}^{2}} = 1\left( {b > 0}\right)$ 的两焦点 ${F}_{1}\text{、}{F}_{2}$ 在 $x$ 轴上, $P$ 为椭圆上任意一点,且 $\angle {F}_{1}P{F}_{2}$ 的最大值为 $\dfrac{\pi }{3}$ ,求 $b$ 的值.
    \begin{jieda}
        当$P$点在椭圆的上顶点时，角度最大，从而可得$b=2\sqrt{3}$
    \end{jieda}


    \item 已知椭圆 $C : \dfrac{{x}^{2}}{{m}^{2}} + {y}^{2} = 1$ (常数 $m > 1$ ),点 $P$ 是 $C$ 上的动点, $M$ 是右顶点,定点 $A$ 的坐标为(2,0). (1) 若 $M$ 与 $A$ 重合,求 $C$ 的焦点坐标；(2) 若 $m = 3$ ,求 $\left| {PA}\right|$ 的最大值与最小值.
    \begin{jieda}
        (1)根据题意,若 $ M $ 与 $ A $ 重合,即椭圆的右顶点
的坐标为(2,0)；
则 $ a = 2 $ ；椭圆的焦点在 $ x $ 轴上,则 $ c = \sqrt{3} $ ；
则椭圆焦点的坐标为 $ \left( {\sqrt{3},0}\right) $ , $ \left( {-\sqrt{3},0}\right) $ ；
(2)若 $ m = 3 $ ,则椭圆的方程为 $ \dfrac{{x}^{2}}{9} + {y}^{2} = 1 $ ,变
形可得 $ {y}^{2} = 1 - \dfrac{{x}^{2}}{9} $ ,
$ {\left| PA\right| }^{2} = {\left( x - 2\right) }^{2} + {y}^{2} = {x}^{2} - {4x} + 4 + {y}^{2} = $
$ \dfrac{8{x}^{2}}{9} - {4x} + 5 $
又由 $ - 3 \leq  x \leq  3 $ ,
根据二次函数的性质, 分析可得,
$ x =  - 3 $ 时, $ {\left| PA\right| }^{2} = \dfrac{8{x}^{2}}{9} - {4x} + 5 $ 取得最大
值,且最大值为 25；
$ x = \dfrac{9}{4} $ 时, $ {\left| PA\right| }^{2} = \dfrac{8{x}^{2}}{9} - {4x} + 5 $ 取得最小
值,且最小值为 $ \dfrac{1}{2} $ ；
则 $ \left| {PA}\right| $ 的最大值为5, $ \left| {PA}\right| $ 的最小值为 $ \dfrac{\sqrt{2}}{2} $ ；
    \end{jieda}


    \item 已知 ${F}_{1}\text{、}{F}_{2}$ 为椭圆 $\dfrac{{x}^{2}}{9} + \dfrac{{y}^{2}}{4} = 1$ 的两个焦点, $P$ 为椭圆上的一点, $P\text{、}{F}_{1}\text{、}{F}_{2}$ 是一个直角三角形的三个顶点,且 $\left| {P{F}_{1}}\right|  > \left| {P{F}_{2}}\right|$ ,求 $\dfrac{\left| P{F}_{1}\right| }{\left| P{F}_{2}\right| }$ .

    \begin{jieda}
        若 $ \angle P{F}_{2}{F}_{1} $ 为直角,则 $ {\left| P{F}_{1}\right| }^{2} = {\left| P{F}_{2}\right| }^{2} + {\left| {F}_{1}{F}_{2}\right| }^{2} $ ,即 $ {\left| P{F}_{1}\right| }^{2} $
$ = {\left( 6 - \left| P{F}_{1}\right| \right) }^{2} + {20} $ ,
得 $ \left| {P{F}_{1}}\right|  = \dfrac{14}{3},\left| {P{F}_{2}}\right|  = \dfrac{4}{3} $ ,故 $ \dfrac{\left| P{F}_{1}\right| }{\left| P{F}_{2}\right| } = \dfrac{7}{2} $ ;
若 $ \angle {F}_{1}P{F}_{2} $ 为直角,则 $ {\left| {F}_{1}{F}_{2}\right| }^{2} = {\left| P{F}_{1}\right| }^{2} + {\left| P{F}_{2}\right| }^{2} $ ,
即 $ {20} = {\left| P{F}_{1}\right| }^{2} + {\left( 6 - \left| P{F}_{1}\right| \right) }^{2} $ ,
得 $ \left| {P{F}_{1}}\right|  = 4,\left| {P{F}_{2}}\right|  = 2 $ ,故 $ \dfrac{\left| P{F}_{1}\right| }{\left| P{F}_{2}\right| } = 2 $ .
    \end{jieda}

    \item 已知椭圆 $C$ 的两个焦点分别为 ${F}_{1}\left( {-1,0}\right)$ 、 ${F}_{2}\left( {1,0}\right)$ ,短轴的两个端点分别为 ${B}_{1}\text{、}{B}_{2}$ . 
    \begin{enumerate}
        \item 若 $\bigtriangleup {F}_{1}{B}_{1}{B}_{2}$ 为等边三角形,求椭圆 $C$ 的标准方程;
        \item 若椭圆 $C$ 的短轴长为 2,过点 ${F}_{2}$ 的直线 $l$ 与椭圆 $C$ 相交于 $P$ 、 $Q$ 两点,且 $\vv{{F}_{1}P}\bot \vv{{F}_{1}Q}$ ,求直线 $l$ 的方程.
    \end{enumerate}

    \begin{jieda}
    \begin{enumerate}
        \item 由题意知, $ \bigtriangleup {F}_{1}{B}_{1}{B}_{2} $ 为等边三角形,
            所以 $ c = \sqrt{3}b $ ,又 $ c = 1 $ ,
            所以 $ b = \dfrac{\sqrt{3}}{3} $ ,
            又由 $ {a}^{2} = {b}^{2} + {c}^{2} $ ,可得 $ {a}^{2} = \dfrac{4}{3} $ ,
            故椭圆 $ C $ 的方程为 $ \dfrac{3{x}^{2}}{4} + 3{y}^{2} = 1 $ .
        \item 易知椭圆 $ C $ 的方程为 $ \dfrac{{x}^{2}}{2} + {y}^{2} = 1 $ ,
            当直线 $ l $ 的斜率不存在时,其方程为 $ x = 1 $ ,不符合题意;
            当直线 $ l $ 的斜率存在时,设直线 $ l $ 的方程为 $ y = k\left( {x - 1}\right) $ ,
            由 $ \left\{  \begin{array}{l} y = k\left( {x - 1}\right) , \\  \dfrac{{x}^{2}}{2} + {y}^{2} = 1, \end{array}\right. $
            得 $ \left( {2{k}^{2} + 1}\right) {x}^{2} - 4{k}^{2}x + 2\left( {{k}^{2} - 1}\right)  = 0 $ ,
            设 $ P\left( {{x}_{1},{y}_{1}}\right) ,Q\left( {{x}_{2},{y}_{2}}\right) $ ,
            则 $ {x}_{1} + {x}_{2} = \dfrac{4{k}^{2}}{2{k}^{2} + 1},{x}_{1}{x}_{2} = \dfrac{2\left( {{k}^{2} - 1}\right) }{2{k}^{2} + 1} $ ,
            $ \vv{{F}_{1}P} = \left( {{x}_{1} + 1,{y}_{1}}\right) , $
            $ \vv{{F}_{1}Q} = \left( {{x}_{2} + 1,{y}_{2}}\right) , $
            因为 $ \vv{{F}_{1}P} \bot  \vv{{F}_{1}Q} $ ,
            所以 $ \vv{{F}_{1}P} \cdot  \vv{{F}_{1}Q} = 0 $ ,
            即 $ \left( {{x}_{1} + 1}\right) \left( {{x}_{2} + 1}\right)  + {y}_{1}{y}_{2} = {x}_{1}{x}_{2} + \left( {{x}_{1} + {x}_{2}}\right)  + 1 + {k}^{2}\left( {{x}_{1} - }\right. $
            $ 1)\left( {{x}_{2} - 1}\right)  = \left( {{k}^{2} + 1}\right) {x}_{1}{x}_{2} - \left( {{k}^{2} - 1}\right) \left( {{x}_{1} + {x}_{2}}\right)  + {k}^{2} + 1 = $
            $ \dfrac{7{k}^{2} - 1}{2{k}^{2} + 1} = 0, $
            解得 $ {k}^{2} = \dfrac{1}{7} $ ,即 $ k =  \pm  \dfrac{\sqrt{7}}{7} $ ,
            故直线 $ l $ 的方程为 $ x + \sqrt{7}y - 1 = 0 $ 或 $ x - \sqrt{7}y - 1 = 0 $ .
    \end{enumerate}
    \end{jieda}
\end{enumerate}

\clearpage
\subsection*{椭圆的性质(2)}
\begin{enumerate}

 \item 已知椭圆 $4{x}^{2} + 9{y}^{2} = {144}$ 内一点 $P\left( {3,2}\right)$ ,过 $P$ 点的弦恰好以 $P$ 为中点,则这条弦 所在的直线方程是\blkda[3]{$2x+3y-12=0$}

 \item 已知线段 $\left| {AB}\right|  = 4,\left| {PA}\right|  + \left| {PB}\right|  = 6,C$ 是 ${AB}$ 的中点,当动点 $P$ 在同一平面内运动时, $\left| {PC}\right|$ 的最大值为\blkda[2]{3},最小值为\blkda[2]{$\sqrt{5}$}.

 \item 在椭圆 $\dfrac{{x}^{2}}{{a}^{2}} + \dfrac{{y}^{2}}{{b}^{2}} = 1\left( {a > b > 0}\right)$ 中,设左焦点、右顶点、短轴上方的顶点的坐标分别为 $F\left( {-c,0}\right)$ 、 $A\left( {a,0}\right)$ 、 $B\left( {0,b}\right)$ . 若 $a,b,c$ 成等比数列,则 $\angle {ABF}$ 的大小为\blkda[2]{$\dfrac{\pi}{2}$}

 \item 若 ${AB}$ 是椭圆 $\dfrac{{x}^{2}}{{a}^{2}} + \dfrac{{y}^{2}}{{b}^{2}} = 1\left( {a > b > 0}\right)$ 中不平行于对称轴的一条弦, $M$ 是 ${AB}$ 的中点, $O$ 是椭圆的中心,则 ${k}_{AB} \cdot  {k}_{OM} =$ \blkda[2]{$-\dfrac{b^2}{a^2}$}

 \item 椭圆 $\dfrac{{x}^{2}}{{a}^{2}} + \dfrac{{y}^{2}}{{b}^{2}} = 1\left( {a > b > 0}\right)$ 上对两焦点张角为 ${90}^{ \circ  }$ 的点有\blkda[12]{0个（当$\sqrt{2}a<b$）；2个（当$\sqrt{2}a=b$）；4个（当$\sqrt{2}a>b$）；}个. (填出所有可能情况及相应条件)

 \item 设点 ${F}_{1},{F}_{2}$ 分别为椭圆 $C : \dfrac{{x}^{2}}{2} + {y}^{2} = 1$ 的左、右焦点,则椭圆 $C$ 的离心率为\blkda[2]{$\dfrac{\sqrt{2}}{2}$}, 经过原点且斜率不为 0 的直线 $l$ 与椭圆 $C$ 交于 $P,Q$ 两点,当四边形 $P{F}_{1}Q{F}_{2}$ 的面积最大时, $\vv{P{F}_{1}} \cdot  \vv{P{F}_{2}} =$ \blkda[2]{0}.

 \item 若直线 $y = {kx} + 1$ 和椭圆 $\dfrac{{x}^{2}}{25} + \dfrac{{y}^{2}}{m} = 1$ 恒有公共点,则实数 $m$ 的取值范围为\blkda[4]{$[1,25)\cup(25,+\infty)$}
 
 \item 如图,已知椭圆 $\dfrac{{x}^{2}}{12} + \dfrac{{y}^{2}}{3} = 1$ 的焦点为 ${F}_{1}\text{、}{F}_{2}$ ,点 $P$ 在此椭圆上, 若线段 $P{F}_{1}$ 的中点 $M$ 恰在 $y$ 轴上,则 $\left| {P{F}_{1}}\right|$ 是 $\left| {P{F}_{2}}\right|$ 的 \blkda[2]{7}倍
    \begin{center}
        \includegraphics[width=0.3\linewidth]{tupianahdfuagfua.png}
    \end{center}
    
 \item 若 ${AB}$ 为过椭圆 $\dfrac{{x}^{2}}{{a}^{2}} + \dfrac{{y}^{2}}{{b}^{2}} = 1\left( {a > b > 0}\right)$ 中心的弦, ${F}_{1}$ 为此椭圆的焦点,则 $\bigtriangleup  {F}_{1}{AB}$ 的面积的最大值是\dotfill{(\kh{C})}
    \xx{${ac}$}{${ab}$}{${bc}$ }{${b}^{2}$}



 \item 若直线 $y = {2x}$ 与椭圆 $\dfrac{{x}^{2}}{{a}^{2}} + \dfrac{{y}^{2}}{{b}^{2}} = 1\left( {a > b > 0}\right)$ 的一个交点的横坐标恰为 $c$ ( $c$ 为半焦距), 则 $\dfrac{c}{a}$ 的值为\blkda[2]{$\sqrt{2}-1$}.

 \item 过点 $M\left( {-2,0}\right)$ 的直线 $m$ 与椭圆 $\dfrac{{x}^{2}}{2} + {y}^{2} = 1$ 交于 ${P}_{1},{P}_{2}$ ,线段 ${P}_{1}{P}_{2}$ 的中点为 $P$ ,设直线 $m$ 的斜率为 ${k}_{1}\left( {{k}_{1} \neq  0}\right)$ ,直线 ${OP}$ 的斜率为 ${{k}}_{2}$ ,则 ${k}_{1}{k}_{2}$ 的值为\blkda[2]{$-\dfrac{1}{2}$}.

 \item 已知直线 $l : y = x - 3$ 与椭圆 $C : \dfrac{{y}^{2}}{4m} + \dfrac{{x}^{2}}{m} = 1$ 相交,则 $C$ 的长轴长的取值范围是\blkda[3]{$(\dfrac{12\sqrt{5}}{5},+\infty)$}.

 \item 已知椭圆 $\dfrac{{x}^{2}}{20} + \dfrac{{y}^{2}}{k} = 1\left( {{20} > k > 0}\right)$ 的焦距为 8 ,过椭圆的一个焦点,作垂直于长轴的直线交椭圆于 ${AB}$ 两点,则 $\left| {AB}\right|  =$ \blkda[2]{$\dfrac{4\sqrt{5}}{5}$}.

 \item 已知椭圆 $C : \dfrac{{x}^{2}}{{a}^{2}} + \dfrac{{y}^{2}}{{b}^{2}} = 1\left( {a > b > 0}\right)$ 的焦距为 12,长半轴长为 $4\sqrt{3}$ .
 \begin{enumerate}
     \item 求椭圆 $C$ 的方程；
     \item 直线 $l$ 与椭圆 $C$ 相交于 $A,B$ 两点,若线段 ${AB}$ 的中点坐标为$(-4,1)$,求直线 $l$ 的方程.
 \end{enumerate}
\begin{jieda}
    \begin{enumerate}
        \item $\dfrac{x^2}{48}+\dfrac{y^2}{12}=1$
        \item 该题可设直线，或点差法做。由点差法，$\dfrac{x_中}{a^2}+\dfrac{y_中}{b^2}k=0$，将中点坐标代入可得，$k=1$，因此直线方程为$x-y+5=0$.
    \end{enumerate}
\end{jieda}

 \item 已知 $P\left( {{x}_{0},{y}_{0}}\right)$ 为椭圆 $\dfrac{{x}^{2}}{{a}^{2}} + \dfrac{{y}^{2}}{{b}^{2}} = 1\left( {a > b > 0}\right)$ 内且不在 $x$ 轴上的一点,直线 $l$ 过点 $P$ , 且直线 $l$ 与这个椭圆交于 $A$ 、 $B$ 两点, $A$ 、 $B$ 两点的中点恰好为点 $P$ . 求证: 直线 $l$ 的斜率 $k$ 为 $- \dfrac{{b}^{2}{x}_{0}}{{a}^{2}{y}_{0}}$ .
\begin{jieda}
    设 $ A\text{、}B $ 两点的坐标分别为 $ \left( {{x}_{1},{y}_{1}}\right) \text{、}\left( {{x}_{2},{y}_{2}}\right) $ ,则 $\left\{\begin{array}{l}
         {b}^{2}{x}_{1}^{2} + {a}^{2}{y}_{1}^{2} = {a}^{2}{b}^{2}  \\
         {b}^{2}{x}_{2}^{2} + {a}^{2}{y}_{2}^{2} = {a}^{2}{b}^{2} 
    \end{array} \right.$ , 两式相减,得 $ {b}^{2}\left( {{x}_{1} - {x}_{2}}\right) \left( {{x}_{1} + {x}_{2}}\right)  + {a}^{2}\left( {{y}_{1} - {y}_{2}}\right) \left( {{y}_{1} + {y}_{2}}\right)  = 0 $ . 
    
    又点 $ P $ 不在 $ x $ 轴上,故 $ {x}_{1} \neq  {x}_{2} $ . 而$\left\{\begin{array}{l}
         {x}_{1} + {x}_{2} = 2{x}_{0}  \\
         {y}_{1} + {y}_{2} = 2{y}_{0}
    \end{array} \right.$
    
    所以 $ k = \dfrac{{y}_{1} - {y}_{2}}{{x}_{1} - {x}_{2}} =  - \dfrac{{b}^{2}\left( {{x}_{1} + {x}_{2}}\right) }{{a}^{2}\left( {{y}_{1} + {y}_{2}}\right) } =  - \dfrac{{b}^{2}{x}_{0}}{{a}^{2}{y}_{0}} $ .
\end{jieda}

 \item 已知点 $M$ 为椭圆 $\dfrac{{x}^{2}}{{a}^{2}} + \dfrac{{y}^{2}}{{b}^{2}} = 1$ 上的任一点,它与此椭圆的短轴两端点 ${B}_{1}\text{、}{B}_{2}$ 的连线分别交 $x$ 轴于点 $P\text{、}Q$ . 求证: $\left| {OP}\right|  \cdot  \left| {OQ}\right|$ 为定值.
\begin{jieda}
    设$M(x_0,y_0)$，从而$|OP|=\dfrac{bx_0}{b-y_0}$,$|OQ|=\dfrac{bx_0}{b+y_0}$,从而$|OP||OQ|=a^2$
\end{jieda}


 \item 已知椭圆 $\dfrac{{x}^{2}}{{a}^{2}} + \dfrac{{y}^{2}}{{b}^{2}} = 1\left( {a > b > 0}\right) ,A\text{、}B$ 是此椭圆上的两点,且线段 ${AB}$ 的垂直平分线与 $x$ 轴相交于点 $P\left( {{x}_{0},0}\right)$ . 求证: $- \dfrac{{a}^{2} - {b}^{2}}{a} < {x}_{0} < \dfrac{{a}^{2} - {b}^{2}}{a}$ .
\begin{jieda}
    由点差法，$k_{AB}=-\dfrac{x_中}{a^2}\dfrac{b^2}{y_中}$，从而$k_垂=\dfrac{y_中a^2}{x_中b^2}$，所以垂直平分线所在直线为$y-y_中=k_垂(x-x_中)$，另$y=0$，可得$x_0=x_中(1-\dfrac{b^2}{a^2})$，因为中点在椭圆内，因此$x_中\in(-a,a)$，即证.
\end{jieda}


 \item 已知椭圆 $\dfrac{{x}^{2}}{4} + \dfrac{{y}^{2}}{3} = 1$ 上有不同的两点关于直线 $y = {4x} + m$ 对称,求实数 $m$ 的取值范围.
 \begin{jieda}
     由点差法，中点轨迹方程为$3x-y=0$，又因为中点在直线$y = {4x} + m$ 上，因此联立可得中点坐标为$（-m,-3m）$。因为弦中点在椭圆内，所以$\dfrac{m^2}{4}+\dfrac{9m^2}{3}<1$，从而可得$m\in(-\dfrac{2\sqrt{13}}{13},\dfrac{2\sqrt{13}}{13})$
 \end{jieda}

 \item 已知椭圆 $E : \dfrac{{x}^{2}}{{a}^{2}} + \dfrac{{y}^{2}}{{b}^{2}} = 1\left( {a > b > 0}\right)$ 的长轴长为 $2\sqrt{3}$ ,且点 $M\left( {\sqrt{2},\dfrac{\sqrt{3}}{3}}\right)$ 在椭圆 $E$ 上.
\begin{enumerate}
    \item 求椭圆 $E$ 的方程；
    \item 设直线 $y = {kx} + \sqrt{2}$ 与椭圆 $E$ 相交于不同的两点 $P$ 和 $Q$ ,当 $\left| {PQ}\right|  = \sqrt{3}$ 时,求实数 $k$ 的值.
\end{enumerate}

\begin{jieda}
    \begin{enumerate}
        \item $\dfrac{x^2}{3}+y^2=1$
        \item 直线 $ {PQ} $ 的方程为： $ y = {kx} + \sqrt{2} $
联立 $ \left\{  \begin{array}{l} y = {kx} + \sqrt{2} \\  {x}^{2} + 3{y}^{2} = 3 \end{array}\right. $ ,消去 $ y $ 后,得关于 $ x $ 的一元二次方程
$ {x}^{2} + 3{\left( kx + \sqrt{2}\right) }^{2} = 3 $
化简得 $ \left( {1 + 3{k}^{2}}\right) {x}^{2} + 6\sqrt{2}{kx} + 3 = 0 $ ,
由题意知 $ \Delta  = {\left( 6\sqrt{2}k\right) }^{2} - 4 \times  3 \times  \left( {1 + 3{k}^{2}}\right)  = {36}{k}^{2} - {12} > 0 $ ,解得 $ k <  - \dfrac{\sqrt{3}}{3} $ 或
$ k > \dfrac{\sqrt{3}}{3} $
由韦达定理可得 $ {x}_{1} + {x}_{2} =  - \dfrac{6\sqrt{2}k}{1 + 3{k}^{2}},{x}_{1} \cdot  {x}_{2} = \dfrac{3}{1 + 3{k}^{2}} $ ,
所以 $ \left| {PQ}\right|  = \sqrt{1 + {k}^{2}}\left| {{x}_{1} - {x}_{2}}\right|  = \sqrt{1 + {k}^{2}}\sqrt{\dfrac{{72}{k}^{2}}{{\left( 1 + 3{k}^{2}\right) }^{2}} - \dfrac{12}{1 + 3{k}^{2}}} = \sqrt{3} $ ,
所以 $ \left( {1 + {k}^{2}}\right)  \cdot  \dfrac{{36}{k}^{2} - {12}}{1 + 3{k}^{2}} = 3 $ ,化简得 $ 9{k}^{4} + 6{k}^{2} - {15} = 0 $ ,解得 $ {k}^{2} = 1 $ ,即
$ k =  \pm  1, $
    \end{enumerate}
\end{jieda}
\end{enumerate}

C组：
\begin{enumerate}
    \item 如图,半椭圆 $\dfrac{{x}^{2}}{{a}^{2}} + \dfrac{{y}^{2}}{{b}^{2}} = 1\left( {x \geq  0}\right)$ 与半椭圆 $\dfrac{{y}^{2}}{{b}^{2}} + \dfrac{{x}^{2}}{{c}^{2}} = 1\left( {x < 0}\right)$ 组成的曲线称为“果圆”,其中 ${a}^{2} = {b}^{2} + {c}^{2},a > 0,b > c > 0$ . “果圆”与 $x$ 轴的交点分别为 ${A}_{1}\text{、}{A}_{2}$ ,若在“果圆” $y$ 轴右侧部分上存在点 $P$ 使得 $\angle {A}_{1}P{A}_{2} = \dfrac{\pi }{2}$ ,则 $\dfrac{c}{a}$ 的取值范围为\blkda[3]{$(\dfrac{1}{2},\dfrac{\sqrt{5}-1}{2})$}.
    \begin{center}
        \includegraphics[width=0.25\linewidth]{shouhsd[faubf;k.png}
    \end{center}

    \item 已知椭圆 $G : \dfrac{{x}^{2}}{6} + \dfrac{{y}^{2}}{{b}^{2}} = 1\left( {0 < b < \sqrt{6}}\right)$ 的左、右焦点分别为 ${F}_{1}$ 、 ${F}_{2}$ ,短轴的两个端点分别为 ${B}_{1}$ 、 ${B}_{2}$ ,点 $P$ 在椭圆 $G$ 上,且满足 $\left| {P{F}_{1}}\right|  + \left| {P{F}_{2}}\right|  = \left| {P{B}_{1}}\right|  + \left| {P{B}_{2}}\right|$ . 当 $b$ 变化时,给出下列四个命题：
    \begin{dingautolist}{172}
        \item 点 $P$ 的轨迹关于 $y$ 轴对称
        \item 存在 $b$ 使得椭圆 $G$ 上满足条件的点 $P$ 仅有两个
        \item $\left| {OP}\right|$ 的最小值为 2
        \item $\left| {OP}\right|$ 最大值为 $\sqrt{6}$
    \end{dingautolist}
    其中正确命题的序号是\blkda[2]{\ding{172}\ding{174}}.
\end{enumerate}

\section*{双曲线的标准方程}

A 组：
\begin{enumerate}
    \item 设双曲线 $\dfrac{{x}^{2}}{5} - \dfrac{{y}^{2}}{4} = 1$ 的左、右焦点分别为 ${F}_{1}$ 、 ${F}_{2}$ ,点 $P$ 在双曲线的右支上,则 $\left| {P{F}_{1}}\right|  - \left| {P{F}_{2}}\right|  =$ \blkda{$2\sqrt{5}$}.
    
    \item 虚轴长等于实半轴长,且焦距长为 4 的双曲线的标准方程是\blkda[6]{$\dfrac{x^2}{\dfrac{16}{5}}-\dfrac{y^2}{\dfrac{4}{5}} = 1$或$\dfrac{y^2}{\dfrac{16}{5}}-\dfrac{x^2}{\dfrac{4}{5}} = 1$}.
    \begin{jieda}
        $a = 2b, c = 2, a^2+b^2 = c^2$，故$a^2 = \dfrac{16}{5}, b^2 = \dfrac{4}{5}$
    \end{jieda}
    
    \item 过两点 ${P}_{1}\left( {3, - 4\sqrt{2}}\right) ,{P}_{2}\left( {\dfrac{9}{4},5}\right)$ 的双曲线的标准方程是\blkda[3]{$\dfrac{y^2}{16} - \dfrac{x^2}{9} = 1$}.
    \begin{jieda}
        设所求双曲线为$Ax^2+By^2 = 1(AB<0)$代入两点坐标有$\left\{\begin{array}{l}
             9A - 32B = 1 \\
             \dfrac{81}{16}A - 25B = 1
        \end{array}\right.$，解得$\left\{\begin{array}{l}
             A = -\dfrac{1}{9} \\
             B =  \dfrac{1}{16}
        \end{array}\right.$，故双曲线标准方程为$\dfrac{y^2}{16} - \dfrac{x^2}{9} = 1$
    \end{jieda}
    
    \item 已知 $M$ 为平面内一动点,设命题甲: 存在两个定点 ${F}_{1},{F}_{2}$ 使得 $\left|{M{F}_{1}}\right| - \left|{M{F}_{2}}\right|$ 是定值,命题乙: $M$ 的轨迹是双曲线,则命题甲是命题乙的\blkda[3]{必要非充分}条件
    
    \item 两定点 ${F}_{1}\left( {-5,0}\right) ,{F}_{2}\left( {5,0}\right)$ ,动点 $M\left( {x,y}\right)$ 满足 $\left| {M{F}_{1}}\right|  - \left| {M{F}_{2}}\right|  = 8$ ,则动点 $M$ 的轨迹方程为\blkda[4]{$\dfrac{{x}^{2}}{16} - \dfrac{{y}^{2}}{9} = 1\left( {x > 0}\right)$}.
    
    \item “ $m < 1$ ” 是 “方程 $\dfrac{{x}^{2}}{m + 2} + \dfrac{{y}^{2}}{m - 1} = 1$ 表示双曲线” 的\dotfill{(\kh{B})}
    \xx{充分不必要条件}{必要不充分条件}{充分必要条件}{既不充分也不必要条件}
    \begin{jieda}
        $-2 < m < 1 \Longleftrightarrow $方程 $\dfrac{{x}^{2}}{m + 2} + \dfrac{{y}^{2}}{m - 1} = 1$ 表示双曲线
    \end{jieda}
    
    \item 若点 $P$ 是双曲线 $C : \dfrac{{x}^{2}}{4} - \dfrac{{y}^{2}}{12} = 1$ 上一点, ${F}_{1},{F}_{2}$ 分别为 $C$ 的左、右焦点, $\left| {P{F}_{1}}\right|  = 9$ , 则 $\left| {P{F}_{2}}\right|  =$ \blkda{5 或 13}.
    \begin{jieda}
        设双曲线 $C$ 的半焦距为 $c > 0$ ,则 $a = 2,b = 2\sqrt{3},c = \sqrt{{a}^{2} + {b}^{2}} = 4$
        
        因为 $\left|{P{F}_{1}}\right| - \left|{P{F}_{2}}\right| = \left| {9 - \left|{P{F}_{2}}\right|}\right|  = 4$，解得 $\left| {P{F}_{2}}\right|  = 5$ 或 $\left| {P{F}_{2}}\right|  = {13}$，经检验均符合题意.
    \end{jieda}
\end{enumerate}

B组：
\begin{enumerate}
    \item 若双曲线 $C$ 以椭圆 $\dfrac{{x}^{2}}{3} + \dfrac{{y}^{2}}{4} = 1$ 的焦点为实轴顶点,长轴顶点为焦点,则双曲线 $C$ 的方程是\blkda[3]{${y}^{2} - \dfrac{{x}^{2}}{3} = 1$}.
    
    \item 经过两点 $\left( {-3,6\sqrt{2}}\right)$ , $\left( {2,3\sqrt{3}}\right)$ 的双曲线的标准方程为\blkda[3]{${x}^{2} - \dfrac{{y}^{2}}{9} = 1$}.
    \begin{jieda}
        设双曲线的方程为: $m{x}^{2} + n{y}^{2} = 1\left( {{mn} < 0}\right)$ ,故 $\left\{  \begin{array}{l} {9m} + {72n} = 1 \\  {4m} + {27n} = 1 \end{array}\right.$ ,解得 $\left\{  \begin{array}{l} m = 1 \\  n =  - \dfrac{1}{9} \end{array}\right.$ ,故双曲线的方程为: ${x}^{2} - \dfrac{{y}^{2}}{9} = 1$
    \end{jieda}

    \item 若点 $P$ 为椭圆 $\dfrac{{x}^{2}}{5} + \dfrac{{y}^{2}}{9} = 1$ 与双曲线 ${y}^{2} - \dfrac{{x}^{2}}{3} = 1$ 在第二象限的交点, ${F}_{1}\text{ 、 }{F}_{2}$ 分别为双曲线的上、下焦点,则 $\left| {P{F}_{1}}\right|  =$ \blkda{2}, $\left| {P{F}_{2}}\right|  =$ \blkda{4}.

    \item 到点 ${F}_{1}\left( {-4,0}\right)$ , ${F}_{2}\left( {4,0}\right)$ 的距离的差的绝对值等于 6 的点的双曲线的标准方程为\blkda[3]{$\dfrac{{x}^{2}}{9} - \dfrac{{y}^{2}}{7} = 1$}.
    \begin{jieda}
        由题意可设双曲线方程为 $\dfrac{{x}^{2}}{{a}^{2}} - \dfrac{{y}^{2}}{{b}^{2}} = 1,\left( {a > 0,b > 0}\right)$ ,焦距设为 ${2c}$ , 由题意可知所求双曲线的两焦点为 ${F}_{1}\left( {-4,0}\right) ,{F}_{2}\left( {4,0}\right)$ ,故 ${2c} = 8,c = 4$ , 又双曲线上的点到点 ${F}_{1}\left( {-4,0}\right) ,{F}_{2}\left( {4,0}\right)$ 的距离的差的绝对值等于 6, 故 ${2a} = 6,a = 3$ ,所以 ${b}^{2} = {c}^{2} - {a}^{2} = 7$ , 故双曲线标准方程为 $\dfrac{{x}^{2}}{9} - \dfrac{{y}^{2}}{7} = 1$ . 
    \end{jieda}

    \item 已知点 ${F}_{1}\text{ 、 }{F}_{2}$ 是双曲线 $\dfrac{{x}^{2}}{9} - \dfrac{{y}^{2}}{16} = 1$ 的两个焦点,点 $P$ 在此双曲线上,且满足 $\left| {P{F}_{1}}\right|  \cdot  \left| {P{F}_{2}}\right|  = {32}$ ,则 $\angle {F}_{1}P{F}_{2} =$ \blkda{$\dfrac{\pi}{2}$}.
    \begin{jieda}
        在$\triangle F_1PF_2$中列余弦定理有$|F_1F_2|^2 = |PF_1|^2+|PF_2|^2 - 2|PF_1|\cdot |PF_2| cos\angle {F}_{1}P{F}_{2} = \left(|PF_1|-|PF_2|\right)^2-2|PF_1|\cdot |PF_2| (cos\angle {F}_{1}P{F}_{2} - 1)$

        代入具体数值有$100 = 36 - 64 \times (cos\angle {F}_{1}P{F}_{2} - 1)$，故$cos\angle {F}_{1}P{F}_{2} = 0$，即$\angle {F}_{1}P{F}_{2} = \dfrac{\pi}{2}$
    \end{jieda}

    \item 已知双曲线 $\dfrac{{x}^{2}}{40} - \dfrac{{y}^{2}}{9} = 1$ 上一点 $M$ 满足 $M{F}_{1} \bot  M{F}_{2}$ ,其中 ${F}_{1}\text{ 、 }{F}_{2}$ 是双曲线的两个焦点,则 $\bigtriangleup  M{F}_{1}{F}_{2}$ 的面积是\blkda{9}.
    \begin{jieda}
        $S = 9 \times cot\left(\dfrac{\pi}{4}\right) = 9$.
    \end{jieda}

    \item 已知椭圆 ${C}_{1} : \dfrac{{x}^{2}}{m + 2} - \dfrac{{y}^{2}}{n} = 1$ 与双曲线 ${C}_{2} : \dfrac{{x}^{2}}{m} + \dfrac{{y}^{2}}{n} = 1$ 有相同的焦点,则椭圆 ${C}_{1}$ 的离心率 ${e}_{l}$ 的取值范围为\blkda{$\left( {\dfrac{\sqrt{2}}{2},1}\right)$}.
    \begin{jieda}
        椭圆 ${C}_{1} : \dfrac{{x}^{2}}{m + 2} + \dfrac{{y}^{2}}{-n} = 1$ ,双曲线 ${C}_{2} : \dfrac{{x}^{2}}{m} - \dfrac{{y}^{2}}{-n} = 1$ ,
        
        因为椭圆 ${C}_{1}$ 与双曲线 ${C}_{2}$ 有相同的焦点，所以 $m + 2 - \left( {-n}\right)  = m - n$ ,解得 $n =  - 1$ .
        
        所以椭圆 ${C}_{1} : \dfrac{{x}^{2}}{m + 2} + {y}^{2} = 1,{e}^{2} = \dfrac{m + 2 - 1}{m + 2} = 1 - \dfrac{1}{m + 2}$ .
        
        因为 $m > 0,m + 2 > 2$ ,所以 $0 < \dfrac{1}{m + 2} < \dfrac{1}{2}, - \dfrac{1}{2} <  - \dfrac{1}{m + 2} < 0$，所以 $\dfrac{1}{2} < 1 - \dfrac{1}{m + 2} < 1$ ,即 $\dfrac{1}{2} < {e}^{2} < 1,\dfrac{\sqrt{2}}{2} < e < 1$ .

    \end{jieda}

    \item 已知直线 $l : y = {2x} - {10}$ 与双曲线 $\dfrac{{x}^{2}}{{a}^{2}} - \dfrac{{y}^{2}}{{b}^{2}} = 1\left( {a > 0,b > 0}\right)$ 的一条渐近线平行,且经过双曲线的一个焦点,则双曲线的标准方程为\blkda[3]{$\dfrac{{x}^{2}}{5} - \dfrac{{y}^{2}}{20} = 1$}.
    \begin{jieda}
        对于直线 $l : y = {2x} - {10}$ ,令 $y = 0$ ,解得 $x = 5$ ,所以双曲线的一个焦点为$(5,0)$，双曲线 $\dfrac{{x}^{2}}{{a}^{2}} - \dfrac{{y}^{2}}{{b}^{2}} = 1\left( {a > 0,b > 0}\right)$ 的渐近线为 $y =  \pm  \dfrac{b}{a}x$ ,
        
        依题意 $\dfrac{b}{a} = 2$ 且 $c = 5$ 又 ${c}^{2} = {a}^{2} + {b}^{2}$ ,解得 ${a}^{2} = 5,{b}^{2} = {20}$，所以双曲线方程为 $\dfrac{{x}^{2}}{5} - \dfrac{{y}^{2}}{20} = 1$
    \end{jieda}

    \item 已知双曲线 $\dfrac{{x}^{2}}{4} - \dfrac{{y}^{2}}{5} = 1$ 的左焦点为 $F$ ,点 $P$ 在双曲线上且在 $x$ 轴上方,若线段 ${PF}$ 的中点在以 $O$ 为圆心, $\left| {OF}\right|$ 为半径的圆上,则直线 ${PF}$ 的斜率为\blkda[3]{$\dfrac{\sqrt{11}}{5}$ 或 $\sqrt{35}$}.
    \begin{jieda}
        由双曲线 $\dfrac{{x}^{2}}{4} - \dfrac{{y}^{2}}{5} = 1$ 可知 $a = 2,b = \sqrt{5},c = \sqrt{{a}^{2} + {b}^{2}} = 3$ ,
        
        设线段 ${PF}$ 的中点为 $M$ ,双曲线的右焦点为 ${A}$，则 $\left| {OM}\right|  = 3,\left| {AF}\right|  = 6$
        \begin{center}
            \includegraphics[width=0.3\textwidth]{bodyxb/src/bo_d3eri7s601uc738iqr70_2_240_1616_444_409_0.jpg}
        \end{center}
        \begin{dingautolist}{172}
            \item 若点 $P$ 在第一象限,则 $\left| {PA}\right|  = 2\left| {OM}\right|  = 6,\left| {PF}\right|  = \left| {PA}\right|  + 4 = {10}$，可得 $\cos \angle {PFA} = \dfrac{{\left| PA\right| }^{2} + {\left| AF\right| }^{2} - {\left| AP\right| }^{2}}{2\left| {PA}\right|  \cdot  \left| {AF}\right| } = \dfrac{{10}^{2} + {6}^{2} - {6}^{2}}{2 \times  {10} \times  6} = \dfrac{5}{6}$，则 $\sin \angle {PFA} = \sqrt{1 - {\cos }^{2}\angle {PFA}} = \dfrac{\sqrt{11}}{6}$
            
            可得 $\tan \angle {PFA} = \dfrac{\sin \angle {PFA}}{\cos \angle {PFA}} = \dfrac{\sqrt{11}}{5}$ ,所以直线 ${PF}$ 的斜率为 $\dfrac{\sqrt{11}}{5}$ .
            
            \item 若点 $P$ 在第二象限，则$\left| {PA}\right|  = 2\left| {OM}\right|  = 6,\left| {PF}\right|  = \left| {PA}\right|  - 4 = 2$

            ，可得 $\cos \angle {PFA} = \dfrac{{\left| PA\right| }^{2} + {\left| AF\right| }^{2} - {\left| AP\right| }^{2}}{2\left| {PA}\right|  \cdot  \left| {AF}\right| } = \dfrac{2^{2} + {6}^{2} - {6}^{2}}{2 \times  2 \times  6} = \dfrac{1}{6}$，$\sin \angle {PFA} = \sqrt{1 - {\cos }^{2}\angle {PFA}} = \dfrac{\sqrt{35}}{6}$

            可得 $\tan \angle {PFA} = \dfrac{\sin \angle {PFA}}{\cos \angle {PFA}} = \sqrt{35}$ ,所以直线 ${PF}$ 的斜率为 $\sqrt{35}$ .
        
        \end{dingautolist}
        \ps 直接求解$\left\{  \begin{array}{l} \dfrac{{\left( s - 3\right) }^{2}}{4} + \dfrac{{t}^{2}}{4} = 9 \\  \dfrac{{s}^{2}}{4} - \dfrac{{t}^{2}}{5} = 1  \\  t > 0 \end{array}\right.$即可得$P$在一二两个象限的坐标$(s, t)$.
    \end{jieda}

    \item 若过双曲线 $\dfrac{{x}^{2}}{16} - \dfrac{{y}^{2}}{9} = 1$ 的左焦点 ${{F}}_{1}$ 的弦 ${{AB}}$ 长为 6,则 $\bigtriangleup  {AB}{F}_{2}$ ( ${{F}}_{2}$ 为右焦点) 的周长是\blkda{28}.
    \begin{jieda}
        根据双曲线的定义可知$\left\{\begin{array}{l}
             |AF_2| - |AF_1| = 8 \\
             |BF_2| - |BF_1| = 8
        \end{array}\right.$，故周长为$|AB|+|AF_2|+|BF_2| = 6+6+8+8 = 28$.
    \end{jieda}

    \item 在 $\bigtriangleup {ABC}$ 中,已知 $A\left( {-4,0}\right) \text{ 、 }B\left( {4,0}\right)$ 两点. 若 $\sin A - \sin B = \dfrac{1}{2}\sin C$ ,则顶点 $C$ 的轨迹方程是\blkda[4]{$\dfrac{{x}^{2}}{4} - \dfrac{{y}^{2}}{12} = 1\left( {x <  - 2}\right)$}.

    \item 已知椭圆 $\dfrac{{x}^{2}}{{m}^{2}} + \dfrac{{y}^{2}}{{n}^{2}} = 1\left( {\left| m\right|  > \left| n\right| }\right)$ 和双曲线 $\dfrac{{x}^{2}}{{a}^{2}} - \dfrac{{y}^{2}}{{b}^{2}} = 1$ 有相同的两焦点 ${F}_{1}\text{ 、 }{F}_{2}$ : 若点 $P$ 为两曲线的一个交点,则 $\left| {P{F}_{1}}\right|  \cdot  \left| {P{F}_{2}}\right|  =$\dotfill{(\kh{A})}
    \xx{${m}^{2} - {a}^{2}$}{$\dfrac{1}{2}\left( {{m}^{2} - {a}^{2}}\right)$}{$m - n$}{${a}^{2} - {m}^{2}$}
    \begin{jieda}
        根据题意可知$\left| {P{F}_{1}}\right|  +  \left| {P{F}_{2}}\right| = 2m, \left|\left| {P{F}_{1}}\right|  -  \left| {P{F}_{2}}\right|\right| = 2a$
        
        故$\left| {P{F}_{1}}\right|  \cdot  \left| {P{F}_{2}}\right| = \dfrac{(2m)^2 - (2a)^2}{4}$
    \end{jieda}
    
    \item 已知动点 $P$ 到两定点 ${F}_{1}\left( {0, - 5}\right) ,{F}_{2}\left( {0,5}\right)$ 的距离之差的绝对值等于 ${2a}$ ,求动点 $P$ 的轨迹.
    \begin{jieda}
        \begin{dingautolist}{172}
            \item $a = 0$，轨迹为$y = 0$
            \item $0 < a < 5$，轨迹为$\dfrac{{y}^{2}}{{a}^{2}} - \dfrac{{x}^{2}}{{25} - {a}^{2}} = 1$
            \item $a = 5$，轨迹为$x = 0\left( {y \leq   - 5\text{ 或 }y \geq  5}\right)$
            \item $a > 5$ ,不存在轨迹
        \end{dingautolist}

    \end{jieda}

    \item 设双曲线 $C$ 与椭圆 $\dfrac{{x}^{2}}{m} + \dfrac{{y}^{2}}{20} = 1$ 有一个公共点为 ${P}\left( {1, - \sqrt{15}}\right)$ ,且它们有公共的焦点,求双曲线 $C$ 的方程.
    \begin{jieda}
        将点$P$坐标代入椭圆方程可得$m = 4$，故椭圆的焦点为$(0, 4)$和$(0, -4)$，设双曲线的方程为$\dfrac{x^2}{a^2}-\dfrac{y^2}{16-a^2} = 1$，代入点$P$坐标，可得双曲线标准方程为$\dfrac{{y}^{2}}{12} - \dfrac{{x}^{2}}{4} = 1$
    \end{jieda}

    \item 已知圆 ${\left( x + 4\right) }^{2} + {y}^{2} = {25}$ 圆心为 ${M}_{1}$ ,圆 ${\left( x - 4\right) }^{2} + {y}^{2} = 1$ 的圆心为 ${M}_{2}$ ,一动圆 ${M}$ 与这两个圆都外切,求动圆圆心的轨迹方程.
    \begin{jieda}
        $\dfrac{{x}^{2}}{4} - \dfrac{{y}^{2}}{12} = 1\left( {x > 0}\right)$
    \end{jieda}

    \item 已知双曲线 ${C}$ 的方程为 ${x}^{2} - \dfrac{{y}^{2}}{2} = 1$ ,试判断过点 $P\left( {1,1}\right)$ 能否作一条直线 $l$ 与双曲线交于 $A\text{ 、 }B$ 两点,且 $P$ 为 ${AB}$ 的中点,并说明理由.
    \begin{jieda}
        设双曲线上两点$A(x_1, y_1), B(x_2, y_2)$，$P$ 为 ${AB}$ 的中点，则有$\left\{\begin{array}{l}
             {x_1}^{2} - \dfrac{{y_1}^{2}}{2} = 1  \\
             {x_2}^{2} - \dfrac{{y_2}^{2}}{2} = 1
        \end{array}\right.$，两式作差有$(x_1+x_2)(x_1-x_2) - \dfrac{(y_1+y_2)(y_1-y_2)}{2} = 0$，根据$P$为中点有$x_1+x_2 = 2, y_1+y_2 = 2$，故$2(x_1-x_2) - (y_1-y_2) = 0$，即$k_{AB} = 2$，故$AB: y = 2x-1$

        联立直线和双曲线方程有$\left\{\begin{array}{l}
             y = 2x-1  \\
             {x}^{2} - \dfrac{{y}^{2}}{2} = 1
        \end{array}\right.$，消去$y$可得$2x^2-(2x-1)^2 = 2$，即$-2x^2+4x-3=0$，方程无解，故不存在.
    \end{jieda}

\end{enumerate}


\section*{双曲线的性质}

\subsection*{双曲线的性质 (1)}

A组：
\begin{enumerate}
    \item 双曲线 $\dfrac{{x}^{2}}{4} - {y}^{2} = 1$ 的顶点坐标\blkda[3]{$\left( {-2,0}\right) ,\left( {2,0}\right)$},渐近线方程\blkda[3]{$y =  \pm  \dfrac{1}{2}x$}.
    
    \item 经过点$(-2,3)$,且对称轴都在坐标轴上的等轴双曲线的标准方程是\blkda[3]{$\dfrac{{y}^{2}}{5} - \dfrac{{x}^{2}}{5} = 1$}.
    \begin{jieda}
        由题意,所求双曲线为等轴双曲线,可得双曲线的方程为 ${x}^{2} - {y}^{2} = \lambda \left( {\lambda  \neq  0}\right)$ ,
        
        因为所求双曲线过点$(-2,3)$,可得 ${\left( -2\right) }^{2} - {3}^{2} = \lambda$ ,解得 $\lambda  =  - 5$，所以所求双曲线的方程为 $\dfrac{{y}^{2}}{5} - \dfrac{{x}^{2}}{5} = 1$ 
    \end{jieda}
    
    \item 与双曲线 $\dfrac{{x}^{2}}{4} - \dfrac{{y}^{2}}{3} = 1$ 有相同焦点,且经过点 $\left( {2,\dfrac{3\sqrt{3}}{2}}\right)$ 的椭圆的标准方程为\blkda[3]{$\dfrac{{x}^{2}}{16} + \dfrac{{y}^{2}}{9} = 1$}.

    \item 已知双曲线 $C : m{x}^{2} - {y}^{2} = 1\left( {m > 0}\right)$ 的一条渐近线方程为 ${mx} + \sqrt{3}y = 0$ ,则 $C$ 的焦距为\blkda{$\dfrac{4\sqrt{3}}{3}$}.
    \begin{jieda}
        $C : m{x}^{2} - {y}^{2} = 1\left( {m > 0}\right)$的渐近线为$y = \pm \sqrt{m}x$，故$\sqrt{m} = \dfrac{m}{\sqrt{3}}$，故$m = 3$，因此双曲线为$\dfrac{x^2}{\dfrac{1}{3}} - {y}^{2} = 1$，焦距为$2\sqrt{1+\dfrac{1}{3}} = \dfrac{4\sqrt{3}}{3}$
    \end{jieda}

    \item 已知双曲线 $\dfrac{{x}^{2}}{9} - \dfrac{{y}^{2}}{16} = 1$ 与 $\dfrac{{y}^{2}}{16} - \dfrac{{x}^{2}}{9} = 1$ ,下列说法正确的是\dotfill{(\kh{C})}
    \xx{两个双曲线有公共顶点}{两个双曲线有公共焦点}{两个双曲线有公共渐近线}{两个双曲线的离心率相等}
    \begin{jieda}
        双曲线 $\dfrac{{x}^{2}}{9} - \dfrac{{y}^{2}}{16} = 1$ 的焦点和顶点都在 $x$ 轴上，而双曲线 $\dfrac{{y}^{2}}{16} - \dfrac{{x}^{2}}{9} = 1$ 的焦点和顶点都在 $y$ 轴上,故 ${A}$ 、 ${B}$ 错误;
        
        双曲线 $\dfrac{{x}^{2}}{9} - \dfrac{{y}^{2}}{16} = 1$ 的渐近线方程为 $y =  \pm  \dfrac{4}{3}x$ ,双曲线 $\dfrac{{y}^{2}}{16} - \dfrac{{x}^{2}}{9} = 1$ 的渐近线方程为 $y =  \pm  \dfrac{4}{3}x$ ,故 ${C}$ 正确;
        
        双曲线 $\dfrac{{x}^{2}}{9} - \dfrac{{y}^{2}}{16} = 1$ 的离心率 ${e}_{1} = \sqrt{\dfrac{9 + {16}}{9}} = \dfrac{5}{3}$，而双曲线 $\dfrac{{y}^{2}}{16} - \dfrac{{x}^{2}}{9} = 1$ 的离心率 ${e}_{2} = \sqrt{\dfrac{{16} + 9}{16}} = \dfrac{5}{4}$ ,故 ${D}$ 错误.
    \end{jieda}
    
    \item 已知双曲线 ${C}_{1} : {x}^{2} + \dfrac{{y}^{2}}{m} = 1\left( {m \neq  0}\right)$ 与 ${C}_{2} : \dfrac{{x}^{2}}{2} - \dfrac{{y}^{2}}{2} = 1$ 共焦点,则 ${C}_{1}$ 的渐近线方程为\dotfill{(\kh{D})}
    \xx{$x \pm  y = 0$}{$\sqrt{2}x \pm  y = 0$}{$x \pm  \sqrt{3}y = 0$}{$\sqrt{3}x \pm  y = 0$}
    \begin{jieda}
        双曲线 ${C}_{2} : \dfrac{{x}^{2}}{2} - \dfrac{{y}^{2}}{2} = 1$ 中 ${a}^{2} = 2\text{ 、 }{b}^{2} = 2$ ,所以 ${c}^{2} = {a}^{2} + {b}^{2} = 4$ ,即 ${C}_{2}$ 焦点坐标为 $\left( {\pm 2,0}\right)$ ,因为双曲线 ${C}_{1} : {x}^{2} - \dfrac{{y}^{2}}{-m} = 1\left( {m \neq  0}\right)$ 与 ${C}_{2} : \dfrac{{x}^{2}}{2} - \dfrac{{y}^{2}}{2} = 1$ 共焦点,所以 $1 + \left( {-m}\right)  = 4$ ,
        
        解得 $m =  - 3$ ,所以双曲线 ${C}_{1} : {x}^{2} - \dfrac{{y}^{2}}{3} = 1$ ,则 ${C}_{1}$ 的渐近线方程为 $\sqrt{3}x \pm  y = 0$
    \end{jieda}

\end{enumerate}
B 组：
\begin{enumerate}
    \item 已知双曲线的方程为 $\dfrac{{x}^{2}}{9} - \dfrac{{y}^{2}}{16} = 1$ ,点 ${F}_{1},{F}_{2}$ 是其左右焦点, ${A}$ 是圆 ${x}^{2} + {\left( y - 5\right) }^{2} = 4$ 上的一点,点 $M$ 在双曲线的右支上,则 $\left| {M{F}_{1}}\right|  + \left| {MA}\right|$ 的最小值是\blkda{$5\sqrt{2} + 4$}.
    \begin{jieda}
        由题意可得, ${c}^{2} = 9 + {16} = {25}$ ,即 $c = 5$ ,则 ${F}_{1},{F}_{2}$ 的坐标分别为 $\left( {-5,0}\right) ,\left( {5,0}\right)$，由双曲线的定义,得 $\left| {M{F}_{1}}\right|  - \left| {M{F}_{2}}\right|  = {2a} = 6$ 
        
        \begin{center}
        \includegraphics[width=0.2\textwidth]{bodyxb/src//bo_d3esa9s601uc738iqt8g_2_241_522_348_386_0.jpg}
        \end{center}

        又 ${A}$ 是圆 ${x}^{2} + {\left( y - 5\right) }^{2} = 4$ 上的点,设圆的圆心为 $C\left( {0,5}\right)$ ,半径为 $r = 2$
        
        由图可知, $\left| {CM}\right|  + \left| {M{F}_{2}}\right|  \geq  \left| {C{F}_{2}}\right|$，所以 $\left| {M{F}_{1}}\right|  + \left| {MA}\right|  = \left| {M{F}_{2}}\right|  + {2a} + \left| {MA}\right|  \geq  \left| {M{F}_{2}}\right|  + \left| {CM}\right|  + {2a} + {CA} \geq  C{F}_{2} + 6 - 2 = 5\sqrt{2} + 4$ ,
        
        当且仅当 $C$ 、 ${A}$ 、 $M$ 、 ${F}_{2}$ 四点共线( ${A}$ 、 $M$ 在 $C{F}_{2}$ 之间)时取等号,
        
        所以 $\left| {M{F}_{1}}\right|  + \left| {MA}\right|$ 的最小值为 $5\sqrt{2} + 4$
    \end{jieda}

    \item 在相距 1000 米的 ${F}_{1}\text{ 、 }{F}_{2}$ 两地听到炮声的时间差是 2 秒(声速:340 米 / 秒). 若以 ${F}_{1}$ 与 ${F}_{2}$ 的中点为原点, ${F}_{1}$ 与 ${F}_{2}$ 两地所在直线为 $x$ 轴,建立平面直角坐标系, 则炮的位置所在的方程是\blkda[4]{$\dfrac{{x}^{2}}{115600} - \dfrac{{y}^{2}}{134400} = 1$}.

    \item 在平面直角坐标系 ${xOy}$ 中,已知 $\bigtriangleup  {ABC}$ 的顶点 $A\left( {-6,0}\right)$ 和 $C\left( {6,0}\right)$ ,若顶点 $B$ 在双曲线 $\dfrac{{x}^{2}}{25} - \dfrac{{y}^{2}}{11} = 1$ 的左支上,则 $\dfrac{\sin A - \sin C}{\sin B} =$ \blkda{$\dfrac{5}{6}$}.
    \begin{jieda}
        $\dfrac{\sin A - \sin C}{\sin B} = \dfrac{|BC| - |AB|}{|AC|} = \dfrac{10}{12} = \dfrac{5}{6}$.
    \end{jieda}

    \item 已知圆 $O$ 过双曲线 $\dfrac{{x}^{2}}{9} - \dfrac{{y}^{2}}{16} = 1$ 的一个顶点和一个焦点. 若圆心 $O$ 在此双曲线上, 则圆心 $O$ 到双曲线中心的距离是\blkda{$\dfrac{16}{3}$}.
    \begin{jieda}
        \begin{dingautolist}{172}
            \item 若过同侧顶点和焦点，不妨设过右顶点和右焦点，则圆心$O$在中垂线$x = \dfrac{a+c}{2} = 4$上，则$O\left(4, \dfrac{4\sqrt{7}}{3}\right)$，因此圆心到原点距离为$\dfrac{16}{3}$
            \item 若过异侧的顶点和焦点，不妨设过右顶点和左焦点，则圆心$O$在中垂线$x = \dfrac{a-c}{2} = -1$上. 因为中垂线与双曲线无交点，所以不满足要求.
        \end{dingautolist}
    \end{jieda}

    \item 设 $P$ 是双曲线 ${x}^{2} - \dfrac{{y}^{2}}{3} = 1$ 的右支上的动点, $F$ 为双曲线的右焦点,已知 $A\left( {3,1}\right)$ , 则 $\left| {PA}\right|  + \left| {PF}\right|$ 的最小值为\blkda{$\sqrt{26} - 2$}.
    \begin{jieda}
        设左焦点为$F'$，则$|PF| = |PF'| - 2$，$|PF'|+|PA| \geq |AF'| = \sqrt{26}$，故$\left| {PA}\right|  + \left| {PF}\right| \geq \sqrt{26} - 2$，当$A,F',P$三点共线时取得最小值.
    \end{jieda}

    \item 点 $Q$ 是双曲线 $C : \dfrac{{x}^{2}}{16} - \dfrac{{y}^{2}}{4} = 1$ 上一动点,过 $Q$ 作圆 $D : {\left( x - 6\right) }^{2} + {y}^{2} = 1$ 的两条切线,切点为 $A$ ,$B$ ,则 $\left| {QA}\right|$ 的最小值为\blkda{$\dfrac{\sqrt{55}}{5}$}.
    \begin{jieda}
        由题知: 设 $Q\left( {{x}_{0},{y}_{0}}\right) ,{x}_{0} \geq  4$ ,则 $\dfrac{{x}_{0}^{2}}{16} - \dfrac{{y}_{0}^{2}}{4} = 1$，由于 $\bigtriangleup {QAD}$ 是直角三角形,且 $\left| {AD}\right|  = 1$ ,所以当 $\left| {QD}\right|$ 取得最小值时, $\left| {QA}\right|  = \sqrt{{\left| QD\right| }^{2} - 1}$ 取得最小值,则 $\left| {QD}\right|  = \sqrt{{\left( {x}_{0} - 6\right) }^{2} + {y}_{0}^{2}} = \sqrt{{\left( {x}_{0} - 6\right) }^{2} + \dfrac{{x}_{0}^{2}}{4} - 4} = \sqrt{\dfrac{5{x}_{0}^{2}}{4} - {12}{x}_{0} + {32}} = \sqrt{\dfrac{5}{4}{\left( {x}_{0} - \dfrac{24}{5}\right) }^{2} + \dfrac{16}{5}} \geq  \dfrac{4\sqrt{5}}{5}$ ,当 ${x}_{0} = \dfrac{24}{5}$ 时,等号成立,故 $\left| {QA}\right|  = \sqrt{{\left| QD\right| }^{2} - 1} \geq  \dfrac{\sqrt{55}}{5}$
        \begin{center}
            \includegraphics[width=0.5\linewidth]{bodyxb//src/shuangquxianxingzhi1B6.png}
        \end{center}
    \end{jieda}

    \item 若连接双曲线 $\dfrac{{x}^{2}}{{a}^{2}} - \dfrac{{y}^{2}}{{b}^{2}} = 1$ 和 $\dfrac{{y}^{2}}{{b}^{2}} - \dfrac{{x}^{2}}{{a}^{2}} = 1$ 的四个顶点的四边形面积为 ${S}_{1}$ ,连接其四个焦点的四边形的面积为 ${S}_{2}$ ,则 $\dfrac{{S}_{1}}{{S}_{2}}$ 的最大值为\blkda{$\dfrac{1}{2}$}.

    \begin{jieda}
        $S_1 = (2a)\times(2b)\div 2 = 2ab$，$S_2 = (2c) \times (2c) \div 2 = 2c^2$，$\dfrac{S_1}{S_2} = \dfrac{2ab}{2c^2} = \dfrac{ab}{a^2+b^2} \leq \dfrac{1}{2}$（$a = b$时取等） 
    \end{jieda}


    \item 如图,已知 ${OA}$ 是双曲线的实半轴, ${OB}$ 是虚半轴, $F$ 为焦点,且 $\angle {BAO} = {30}^{ \circ  }$ . 若 ${S}_{\bigtriangleup {ABF}} = \dfrac{1}{2}\left( {6 - 3\sqrt{3}}\right)$ ,则该双曲线的离心率为\blkda{$\dfrac{2\sqrt{3}}{3}$}.
    \begin{center}
    \includegraphics[width=0.2\textwidth]{bodyxb/src/bo_d3esa6c601uc738iqt7g_1_171_1137_303_284_0.jpg}
    \end{center}
    \begin{jieda}
        $\angle{BAO} = \ang{30} \Longleftrightarrow a = \sqrt{3}b, c = 2b, e = \dfrac{c}{a} = \dfrac{2\sqrt{3}}{3}$
        
        \ps 无需${S}_{\bigtriangleup {ABF}}$即可确定离心率，但根据${S}_{\bigtriangleup {ABF}}$可以确定双曲线的标准方程，${S}_{\bigtriangleup {ABF}} = (c-a)\times b \div 2 = \dfrac{2-\sqrt{3}}{2}b^2 = \dfrac{1}{2}\left( {6 - 3\sqrt{3}}\right)$，故$b = \sqrt{3}, a = 3, c = 2\sqrt{3}$
    \end{jieda}

    \item 已知圆 ${C}_{1} : {\left( x - 4\right) }^{2} + {y}^{2} = 4$ ,圆 ${C}_{2} : {\left( x + 4\right) }^{2} + {y}^{2} = {16}$ ,若动圆 $M$ 与圆 ${C}_{1},{C}_{2}$ 均外切,则动圆圆心的轨迹方程为\blkda[4]{${x}^{2} - \dfrac{{y}^{2}}{15} = 1\left( {x \geq  1}\right)$}.
    \begin{jieda}
        圆 ${C}_{1} : {\left( x - 4\right) }^{2} + {y}^{2} = 4$ 的圆心为 ${C}_{1}\left( {4,0}\right)$ ,半径 ${r}_{1} = 2$ ；
        
        圆 ${C}_{2} : {\left( x + 4\right) }^{2} + {y}^{2} = {16}$ 的圆心为 ${C}_{2}\left( {-4,0}\right)$ ,半径 ${r}_{2} = 4$ ,
        
        设动圆 $M$ 的半径为 $r$ ,由动圆 $M$ 与圆 ${C}_{1},{C}_{2}$ 都外切,得 $\left| {M{C}_{1}}\right|  = r + 2,\left| {M{C}_{2}}\right|  = r + 4$ ,
        
        则 $\left| {M{C}_{2}}\right|  - \left| {M{C}_{1}}\right|  = 2 < \left| {{C}_{1}{C}_{2}}\right|$ ,因此 $M$ 点的轨迹是以 ${C}_{1},{C}_{2}$ 为焦点的双曲线的右支,
        
        设方程为 $\dfrac{{x}^{2}}{{a}^{2}} - \dfrac{{y}^{2}}{{b}^{2}} = 1\left( {a > 0,b > 0}\right)$ ,则 $a = 1,c = 4,b = \sqrt{{c}^{2} - {a}^{2}} = \sqrt{15}$ ,
        
        所以 $M$ 的轨迹方程为 ${x}^{2} - \dfrac{{y}^{2}}{15} = 1\left( {x \geq  1}\right)$
    \end{jieda}

    \item 若 ${ab} \neq  0$ 且 $a \neq  b$ ,则 ${ax} - y + b = 0$ 和 $b{x}^{2} + a{y}^{2} = {ab}$ 所表示的曲线只可能是下图中的\dotfill{(\kh{D})}
    \xx{\includegraphics[width=0.35\textwidth]{bodyxb/src/bo_d3esa6c601uc738iqt7g_1_351_1714_208_253_0.jpg}}{\includegraphics[width=0.35\textwidth]{bodyxb/src/bo_d3esa6c601uc738iqt7g_1_558_1709_220_258_0.jpg}}{\includegraphics[width=0.35\textwidth]{bodyxb/src/bo_d3esa6c601uc738iqt7g_1_775_1706_253_259_0.jpg}}{\includegraphics[width=0.35\textwidth]{bodyxb/src/bo_d3esa6c601uc738iqt7g_1_1029_1709_264_258_0.jpg}}
    \item 已知双曲线 $\dfrac{{x}^{2}}{{a}^{2}} - \dfrac{{y}^{2}}{{b}^{2}} = 1\left( {a > 0,b > 0}\right)$ 的两焦点为 ${F}_{1}\text{ 、 }{F}_{2}$ . 若弦 ${AB}$ 过点 ${F}_{1}$ ,且 $A\text{ 、 }B$ 均在此双曲线的左支上, $\left| {AB}\right|  = l$ ,则 $\bigtriangleup  {AB}{F}_{2}$ 的周长为\dotfill{(\kh{B})}
    \xx{${4a} + l$}{${4a} + {2l}$}{${2a} + l$}{${4a} - l$} 
    
    \item 已知点 $P$ 在焦点为 ${F}_{1}$ 、 ${F}_{2}$ 的双曲线 $\dfrac{{x}^{2}}{{a}^{2}} - \dfrac{{y}^{2}}{{b}^{2}} = 1\left( {a > 0,b > 0}\right)$ 的右支上运动, 则$\triangle P{F}_{1}{F}_{2}$ 的内切圆的圆心一定在\dotfill{(\kh{A})}
    \xx{一条直线上}{一个圆上}{一个椭圆上}{一条抛物线上}
    \begin{jieda}
        设内心为$I$，过$I$做$IH \bot x$轴，根据几何关系可知$|F_1H| - |F_2H| = |PF_1| - |PF_2| = 2a$，故$H(a, 0)$，因此$I$在直线$x = a$上.
    \end{jieda}
    
    \item 若平面上到两定点距离的差的绝对值为常数的点的轨迹存在, 则轨迹可以是\dotfill{(\kh{D})}
    \xx{双曲线}{双曲线,直线}{双曲线,两条射线}{双曲线,直线,两条射线}

    \item 已知点 $M\left( {{x}_{0},{y}_{0}}\right)$ 是双曲线 $\dfrac{{x}^{2}}{{a}^{2}} - \dfrac{{y}^{2}}{{b}^{2}} = 1$ 上的任意一点,直线 $l$ 的方程是 $x = \dfrac{{a}^{2}}{c}$ , 其中 $c = \sqrt{{a}^{2} + {b}^{2}}$ . 求点 $M$ 到焦点 $F\left( {c,0}\right)$ 的距离与它到直线 $l$ 的距离之比.
    \begin{jieda}
        $\dfrac{\sqrt{(x_0-c)^2+y_0^2}}{\left|x-\dfrac{a^2}{c}\right|} = \dfrac{\sqrt{x_0^2-2cx_0+c^2 + \dfrac{b^2}{a^2}x_0^2-b^2}}{\left|x-\dfrac{a^2}{c}\right|} = \dfrac{\sqrt{\dfrac{c^2}{a^2}x_0^2-2cx_0+a^2}}{\left|x-\dfrac{a^2}{c}\right|} = \\
        \dfrac{c}{a} \cdot \dfrac{\sqrt{x_0^2 - 2\dfrac{a^2}{c}x_0 + \dfrac{a^4}{c^2}}}{\left|x-\dfrac{a^2}{c}\right|} = \dfrac{c}{a}$
    \end{jieda}

    \item 已知直线 $y = {ax} + 1$ 与双曲线 $3{x}^{2} - {y}^{2} = 1$ 相交于 $A$ 、 $B$ 两点. 当 $a$ 为何值时,以 ${AB}$ 为直径的圆过坐标原点? 
    \begin{jieda}
        设$A(x_1, y_1)$，$B(x_2, y_2)$，以 ${AB}$ 为直径的圆过坐标原点即$x_1x_2+y_1y_2 = 0$，代入直线方程$y = {ax} + 1$可得$x_1x_2+({ax_1} + 1)({ax_2} + 1) = 0$，即$(1+a^2)x_1x_2 + a(x_1+x_2) + 1 = 0$

        联立直线和双曲线方程有$\left\{\begin{array}{l}
             y = {ax} + 1  \\
             3{x}^{2} - {y}^{2} = 1 
        \end{array}\right.$，消去$y$可得$3{x}^{2} - ({ax} + 1)^2 = 1 $，即$(3-a^2)x^2-2ax -2 = 0$，$\Delta = 4a^2+8(3-a^2) > 0$，即$a^2 < 6$，根据韦达定理有$\left\{\begin{array}{l}
             x_1+x_2 = \dfrac{2a}{3-a^2} \\
             x_1x_2 = \dfrac{-2}{3-a^2}
        \end{array}\right.$
        
        代入$(1+a^2)x_1x_2 + a(x_1+x_2) + 1 = 0$有$(1+a^2)\left(\dfrac{-2}{3-a^2}\right) + a\left(\dfrac{2a}{3-a^2}\right) + 1 = 0$，整理得$-2-2a^2+2a^2+3-a^2 = 0$，故$a = \pm 1$
    \end{jieda}
    \item 已知 $M\left( {-2,0}\right) \text{ 、 }N\left( {2,0}\right)$ 两点,动点 $P$ 满足条件 $\left| {PM}\right|  - \left| {PN}\right|  = 2\sqrt{2}$ ,设动点 $P$ 的轨迹为 $W$ . 
    \begin{enumerate}
        \item 求 $W$ 的方程
        \item 若 $A$ 、 $B$ 是 $W$ 上的不同两点, $O$ 是坐标原点,求 $\vv{OA} \cdot  \vv{OB}$ 的最小值
    \end{enumerate}
    \begin{jieda}
        \begin{enumerate}
            \item $\because \left| {PM}\right|  - \left| {PN}\right|  = 2\sqrt{2} < 4 = \left| {MN}\right|$
            
            由双曲线定义可知: $P$ 的轨迹是以 $M$ , $N$ 为焦点的双曲线的右支
            
            $c = 2,a = \sqrt{2},{b}^{2} = {c}^{2} - {a}^{2} = 2$
            
            因此$W$ 的方程为: $\dfrac{{x}^{2}}{2} - \dfrac{{y}^{2}}{2} = 1\left( {x > 0}\right)$
            
            \item 
            \begin{dingautolist}{172}
                \item 当直线 ${AB}$ 斜率不存在时,设直线 ${AB}$ 方程为： $x = m\left( {m > \sqrt{2}}\right)$
                此时 $A\left( {m,\sqrt{{m}^{2} - 2}}\right) ,B\left( {m, - \sqrt{{m}^{2} - 2}}\right) \therefore \vv{OA} \cdot  \vv{OB} = {m}^{2} - \left( {{m}^{2} - 2}\right)  = 2$
                \item 
                当直线 ${AB}$ 斜率存在时,设直线 ${AB}$ 方程为： $y = {kx} + b$，设$A(x_1, y_1)$，$B(x_2, y_2)$，

                $\vv{OA} \cdot  \vv{OB} = {x}_{1}{x}_{2} + {y}_{1}{y}_{2} = {x}_{1}{x}_{2} + \left( {k{x}_{1} + b}\right) \left( {k{x}_{2} + b}\right)  = \left( {{k}^{2} + 1}\right) {x}_{1}{x}_{2} + {kb}\left( {{x}_{1} + {x}_{2}}\right)  + {b}^{2}$
                
                联立直线和双曲线方程可得： $\left( {1 - {k}^{2}}\right) {x^2} - {2kbx} - \left( {{b}^{2} + 2}\right)  = 0$
                
                因为$W$为双曲线的右一支，故上式有两个不等的正实数根
                
                $\therefore \left\{  \begin{array}{l} \Delta  = 4{k}^{2}{b}^{2} + 4\left( {1 - {k}^{2}})(b^2+2) > 0\right. \\ 
                {x}_{1} + {x}_{2} = \dfrac{2kb}{1 - {k}^{2}} > 0 \\  
                {x}_{1}{x}_{2} = \dfrac{-{b}^{2} - 2}{1 - {k}^{2}} > 0 \end{array}\right.$
                解得: ${k}^{2} > 1$
                
                
                $\vv{OA} \cdot  \vv{OB} = \dfrac{\left( {{b}^{2} + 2}\right) \left( {{k}^{2} + 1}\right) }{{k}^{2} - 1} + \dfrac{2{k}^{2}{b}^{2}}{1 - {k}^{2}} + {b}^{2} = \dfrac{2{k}^{2} + 2}{{k}^{2} - 1} = 2 + \dfrac{4}{{k}^{2} - 1}$
                
                由 ${k}^{2} > 1$ 得： $\therefore \vv{OA} \cdot  \vv{OB} = 2 + \dfrac{4}{{k}^{2} - 1} > 2$
            \end{dingautolist}
            综上所述, $\vv{OA} \cdot  \vv{OB}$ 的最小值为 2
        \end{enumerate}
    \end{jieda}
\end{enumerate}

\clearpage
\subsection*{双曲线的性质 (2)}
\begin{enumerate}
    \item 双曲线 $C : {y}^{2} - \dfrac{{x}^{2}}{4} = 1$ 的渐近线方程为\blkda[3]{$y =  \pm  \dfrac{1}{2}x$},设双曲线${C}_{1} : \dfrac{{x}^{2}}{{a}^{2}} - \dfrac{{y}^{2}}{{b}^{2}} = 1\left( {a > 0,b > 0}\right)$ 经过点$(4,1)$,且与双曲线 $C$ 具有相同渐近线,则双曲线 ${C}_{1}$ 的标准方程为\blkda[3]{$\dfrac{{x}^{2}}{12} - \dfrac{{y}^{2}}{3} = 1$}.
    \begin{jieda}
        \begin{enumerate}
            \item 双曲线 $C : {y}^{2} - \dfrac{{x}^{2}}{4} = 1$ 的焦点在 $y$ 轴上,且 $a = 1,b = 2$ ,故渐近线方程为 $y =  \pm  \dfrac{1}{2}x$ 故答案为: $y =  \pm  \dfrac{1}{2}x$
            \item 由双曲线 ${C}_{1}$ 与双曲线 $C$ 具有相同渐近线,可设 ${C}_{1} : {y}^{2} - \dfrac{{x}^{2}}{4} = \lambda \left( {\lambda  \neq  0}\right)$ ,

            代入$(4,1)$有 ${1}^{2} - \dfrac{{4}^{2}}{4} = \lambda$ ,所以 $\lambda  =  - 3$ ,故 ${C}_{1} : {y}^{2} - \dfrac{{x}^{2}}{4} =  - 3$ ,化简得 $\dfrac{{x}^{2}}{12} - \dfrac{{y}^{2}}{3} = 1$ .
        \end{enumerate}
    \end{jieda}

    \item 已知双曲线 $\dfrac{{x}^{2}}{{a}^{2}} - \dfrac{{y}^{2}}{{b}^{2}} = 1\left( {a > 0,b > 0}\right)$ 的离心率为 $\sqrt{2}$ ,其一条渐近线被圆 ${\left( x - m\right) }^{2} + {y}^{2} = 4\left( {m > 0}\right)$ 截得的线段长为 $2\sqrt{2}$ ,则实数 $m$ 的值为\blkda{2}.
    \begin{jieda}
        由题可得: $c = \sqrt{2}a$ ,即有 $a = b$ ,渐近线方程为 $y =  \pm  x$ ,圆 ${\left( x - m\right) }^{2} + {y}^{2} = 4\left( {m > 0}\right)$ 的圆心为$(m,0)$,半径为 2,可得圆心到直线的距离为 ${d} = \dfrac{\left| m\right| }{\sqrt{2}}$ ,则直线被圆截得的弦长为 $2\sqrt{4 - \dfrac{{m}^{2}}{2}} = 2\sqrt{2}$ ,解得 ${m} = 2$ ($-2$ 舍去).
    \end{jieda}

    \item 已知 ${F}_{1}\left( {-3,0}\right) \text{ 、 }{F}_{2}\left( {3,0}\right)$ 两点. 若满足条件 $\left| {M{F}_{1}}\right|  + \left| {M{F}_{2}}\right|  = {2m} + 1$ 的动点 $M$ 的轨迹是椭圆,满足条件 $\left|\left|{N{F}_{1}}\right| - \left| {N{F}_{2}}\right|\right|  = {2m} - 1$ 的动点 $N$ 的轨迹是双曲线,则实数 $m$ 的取值范围是\blkda{$\left(\dfrac{5}{2}, \dfrac{7}{2}\right)$}.
    \begin{jieda}
        因为$M$ 的轨迹是椭圆，所以$2m+1 > 6$，故$m > \dfrac{5}{2}$

        因为$N$ 的轨迹是双曲线，所以$0 < 2m-1 < 6$，故$\dfrac{1}{2} < m < \dfrac{7}{2}$

        综上，$m \in \left(\dfrac{5}{2}, \dfrac{7}{2}\right)$
    \end{jieda}

    \item 若双曲线 $\dfrac{{x}^{2}}{4} - \dfrac{{y}^{2}}{m} = 1$ 的渐近线方程为 $y =  \pm  \dfrac{\sqrt{3}}{2}x$ ,则此双曲线的焦点坐标是\blkda[4]{$\left(\sqrt{7}, 0\right)$和$\left(-\sqrt{7}, 0\right)$}.
    \begin{jieda}
         $\dfrac{{x}^{2}}{4} - \dfrac{{y}^{2}}{m} = 1$ 的渐近线方程为$y = \pm \dfrac{\sqrt{m}}{2}x$，故$m = 3$.

         因此$a = 2, b = \sqrt{3}, c = \sqrt{7}$，故焦点坐标为$\left(\sqrt{7}, 0\right)$和$\left(-\sqrt{7}, 0\right)$
    \end{jieda}

    \item 设 $P$ 是双曲线 $\dfrac{{x}^{2}}{{a}^{2}} - \dfrac{{y}^{2}}{9} = 1$ 上一点,双曲线的一条渐近线方程为 ${3x} - {2y} = 0$ ,${F}_{1}$ 、 ${F}_{2}$ 分别是双曲线的左、右焦点,若 $\left| {P{F}_{1}}\right|  = 3$ ,则 $\left| {P{F}_{2}}\right|  =$ \blkda{7}.
    \begin{jieda}
        双曲线 $\dfrac{{x}^{2}}{{a}^{2}} - \dfrac{{y}^{2}}{9} = 1$的渐近线为$y = \pm \dfrac{3}{a}x$，故根据题意可知$a^2 = 4$，故根据双曲线定义可知$\left| {P{F}_{2}}\right| = 3+4 = 7$
    \end{jieda}

    \item 过点 $P\left( {4,3}\right)$ ,且渐近线方程为 $x \pm  {2y} = 0$ 的双曲线的标准方程是\blkda[3]{$\dfrac{{y}^{2}}{5} - \dfrac{{x}^{2}}{20} = 1$}.

    \item 已知双曲线 $\dfrac{{x}^{2}}{12} - \dfrac{{y}^{2}}{4} = 1$ 的右焦点为 $F$ . 若过点 $F$ 的直线 $l$ 与双曲线的右支有且只有一个交点,则直线 $l$ 的斜率的取值范围是\blkda{$\left[-\dfrac{\sqrt{3}}{3}, \dfrac{\sqrt{3}}{3}\right]$}.

    \item 双曲线 $\dfrac{{x}^{2}}{{a}^{2}} - \dfrac{{y}^{2}}{{b}^{2}} = 1\left( {a > b > 0}\right)$ 的实轴长 4 ,则双曲线上的一点 $\left( {4,\sqrt{3}}\right)$ 到两渐近线的距离的乘积等于\blkda{$\dfrac{4}{5}$}.
    \begin{jieda}
        设双曲线上一点为$(x_0, y_0)$，则其到双曲线的两条渐近线$y = \pm \dfrac{b}{a}x$的距离乘积为$\dfrac{\left|y_0 + \dfrac{b}{a}x_0\right|}{\sqrt{1+\dfrac{b^2}{a^2}}} \times \dfrac{\left|y_0 - \dfrac{b}{a}x_0\right|}{\sqrt{1+\dfrac{b^2}{a^2}}} = \dfrac{\left|y_0^2 - \dfrac{b^2}{a^2}x_0^2\right|}{\dfrac{c^2}{a^2}} = \dfrac{a^2b^2}{c^2}$

        实轴长 4$\Longleftrightarrow a = 2$，将$\left( {4,\sqrt{3}}\right)$代入方程，可得$b = 1$，因此$c = \sqrt{5}$，故目标式为$\dfrac{4}{5}$
    \end{jieda}

    \item 已知双曲线 $C : \dfrac{{x}^{2}}{{a}^{2}} - \dfrac{{y}^{2}}{{b}^{2}} = 1\left( {a > 0,b > 0}\right)$ 的左、右顶点分别为 ${A}_{1},{A}_{2},F$ 为 $C$ 的右焦点,$C$ 的离心率为 2,若 $P$ 为 $C$ 右支上一点, ${PF} \bot  F{A}_{2}$ ,记 $\angle {A}_{1}P{A}_{2} = \theta \left( {0 < \theta  < \dfrac{\pi }{2}}\right)$ ,则 $\tan \theta  =$ \blkda{$\dfrac{1}{2}$}.
    \begin{jieda}
        设 $C$ 的焦距为 ${2c}$ ,点 $P\left( {{x}_{0},{y}_{0}}\right)$ ,由 $C$ 的离心率为 2 可知 $c = {2a},b = \sqrt{3}a$ ,因为 ${PF}\bot {F{A}_{2}}$ ,所以 ${x}_{0} = c$ ,将 $P\left( {c,{y}_{0}}\right)$ 代入 $C$ 的方程得 $\dfrac{{c}^{2}}{{a}^{2}} - \dfrac{{y}_{0}^{2}}{{b}^{2}} = 1$ ,即 $\left| {y}_{0}\right|  = \sqrt{3}b$，所以 $\tan \angle P{A}_{2}F = \dfrac{\sqrt{3}b}{c - a} = 3,\tan \angle P{A}_{1}F = \dfrac{\sqrt{3}b}{c - \left( {-a}\right) } = 1$
        
        故 $\tan \theta  = \tan \left( {\angle P{A}_{2}F - \angle P{A}_{1}F}\right)  = \dfrac{3 - 1}{1 + 3 \times  1} = \dfrac{1}{2}$ .
    \end{jieda} 

    \item 已知双曲线 $c : \dfrac{{x}^{2}}{2} - \dfrac{{y}^{2}}{{b}^{2}} = 1\left( {b > 0}\right)$ 的离心率为 $e$ ,若 $e \in  \left( {\sqrt{5},\sqrt{10}}\right)$ 则 $C$ 的焦点到一条渐近线的距离的取值范围为\blkda{$\left( {2\sqrt{2},3\sqrt{2}}\right)$}.
    \begin{jieda}
        因为 $e = \sqrt{1 + \dfrac{{b}^{2}}{{a}^{2}}} = \sqrt{1 + \dfrac{{b}^{2}}{2}} \in  \left( {\sqrt{5},\sqrt{10}}\right)$ ,所以 $b \in  \left( {2\sqrt{2},3\sqrt{2}}\right)$，而 $C$ 的焦点 $\left( {\pm c,0}\right)$ 到渐近线 ${bx} \pm  {ay} = 0$ 的距离为 $d = \dfrac{bc}{\sqrt{{a}^{2} + {b}^{2}}} = b$ .
        
        所以距离的取值范围为 $\left( {2\sqrt{2},3\sqrt{2}}\right)$ .
    \end{jieda}

    \item 已知双曲线 $\dfrac{{y}^{2}}{{a}^{2}} - \dfrac{{x}^{2}}{{b}^{2}} = 1\left( {a > 0,b > 0}\right)$ 的两个焦点分别为 ${F}_{1}\left( {0, - c}\right) ,{F}_{2}\left( {0,c}\right)$ , $c > 0$ ,以 ${F}_{1}{F}_{2}$ 为直径的圆与双曲线在第一象限的交点为 $P$ ,若直线 $P{F}_{1}$ 与圆 $E$ : ${x}^{2} + {\left( y - \dfrac{c}{3}\right) }^{2} = \dfrac{{b}^{2}}{9}$ 相切,则双曲线的离心率是\blkda{$\dfrac{\sqrt{65}}{7}$}.
    \begin{jieda}
        设直线 $P{F}_{1}$ 与圆 $E : {x}^{2} + {\left( y - \dfrac{c}{3}\right) }^{2} = \dfrac{{b}^{2}}{9}$ 相切于点 $M$ ,
        
        \begin{center}
        \includegraphics[width=0.3\textwidth]{bodyxb/src/bo_d3erhjbef24c73cu4rp0_2_241_201_364_336_0.jpg}
        \end{center}
        
        因为以 ${F}_{1}{F}_{2}$ 为直径的圆过点 $P$ ,所以 $P{F}_{1} \bot  P{F}_{2}$ ,
        
        又圆 $E$ 与直线 $P{F}_{1}$ 的切点为 $M$ ,所以 ${EM} \bot  P{F}_{1}$ ,从而 ${EM}//P{F}_{2}$ .
        
        由 $\bigtriangleup {F}_{1}{EM} \backsim  \bigtriangleup {F}_{1}{F}_{2}P$ ,得 $\dfrac{P{F}_{2}}{EM} = \dfrac{{F}_{1}{F}_{2}}{E{F}_{1}}$ ,所以 $P{F}_{2} = \dfrac{{EM} \cdot  {F}_{1}{F}_{2}}{E{F}_{1}} = \dfrac{\dfrac{1}{3}b \cdot  {2c}}{\dfrac{4c}{3}} = \dfrac{1}{2}b$ .
        
        又 $P{F}_{1} = P{F}_{2} + {2a} = \dfrac{1}{2}b + {2a}$ ,所以 ${\left( \dfrac{1}{2}b\right) }^{2} + {\left( \dfrac{1}{2}b + 2a\right) }^{2} = 4{c}^{2}$ ,解得 ${7b} = {4a}$ ,
        
        因此离心率为 $e = \dfrac{c}{a} = \sqrt{1 + {\left( \dfrac{b}{a}\right) }^{2}} = \dfrac{\sqrt{65}}{7}$
    \end{jieda}

    \item 如果椭圆 $\dfrac{{x}^{2}}{3{m}^{2}} + \dfrac{{y}^{2}}{5{n}^{2}} = 1$ 和双曲线 $\dfrac{{x}^{2}}{2{m}^{2}} - \dfrac{{y}^{2}}{3{n}^{2}} = 1$ 有公共焦点,那么双曲线的渐近线方程是\dotfill{(\kh{B})}
    
    \xx{$y =  \pm  \dfrac{\sqrt{2}}{2}x$}{$y =  \pm  \dfrac{\sqrt{3}}{4}x$}{$y =  \pm  \sqrt{2x}$}{$y =  \pm  \dfrac{2\sqrt{3}}{3}x$}

    \item 若双曲线 $\dfrac{{x}^{2}}{{a}^{2}} - \dfrac{{y}^{2}}{2} = 1\left( {a > 2}\right)$ 的两条渐近线的夹角为 $\dfrac{\pi }{3}$ ,则此双曲线的焦距与实轴长之比 $\dfrac{c}{a}$ 等于\dotfill{(\kh{D})}
    \xx{$\sqrt{2}$}{$\sqrt{3}$}{$\dfrac{2\sqrt{6}}{3}$}{$\dfrac{2\sqrt{3}}{3}$}
    \begin{jieda}
        渐近线$y = \pm \dfrac{\sqrt{2}}{a}x$，故根据两条渐近线的夹角为 $\dfrac{\pi }{3}$可知$\dfrac{\sqrt{2}}{a} = \dfrac{\sqrt{3}}{3}$，故$a = \sqrt{6}$，故$c = 2\sqrt{2}$，因此$\dfrac{c}{a} = \dfrac{2\sqrt{3}}{3}$
    \end{jieda}
    
    \item 过双曲线 $2{x}^{2} - {y}^{2} - 2 = 0$ 的右焦点作直线 $l$ 交双曲线于 $A\text{ 、 }B$ 两点,若 $\left| {AB}\right|  = 4$ , 则这样的直线有\dotfill{(\kh{B})}
    \xx{4 条}{3 条}{2 条}{1 条}
    \begin{jieda}
        由 $2{x}^{2} - {y}^{2} = 2$ ,得 ${x}^{2} - \dfrac{{y}^{2}}{2} = 1$ .
        
        当 $l$ 斜率不存在时, $\left| {AB}\right|  = \dfrac{2{b}^{2}}{a} = 4$ 符合要求.
        
        当 $l$ 斜率存在时,若 $A,B$ 两点都在右支上,则 $\left| {AB}\right|  > 4$ 不符合要求.
        
        若 $A,B$ 在左、右两支上,满足题意的直线有两条. 所以共3条.
    \end{jieda}
    
    \item 某中心接到其正东、正西、正北方向三个观测点的报告；正西、正北两个观测点同时听到了一声巨响,正东观测点听到的时间比其它两观测点晚 2s,已知各观测点到该中心的距离是 ${680}{m}$ ,则该巨响发生在接报中心的(   )处(假定当时声音传播的速度为 ${340}{m}/{s}$ ,相关各点均在同一平面上)\dotfill{(\kh{A})}
    \xx{西偏北 ${45}^{ \circ  }$ 方向,距离 ${340}\sqrt{3}{m}$}{东偏南 ${45}^{ \circ  }$ 方向,距离 ${340}\sqrt{3}{m}$}{西偏北 45° 方向,距离 ${170}\sqrt{3}{m}$}{东偏南 ${45}^{ \circ  }$ 方向,距离 ${170}\sqrt{3}{m}$}
    \begin{jieda}
        如图,
        
        \begin{center}
        \includegraphics[width=0.2\textwidth]{bodyxb/src/bo_d3erhjbef24c73cu4rp0_2_244_1796_349_326_0.jpg}
        \end{center}
        
        以接报中心为原点 $O$ ,正东、正北方向为 $x$ 轴、 $y$ 轴正向,建立直角坐标系. 设 $A\text{ 、 }B\text{ 、 }C$ 分别是西、东、北观测点,则 $A\left( {-{680},0}\right)$ , $B\left( {{680},0}\right)$ , $C\left( {0,{680}}\right)$ .
        
        设 $P\left( {x,y}\right)$ 为巨响为生点,由 $A\text{ 、 }C$ 同时听到巨响声,得 $\left| {PA}\right|  = \left| {PC}\right|$ ,故 $P$ 在 ${AC}$ 的垂直平分线 ${PO}$ 上, ${PO}$ 的方程为 $y =  - x$ ,因 $B$ 点比 ${A}$ 点晚 $2{s}$ 听到爆炸声,故, $\left| {PB}\right|  - \left| {PA}\right|  = {340} \times  2 = {680}$ 由双曲线定义知 $P$ 点在以 $A,B$ 为焦点的双曲线左支 $\dfrac{{x}^{2}}{{a}^{2}} - \dfrac{{y}^{2}}{{b}^{2}} = 1\left( {x < 0}\right)$ 上
        
        依题意得 $a = {340},c = {680},\therefore {b}^{2} = {c}^{2} - {a}^{2} = {680}^{2} - {340}^{2} = 3 \times  {340}^{2}$ ，故双曲线方程为 $\dfrac{{x}^{2}}{{340}^{2}} - \dfrac{{y}^{2}}{3 \times  {340}^{2}} = 1$，将 $y =  - x$ 代入上式,得$x =  \pm  {170}\sqrt{6},\because x < 0,\therefore x =  - {170}\sqrt{6},y = {170}\sqrt{6}$ ,即 $P\left( {-{170}\sqrt{6},{170}\sqrt{6}}\right)$，故 ${PO} = {340}\sqrt{3}$ .
        
        故巨响发生在接报中心的西偏北 ${45}^{ \circ  }$ 距中心 ${340}\sqrt{3}{m}$ 处.
    \end{jieda} 

    \item 已知 ${F}_{1}\text{ 、 }{F}_{2}$ 为双曲线 $\dfrac{{x}^{2}}{{a}^{2}} - \dfrac{{y}^{2}}{{b}^{2}} = 1\left( {a > 0,b > 0}\right)$ 的左、右焦点,过 ${F}_{2}$ 作垂直于 $x$ 轴的直线,交此双曲线于点 $P$ ,且 $\angle P{F}_{1}{F}_{2} = {30}^{ \circ  }$ ,求此双曲线的渐近线方程.
    \begin{jieda}
        $S_{\triangle PF_1F_2} = b^2cot\ang{30} = \sqrt{3}b^2 = c \times \dfrac{b^2}{a}$，故$e = \sqrt{3}$，即$\dfrac{b}{a} = \sqrt{2}$，即渐近线方程为$y = \pm \sqrt{2}x$
    \end{jieda}

    \item 已知直线 $l$ 交双曲线 $\dfrac{{x}^{2}}{{a}^{2}} - \dfrac{{y}^{2}}{{b}^{2}} = 1$ 及其渐近线于 $A\text{ 、 }B\text{ 、 }C\text{ 、 }D$ 四点. 求证: 夹在双曲线及其渐近线间的线段相等.
    \begin{jieda}
        证明: 设直线 $l : y = {kx} + m$ ,与双曲线和渐近线从左到右依次交于 ${A},{B},{C},{D}$设横坐标依次为 ${x}_{1},{x}_{2},{x}_{3},{x}_{4}$ 
        
        由 $\left\{  {\begin{matrix} y = {kx} + m \\  {b}^{2}{x}^{2} - {a}^{2}{y}^{2} = {a}^{2}{b}^{2} \end{matrix}\text{ 得 }{x}_{2} + {x}_{3} = \dfrac{2{a}^{2}{km}}{{b}^{2} - {a}^{2}{k}^{2}}}\right.$
        同理得${x}_{1} + {x}_{4} = \dfrac{2{a}^{2}{km}}{{b}^{2} - {a}^{2}{k}^{2}}$
        所以 ${x}_{1} + {x}_{4} = {x}_{2} + {x}_{3}$
        即 ${{AD}}$ 和 ${{BC}}$ 的中点重合,即 $\left| {AB}\right|  = \left| {CD}\right|$
    \end{jieda}

    \item (选做) 过双曲线 $y = \dfrac{k}{x}$ (常数 $k > 0$ )上任意一点 $A$ 作 ${AE}//x$ 轴,交 $y$ 轴于点 $E$ ,作 ${AF}//y$ 轴,交 $x$ 轴于点 $F$ ,得到矩形 ${AEOF}$ ,则它的面积 $S = k,k$ 是与点 $A$ 位置无关的常数. 试把这个结论推广到一般双曲线 $\dfrac{{x}^{2}}{{a}^{2}} - \dfrac{{y}^{2}}{{b}^{2}} = 1\left( {a > 0,b > 0}\right)$ ,并证明你的推广.
    \begin{jieda}
        推广结论：设 $A$ 是双曲线 $\dfrac{{x}^{2}}{{a}^{2}} - \dfrac{{y}^{2}}{{b}^{2}} = 1\left( {a > 0,b > 0}\right)$ 上任意一点,过点 $A$ 分别作渐近线 ${bx} \pm  {ay} = 0$ 的平行线 ${AE}$ 、 ${AF}$ ,并分别交渐近线于 $E$ 、 $F$ ,得到平行四边形 ${AEOF}$ ,则平行四边形 ${AEOF}$ 的面积 $S$ 是与点 $A$ 位置无关的常数.
    
        证明: 设 $A\left( {{x}_{0},{y}_{0}}\right)$ ,直线 ${AE}$ 的方程为 $y - {y}_{0} =  - \dfrac{b}{a}\left( {x - {x}_{0}}\right)$ 
        
        联立方程组 $\left\{  \begin{array}{l} y - {y}_{0} =  - \dfrac{b}{a}\left( {x - {x}_{0}}\right) \\  y = \dfrac{b}{a}x \end{array}\right.$ ,解得交点 $E\left( {\dfrac{a{y}_{0} + b{x}_{0}}{2b},\dfrac{a{y}_{0} + b{x}_{0}}{2a}}\right)$ 
        
        同理可得$F\left( {\dfrac{a{y}_{0} - b{x}_{0}}{2b}, -\dfrac{a{y}_{0} - b{x}_{0}}{2a}}\right)$

        $S = \left|\dfrac{a{y}_{0} - b{x}_{0}}{2b} \times \dfrac{a{y}_{0} + b{x}_{0}}{2a} + \dfrac{a{y}_{0} - b{x}_{0}}{2a} \times \dfrac{a{y}_{0} + b{x}_{0}}{2b}\right| = \dfrac{2|a^2y_0^2-b^2x_0^2|}{4ab} = \dfrac{2a^2b^2}{4ab} = \dfrac{ab}{2}$
        
        所以 $S = \dfrac{1}{2}a{b}$ 是与点 $A$ 位置无关的常数.
    \end{jieda}

\end{enumerate}

\clearpage
\subsection*{双曲线的性质 (3)}
\begin{enumerate}
    \item 若方程 $\left( {k + 1}\right) {y}^{2} - {x}^{2} = k - 1$ 表示焦点在 $x$ 轴上的双曲线,则实数 $k$ 的取值范围是\blkda{$(-1,1)$}.

    \item 若点 $F\left( {0,3}\right)$ 是双曲线 ${8k}{x}^{2} - k{y}^{2} = 8$ 的一个焦点,则实数 $k =$ \blkda{$-1$}.

    \item 若双曲线 $2{x}^{2} - {y}^{2} = k$ 的焦距是 6 ,则离心率 ${e} =$ \blkda{$\sqrt{3}$ 或 $\dfrac{\sqrt{6}}{2}$}.

    \item 若方程 $\sqrt{{x}^{2} - 1} = x + m$ 无实数解,则实数 $m$ 的取值范围是\blkda[4]{$\left( {-\infty , - 1}\right)  \cup  \lbrack 0,1)$}.

    \item 若定义双曲线对称轴与双曲线交点即双曲线顶点,则等轴双曲线 ${xy} = 4$ 的焦距为\blkda{8}.
    \begin{jieda}
        $2a = 4\sqrt{2}, e = \sqrt{2}, 2c = 8$.
    \end{jieda}

    \item 从双曲线 $\dfrac{{x}^{2}}{{a}^{2}} - \dfrac{{y}^{2}}{{b}^{2}} = 1$ 上任意一点 ${M}$ 分别作实轴的平行线,与渐近线交于 ${P}$ 和 ${Q}$ 两点, 则 $\left| {MP}\right|  \cdot  \left| {MQ}\right|$ 的值是\blkda{${a}^{2}$}.
    \begin{jieda}
        双曲线$\dfrac{{x}^{2}}{{a}^{2}} - \dfrac{{y}^{2}}{{b}^{2}} = 1$的渐近线为$y = \pm \dfrac{b}{a}x$
    
        设$M(x_0, y_0)$，则$P\left(\dfrac{a}{b}y_0, y_0\right)$，$Q\left(-\dfrac{a}{b}y_0, y_0\right)$，$|MP| \cdot |MQ| = \left|x_0 - \dfrac{a}{b}y_0\right| \cdot \left|x_0 + \dfrac{a}{b}y_0\right| = \left|x_0^2 - \dfrac{a^2}{b^2}y_0^2\right| = a^2$
    \end{jieda}

    \item 设曲线 ${C}_{1} : \dfrac{{x}^{2}}{{a}^{2}} - \dfrac{{y}^{2}}{{b}^{2}} = 1$ ,曲线 ${C}_{2} : \dfrac{{y}^{2}}{{b}^{2}} - \dfrac{{x}^{2}}{{a}^{2}} = 1$ ,曲线 ${C}_{3} : \dfrac{{x}^{2}}{{b}^{2}} - \dfrac{{y}^{2}}{{a}^{2}} = 1,{a}^{2} \neq  {b}^{2}$ ,则\dotfill{(\kh{B})}
    \xx{${C}_{1}$ 和 ${C}_{2}$ 有公共焦点}{${C}_{1}$ 和 ${C}_{3}$ 有公共焦点}{${C}_{3}$ 和 ${C}_{2}$ 有公共渐近线}{${C}_{1}$ 和 ${C}_{3}$ 有公共渐近线}
    
    \item 已知双曲线 $\dfrac{{x}^{2}}{{a}^{2}} - \dfrac{{y}^{2}}{{b}^{2}} = 1\left( {a > 0,b > 0}\right)$ 的左、右焦点分别是 ${F}_{1}\text{ 、 }{F}_{2}$ ,它的焦距为 ${2c}$ . 若点 $P$ 为此双曲线的右支上一点,则 $\bigtriangleup P{F}_{1}{F}_{2}$ 的内切圆的圆心横坐标为\dotfill{(\kh{A})}
    \xx{$a$}{$b$}{$c$}{$a + b - c$}
    \begin{jieda}
        设内切圆与$F_1F_2$切于$M$，则$2a = |PF_1| - |PF_2| = |MF_1| - |MF_2|$，故$M(a, 0)$，因此内切圆的圆心横坐标为$a$.
    \end{jieda}
    
    \item 已知双曲线 $\dfrac{{x}^{2}}{{a}^{2}} - \dfrac{{y}^{2}}{{b}^{2}} = 1\left( {a > 0,b > 0}\right)$ 的右焦点为 $F$ ,直线 $x = \dfrac{{a}^{2}}{c}$ ( $c$ 为半焦距)与它的一条渐近线交于点 $A$ . 若 $\bigtriangleup {OAF}$ 的面积为 $\dfrac{{a}^{2}}{2}$ ($O$ 为坐标原点),则此双曲线的两条渐近线的夹角为\dotfill{(\kh{D})}
    \xx{${30}^{ \circ  }$}{${45}^{ \circ  }$}{${60}^{ \circ  }$}{${90}^{ \circ  }$}
    \begin{jieda}
        不妨设$A$在第一象限，则$A\left(\dfrac{a^2}{c}, \dfrac{ab}{c}\right)$，故根据$\bigtriangleup {OAF}$ 的面积为 $\dfrac{{a}^{2}}{2}$可知$\dfrac{a^2}{2} = \dfrac{1}{2} \times c \times \dfrac{ab}{c}$，故$a = b$，因此两渐近线垂直.
    \end{jieda}
    
    \item 若 $P$ 为双曲线 $\dfrac{{x}^{2}}{9} - \dfrac{{y}^{2}}{16} = 1$ 的右支上一点, $M\text{ 、 }N$ 分别是圆 ${\left( x + 5\right) }^{2} + {y}^{2} = 4$ 和 ${\left( x - 5\right) }^{2} + {y}^{2} = 1$ 上的点,则 $\left| {PM}\right|  - \left| {PN}\right|$ 的最大值为\dotfill{(\kh{C})}
    \xx{7}{8}{9}{10}
    \begin{jieda}
        $\left| {PM}\right|  - \left| {PN}\right| \leq (2+|PF_1|) - (|PF_2| - 1) = 3+6 = 9$
    \end{jieda}
    
    \item 已知过定点 $P\left( {0,2}\right)$ 的直线 $l$ 斜率为$k$，$l$与双曲线 $C : {x}^{2} - \dfrac{{y}^{2}}{4} = 1$ 相交于 $A,B$ 两点,问:
    \begin{enumerate}
        \item $k$ 为何值时, $A,B$ 分别在双曲线的两支上？
        \item $k$ 为何值时, $A,B$ 都在双曲线的左支上?
    \end{enumerate}
    \begin{jieda}
        双曲线 $C : {x}^{2} - \dfrac{{y}^{2}}{4} = 1$的渐近线为$y = \pm 2x$
        \begin{enumerate}
            \item 数形结合可知：$k \in (-2,2)$时满足要求.
            \item 数形结合可知：$A,B$ 都在双曲线的左支的必要条件是$k > 2$，直线$l$的方程为$y = kx + 2$，与双曲线方程联立，消去$y$可得$4x^2 - (kx+2)^2 = 4$，即$(4-k^2)x^2-4kx - 8 = 0$，$\Delta > 0 \Longleftrightarrow k^2 < 8$
            
            综上，$k \in \left( {2,2\sqrt{2}}\right)$时满足要求.
        \end{enumerate}
    \end{jieda}
    
    \item 双曲线 $\dfrac{{y}^{2}}{12} - \dfrac{{x}^{2}}{13} = 1$ 的一支上的三点 $A\left( {{x}_{1},{y}_{1}}\right)$ 、 ${B}\left( {\sqrt{26},6}\right)$ 、 ${C}\left( {{x}_{2},{y}_{2}}\right)$ 与焦点 $F\left( {0,5}\right)$ 的距离成等差数列. 
    \begin{enumerate}
        \item 求 ${y}_{1} + {y}_{2}$
        \item 证明:线段 ${AC}$ 的垂直平分线经过一定点
    \end{enumerate}
    \begin{jieda}
        \begin{enumerate}
            \item 根据双曲线第二定义可知$|AF| = e\left|y_1-\dfrac{a^2}{c}\right| = \dfrac{5}{2\sqrt{3}} \times \left(y_1-\dfrac{12}{5}\right)$，$|CF| = e\left|y_2-\dfrac{a^2}{c}\right| = \dfrac{5}{2\sqrt{3}} \times \left(y_2-\dfrac{12}{5}\right)$

            根据题意可知$|AF|+|CF| = 2|BF| = 6\sqrt{3}$，即$\dfrac{5}{2\sqrt{3}} \times \left(y_1+y_2-\dfrac{24}{5}\right) = 6\sqrt{3}$，故$y_1+y_2 = 12$
            
            \item $\left\{\begin{array}{l}
                 \dfrac{{y_1}^{2}}{12} - \dfrac{{x_1}^{2}}{13} = 1  \\
                 \dfrac{{y_2}^{2}}{12} - \dfrac{{x_2}^{2}}{13} = 1
            \end{array}\right.$，两式作差得到$\dfrac{(y_1+y_2)(y_1-y_2)}{12} - \dfrac{(x_1+x_2)(x_1-x_2)}{13} = 0$，即$\dfrac{y_1-y_2}{x_1-x_2} = \dfrac{12}{13} \cdot \dfrac{x_1+x_2}{y_1+y_2}$
            设$AC$中点为$M(x_0, y_0)$，则$k_{AC} = \dfrac{12}{13} \cdot \dfrac{x_0}{y_0}$

            $AC$垂直平分线方程为$y = -\dfrac{1}{k_{AC}} (x - x_0) + y_0 = -\dfrac{13y_0}{12x_0}(x-x_0)+y_0$，根据上一问可知$y_0 = 6$，故${AC}$ 的垂直平分线经过定点$\left( {0,\dfrac{25}{2}}\right)$
        \end{enumerate}
    \end{jieda}
    
    \item 已知 $M\left( {-3,0}\right) \text{ 、 }N\left( {3,0}\right) \text{ 、 }B\left( {1,0}\right)$ 三点,圆 $C$ 与直线 ${MN}$ 相切于点 $B$ ,分别过点 $M\text{ 、 }N$ ,且与圆 $C$ 相切的两条直线相交于点 $P$ ,求点 $P$ 的轨迹方程.
    \begin{jieda}
        设$PM, PN$切圆$C$于$M'N'$，根据题意可知$|PM'| = |PN'|$，$MM' = MB$，$NN' = NB$，故$|PM| - |PN| = |MB| - |NB| = 2$，故$P$点轨迹符合双曲线定义，为双曲线的右支.

        根据题意可知$a = 1, c = 3$，故$b = 2\sqrt{2}$，故轨迹方程为${x}^{2} - \dfrac{{y}^{2}}{8} = 1\left( x > 0 \text{且} {y \neq  0}\right)$
    \end{jieda}
\end{enumerate}
\section*{抛物线的标准方程}

\begin{center}
\adjustbox{width=\textwidth}{
\begin{tabular}{|c|c|c|c|c|c|}
\hline
\phantom{X} & 标准方程 & ${y}^{2} = {2px}\left( {p > 0}\right)$ & ${y}^{2} =  - {2px}\left( {p > 0}\right)$ & ${x}^{2} = {2py}\left( {p > 0}\right)$ & ${x}^{2} =  - {2py}(p$  $> 0)$ \\
\cline{1-6}
\multicolumn{2}{|c|}{图形} &  \includegraphics[width=0.2\textwidth]{bodyxb/src/bo_d3q39gjef24c73d0aopg_0_498_340_181_146_0.jpg}  &  \includegraphics[width=0.2\textwidth]{bodyxb/src/bo_d3q39gjef24c73d0aopg_0_739_337_183_149_0.jpg}  &  \includegraphics[width=0.2\textwidth]{bodyxb/src/bo_d3q39gjef24c73d0aopg_0_968_342_180_139_0.jpg}  &  \includegraphics[width=0.2\textwidth]{bodyxb/src/bo_d3q39gjef24c73d0aopg_0_1189_341_167_141_0.jpg}  \\
\cline{1-6}
\multirow{6}{*}{性质} & 焦点 & \blkda{$\left(\dfrac{p}{2}, 0\right)$} & \blkda{$\left(-\dfrac{p}{2}, 0\right)$} & \blkda{$\left(0,\dfrac{p}{2}\right)$} & \blkda{$\left(0,-\dfrac{p}{2}\right)$} \\
\cline{2-6}
 & 准线 & $x =  - \dfrac{p}{2}$ & \blkda{$x = \dfrac{p}{2}$} & $y =  - \dfrac{p}{2}$ & \blkda{$y = \dfrac{p}{2}$} \\
\cline{2-6}
 & 范围 & $x \geq  0,y \in  \mathbf{R}$ & $x \leq  0,y \in  \mathbf{R}$ & $y \geq  0,x \in  \mathbf{R}$ & $y \leq  0,x \in  \mathbf{R}$ \\
\cline{2-6}
 & 对称轴 & \multicolumn{2}{|c|}{$x$ 轴} & \multicolumn{2}{|c|}{$y$ 轴} \\
\cline{2-6}
 & 顶点 & \multicolumn{4}{|c|}{\blkda{$(0, 0)$}} \\
\cline{2-6}
 & 离心率 & \multicolumn{4}{|c|}{$e = $\blkda{1}} \\
\cline{1-6}
\hline
\end{tabular}
}
\end{center}

A 组：

\begin{enumerate}
    \item 抛物线 $y = 2{x}^{2}$ 的焦点坐标为\blkda{$\left(0, \dfrac{1}{8}\right)$}.
    \item  抛物线 ${y}^{2} =  - {2x}$ 的准线方程为\blkda{$x = \dfrac{1}{2}$}.
    \item 设抛物线的顶点在原点,准线方程为 $x =  - 2$ ,则抛物线的方程是\blkda{$y^2 = 8x$}.
    \item 若抛物线 ${y}^{2} = {2px}$ 的焦点与双曲线 ${x}^{2} - \dfrac{{y}^{2}}{3} = 1$ 的右焦点重合,则 $p =$ \blkda{4}.
    \begin{jieda}
        根据双曲线标准方程可知其右焦点为$(2, 0)$，故$p = 4$.
    \end{jieda}
    \item 已知抛物线 $C : {y}^{2} = {4x}$ 的焦点为 $F$ ,点 $P\left( {m,n}\right)$ 在抛物线 $C$ 上. 若 $\left| {PF}\right|  = 3$ ,则 $m =$ \blkda{2}.
    \item 已知抛物线 $C : {y}^{2} = {4x}$ 的焦点为 $F$ ,点 $A\text{ 、 }B$ 是抛物线 $C$ 上不同的两点,且$A\text{ 、 }B$ 中点的横坐标为 2,则 $\left| {AF}\right|  + \left| {BF}\right|  =$ \blkda{6}.
    \begin{jieda}
        设 $A\left( {{x}_{1},{y}_{1}}\right) ,B\left( {{x}_{2},{y}_{2}}\right)$ ,由 $A,B$ 中点的横坐标为 2,可得 ${x}_{1} + {x}_{2} = 4$，所以 $\left| {AF}\right|  + \left| {BF}\right|  = {x}_{1} + 1 + {x}_{2} + 1 = 6$ .
    \end{jieda}
    \item 抛物线 $y = a{x}^{2}\left( {a < 0}\right)$ 的焦点到准线的距离是\blkda{$-\dfrac{1}{2a}$}.
    \begin{jieda}
        $x^2=\dfrac{1}{a}y$，$p = -\dfrac{1}{2a}$
    \end{jieda}
    \item 若顶点在原点,焦点在 $y$ 轴上的抛物线上点 $P\left( {m, - 2}\right)$ 到焦点距离为 4 ,则抛物线的焦点坐标是\blkda{$(0, -2)$}.
    \begin{jieda}
        根据抛物线定义可知$P$到准线的距离为$4$，即准线为$y = 2$，故焦点坐标为$(0, -2)$.
    \end{jieda}

\end{enumerate}
B 组：
\begin{enumerate}
    \item 已知抛物线 ${y}^{2} = {8x}$ 上一点 $P$ 到焦点的距离为 5,则点 $P$ 到 $x$ 轴的距离为\blkda{$2\sqrt{6}$}.
    \begin{jieda}
        抛物线方程 ${y}^{2} = {8x}$ ,则焦点坐标为 $F\left( {2,0}\right)$ ,准线方程为 $x =  - 2$，由抛物线的定义可知,点 $P$ 到准线 $x =  - 2$ 的距离为 5，所以 ${x}_{P} + 2 = 5$ 
        
        故${x}_{P} = 3$ ,代入 ${y}^{2} = {8x}$，则 ${y}_{P} =  \pm  \sqrt{8 \times  3} =  \pm  2\sqrt{6}$
    \end{jieda}
    
    \item 若直线 $l$ 经过抛物线 $C : {y}^{2} = {4x}$ 的焦点且其一个方向向量为 $\vv{d} = \left( {1,1}\right)$ ,则直线 $l$ 的方程为\blkda[3]{$x - y - 1 = 0$}.
    \begin{jieda}
        因为抛物线 ${y}^{2} = {4x}$ 的焦点为$(1,0)$,方向向量为 $\vv{d} = \left( {1,1}\right)$ 的直线 $l$ 的斜率为 1，故直线 $l$ 的方程是 $y - 0 = 1 \cdot  \left( {x - 1}\right)$ ,即 $x - y - 1 = 0$
    \end{jieda}
    
    \item 已知抛物线 ${y}^{2} = {2px}$ 的上点 $M$ 的横坐标为 2. 若点 $M$ 到焦点 $F$ 的距离为 5,则 $p$ 的值为\blkda{6}.
    \begin{jieda}
        $5 = |MF| = x_M + \dfrac{p}{2} = 2 + \dfrac{p}{2}$
    \end{jieda}
    \item 设圆 ${x}^{2} + {y}^{2} - {6x} - 7 = 0$ 与抛物线 ${y}^{2} = {2px}\left( {p > 0}\right)$ 的准线相切,则 $p =$ \blkda{2}.
    \begin{jieda}
        由题意可知准线为$x = -1$，故$p = 2$.
    \end{jieda}
    \item 若动圆 $M$ 经过点 $A\left( {3,0}\right)$ ,且与直线 $l : x =  - 3$ 相切,则圆心 $M$ 的轨迹方程为\blkda{$y^2 = 12x$}.
    
    \item 已知点 $F\left( {-\dfrac{1}{4},0}\right)$ ,直线 $l : x = \dfrac{1}{4}$ ,点 $B$ 是直线 $l$ 上的动点,若过 $B$ 垂直于 $y$ 轴的直线与线段 ${BF}$ 的垂直平分线交于点 $M$ ,则点 $M$ 的轨迹方程是\blkda{$y^2 = -x$}.
    \begin{jieda}
        根据垂直平分线的性质可知$|MF| = |MB|$，又因为$M$在过 $B$ 垂直于 $y$ 轴的直线上，所以$|MB|$即$M$到直线$l$的距离，故$M$的轨迹为抛物线，$F$为其焦点，$l$为其准线. 故轨迹方程为$y^2 = -x$.
    \end{jieda}

    \item 已知抛物线 ${x}^{2} = {2y}$ 上距离点 $A\left( {0,a}\right) \left( {a > 0}\right)$ 最近的点恰好是其顶点,则 $a$ 的取值范围是\blkda{$(0, 1]$}.
    \begin{jieda}
        设 $P\left( {x,y}\right)$ 为抛物线 ${x}^{2} = {2y}$ 上任意一点，则 $\left| {PA}\right|  = \sqrt{{x}^{2} + {\left( y - a\right) }^{2}} = \sqrt{{2y} + {y}^{2} - {2ay} + {a}^{2}} = \sqrt{[y - \left( {a - 1}\right)]^{2} + {2a} - 1}$ 
        
        因为 $a > 0$，所以对称轴 $a - 1 >  - 1$，又由于 $y \geq  0$ ,且 $\left| {PA}\right|$ 最小时$y = 0$，所以 $- 1 < a - 1 \leq  0$，所以 $0 < a \leq  1$ .
    \end{jieda}
    
    \item 圆心在抛物线 ${y}^{2} = {2x}$ 上,且与 $x$ 轴和该抛物线的准线都相切的一个圆的方程是\dotfill{(\kh{A})}
    \xx{${x}^{2} + {y}^{2} - x - {2y} + \dfrac{1}{4} = 0$}{${x}^{2} + {y}^{2} - x - {2y} - \dfrac{1}{4} = 0$}{${x}^{2} + {y}^{2} + x - {2y} + 1 = 0$}{${x}^{2} + {y}^{2} - x - {2y} + 1 = 0$}
    
    \item 已知过抛物线 ${y}^{2} = {2px}\left( {p > 0}\right)$ 的焦点 $F$ 的直线与抛物线交于 $A\text{ 、 }B$ 两点,以 ${AB}$ 为直径的圆 $M$ 与抛物线的准线 $l$ 的位置关系为\dotfill{(\kh{B})}
    \xx{相交}{相切}{相离}{不确定}
    \begin{jieda}
        如图所示,设线段 ${AB}$ 的中点为 $M$ ,且点 $A$ 、 $B$ 、 $M$ 在抛物线的准线 $l$ 上的射影依次是点 ${A}_{1}\text{ 、 }{B}_{1}\text{ 、 }{M}_{1}$ ,则 $M{M}_{1}$ 是直角梯形 $A{A}_{1}{B}_{1}B$ 的中位线.
        
        \begin{center}
        \includegraphics[width=0.25\textwidth]{bodyxb/src/bo_d19nc3f7aajc73895nqg_3_1150_1122_357_360_0.jpg}
        \end{center}
        
        因为点 $A$ 、 $B$ 在抛物线上,则由抛物线定义可知$\left| {A{A}_{1}}\right|  = \left| {AF}\right| ,\left| {B{B}_{1}}\right|  = \left| {BF}\right|$
        
        于是$\left| {M{M}_{1}}\right|  = \dfrac{1}{2}\left( {\left| {A{A}_{1}}\right|  + \left| {B{B}_{1}}\right| }\right)  = \dfrac{1}{2}\left( {\left| {AF}\right|  + \left| {BF}\right| }\right)  = \dfrac{1}{2}\left| {AB}\right|$
        
        即点 $M$ 到准线的距离等于圆的半径. 因此,以线段 ${AB}$ 为直径的圆与该抛物线的准线相切.
    \end{jieda}
    
    \item 若点 $P$ 是抛物线 ${y}^{2} = {2x}$ 上的一个动点,则点 $P$ 到点$(0,2)$的距离与 $P$ 到该抛物线准线的距离之和的最小值为\dotfill{(\kh{A})}
    \xx{$\dfrac{\sqrt{17}}{2}$}{$3$}{$\sqrt{5}$}{$\dfrac{9}{2}$}
    \begin{jieda}
        设抛物线焦点为$F\left(\dfrac{1}{2}, 0\right)$，因为$P,F$位于抛物线两侧，则根据抛物线定义可知题中目标值即$|PF| = \dfrac{\sqrt{17}}{2}$
    \end{jieda}
    
    \item 若动点 $M$ 的坐标$(x, y)$满足方程 $5\sqrt{{\left( x - 1\right) }^{2} + {\left( y - 2\right) }^{2}} = \left| {{3x} + {4y} + {12}}\right|$ ,则点 $M$ 的轨迹是\dotfill{(\kh{D})}
    \xx{圆}{椭圆}{双曲线}{抛物线}
    \begin{jieda}
        $5\sqrt{{\left( x - 1\right) }^{2} + {\left( y - 2\right) }^{2}} = \left| {{3x} + {4y} + {12}}\right| \Longleftrightarrow \sqrt{{\left( x - 1\right) }^{2} + {\left( y - 2\right) }^{2}} = \dfrac{\left| {{3x} + {4y} + {12}}\right|}{5}$，即$M$到$(1,2)$的距离等于其到直线$3x+4y+12 = 0$的距离.
    \end{jieda}
    
    \item 已知抛物线 $C : {y}^{2} = {4x}$ 的焦点为 $F$ .
    \begin{enumerate}
        \item 求抛物线 $C$ 的焦点坐标和准线方程
        \item 过焦点 $F$ 的直线 $l$ 与抛物线交于 ${A}$ 、 $B$ 两点,若 $\left| {AF}\right|  = 3$ ,求线段 $AB$ 的长
    \end{enumerate}
    \begin{jieda}
        \begin{enumerate}
            \item 因为 ${2p} = 4$ ,解得 $p = 2$ ,则抛物线 $C$ 的焦点坐标 $F\left( {1,0}\right)$ ,准线方程为 $x =  - 1$ ;
            
            \item 不妨设 $A\left( {{x}_{1},{y}_{1}}\right) ,B\left( {{x}_{2},{y}_{2}}\right)$ ,
            
            因为 $\left| {AF}\right|  = {x}_{1} + 1 = 3$ ,所以 ${x}_{1} = 2$，当 $x = 2$ 时,解得 $y =  \pm  2\sqrt{2}$ ,
            
            不妨令 $A\left( {2,2\sqrt{2}}\right) ,{k}_{AF} = \dfrac{2\sqrt{2}}{2 - 1} = 2\sqrt{2}$，此时直线 $l$ 的方程为 $y = 2\sqrt{2}\left( {x - 1}\right)$
            
            联立 $\left\{  \begin{array}{l} {y}^{2} = {4x} \\  y = 2\sqrt{2}\left( {x - 1}\right)  \end{array}\right.$ ,消去 $y$ 并整理得 $2{x}^{2} - {5x} + 2 = 0$ ,
            
            由韦达定理得 ${x}_{1} + {x}_{2} = \dfrac{5}{2}$，则 $\left| {AB}\right|  = \left| {AF}\right|  + \left| {BF}\right|  = {x}_{1} + 1 + {x}_{2} + 1 = \dfrac{9}{2}$ .
        \end{enumerate}
    \end{jieda}
    \item 如图是抛物线形拱桥,设水面宽 $\left| {AB}\right|  = {18}$ 米,拱顶距离水面 8 米,一货船在水面上的部分横断面为一矩形 ${CDEF}$ . 若 $\left| {CD}\right|  = 9$ 米,那么 $\left| {DE}\right|$ 不超过多少米才能使货船通过拱桥？
    
    \begin{center}
    \includegraphics[width=0.3\textwidth]{bodyxb/src/bo_d3q39gjef24c73d0aopg_2_241_428_391_185_0.jpg}
    \end{center}
    \begin{jieda}
        如图所示,以点 $O$ 为坐标原点,过点 $O$ 且平行于 ${AB}$ 的直线为 $x$ 轴,线段 ${AB}$ 的垂直平分线为 $y$ 轴建立平面直角坐标系,则 $B\left( {9, - 8}\right)$ .
        
        \begin{center}
        \includegraphics[width=0.2\textwidth]{bodyxb/src/bo_d3q39tk601uc738l6gig_3_244_496_332_264_0.jpg}
        \end{center}
        
        设抛物线方程为 ${x}^{2} =  - {2py}\left( {p > 0}\right)$ .
        
        因为$B$ 点在抛物线上, 故${81} =  - {2p} \cdot  \left( {-8}\right)$，因此$p = \dfrac{81}{16}$，所以抛物线的方程为 ${x}^{2} =  - \dfrac{81}{8}y$ .
        
        当 $x = \dfrac{9}{2}$ 时, $y =  - 2$ ,即 $\left| {DE}\right|  = 8 - 2 = 6$，所以$\left| {DE}\right|$ 不超过 6 米才能使货船通过拱桥.
    \end{jieda}
    
    \item 如图,直线 $l_{1}$ 和 $l_{2}$ 相交于点 ${M},l_{1} \bot  l_{2}$ ,点 ${N} \in  l_{2}$ ,以 $A$ 、 $B$ 为端点的曲线段 $C$ 上任一点到 ${l}_1$ 的距离与到点 $N$ 的距离相等. 若 $\bigtriangleup  {AMN}$ 为锐角三角形, $\left| {AM}\right|  = \sqrt{17}$ , $\left| {AN}\right|  = 3$ ,且 $\left| {NB}\right|  = 6$ ,建立适当的坐标系,求曲线段 $C$ 的方程.
    
    \begin{center}
    \includegraphics[width=0.2\textwidth]{bodyxb/src/bo_d3q39gjef24c73d0aopg_2_1118_762_293_278_0.jpg}
    \end{center}
    \begin{jieda}
        根据抛物线定义可知$C$为以$N$为焦点，以$l_1$为准线的抛物线的一部分，以$MN$中点为原点，以$\vv{MN}$为$x$轴正方向建立平面直角坐标系，设抛物线方程为$y^2 = 2px$
        
        故$A$点坐标为$\left(3-\dfrac{p}{2}, \sqrt{6p - p^2}\right)$，故$17 = |AM|^2 = 9 + 6p - p^2$，解得$p = 2$或$4$，即$A(2, 1)$或$(1, 2\sqrt{2})$，根据$\bigtriangleup  {AMN}$ 为锐角三角形可知$A(1, 2\sqrt{2})$

        故$y^2 = 8x, N(2, 0)$，根据$\left| {NB}\right|  = 6$可知$x_B = 4$，故$B(4, 4\sqrt{2})$
        
        因此$C$为$y^2 = 8x, y \in [2\sqrt{2}, 4\sqrt{2}]$
    \end{jieda}
    
    \item 已知点 $M$ 到点 $F\left( {1,0}\right)$ 的距离与到直线 $x = 3$ 的距离之和等于 4,求点 $M$ 的轨迹方程.
    \begin{jieda}
        $\left\{  \begin{array}{ll} {y}^{2} =  - {12x} + {48}, & 3 \leq  x \leq  4 \\  {y}^{2} = {4x}, & 0 \leq  x \leq  3 \end{array}\right.$
    \end{jieda}
    
    \item 如图,过抛物线 ${y}^{2} = {2px}\left( {p > 0}\right)$ 上一点 $P\left( {{x}_{0},{y}_{0}}\right) \left( {{y}_{0} > 0}\right)$ 作两条直线分别交抛物线于 $A\left( {{x}_{1},{y}_{1}}\right) \text{ 、 }B\left( {{x}_{2},{y}_{2}}\right)$ . 
    
    \begin{center}
    \includegraphics[width=0.2\textwidth]{bodyxb/src/bo_d3q39gjef24c73d0aopg_2_197_1610_292_304_0.jpg}
    \end{center}
    \begin{enumerate}
        \item 求该抛物线上纵坐标为 $\dfrac{p}{2}$ 的点到其焦点 $F$ 的距离
        \item 当 ${PA}$ 与 ${PB}$ 的斜率存在且倾斜角互补时,求 $\dfrac{{y}_{1} + {y}_{2}}{{y}_{0}}$ 的值,并证明此时直线 ${AB}$ 的斜率是非零常数.
    \end{enumerate}
    \begin{jieda}
        \begin{enumerate}
            \item $\left(\dfrac{p}{2}\right)^2 = 2px$，解得$x = \dfrac{p}{8}$，故其到焦点距离为$\dfrac{p}{8} + \dfrac{p}{2} = \dfrac{5p}{8}$
            \item 
            $y_0^2 = 2px_0$，故$P\left(\dfrac{y_0^2}{2p}, y_0\right)$，设直线$PA$为$x = m(y - y_0) + \dfrac{y_0^2}{2p}$，则$PB$为$x = -m(y - y_0) + \dfrac{y_0^2}{2p}$

            联立$PA$与抛物线方程$\left\{\begin{array}{l}
                 y^2 = 2px \\
                 x = m(y - y_0) + \dfrac{y_0^2}{2p}
            \end{array}\right.$，消去$x$可得$y^2 = 2p\left(m(y - y_0) + \dfrac{y_0^2}{2p}\right)$，即$y^2 -2pmy + (2pmy_0-y_0^2) = 0$，故根据韦达定理以及$P$点纵坐标可知$A$点纵坐标$y_1 = 2pm - y_0$

            同理$y_2 = -2pm-y_0$，故$\dfrac{{y}_{1} + {y}_{2}}{{y}_{0}} =  - 2$
            
            ${k}_{AB} = \dfrac{y_2-y_1}{x_2-x_1} = \dfrac{y_2-y_1}{\dfrac{y_2^2}{2p} - \dfrac{y_1^2}{2p},} = \dfrac{2p}{y_1+y_2}= - \dfrac{p}{{y}_{0}}$
        \end{enumerate}
    \end{jieda}
\end{enumerate}

C 组：
\begin{enumerate}
    \item 过抛物线 ${y}^{2} = 2px\left( {p > 0}\right)$ 的焦点 $F$ 作直线交抛物线于 $A,B$ 两点,以 ${AB}$ 为直径的圆与准线 $l$ 的公共点为 $M$ ,若 $\angle {AMF} = {60}^{ \circ  }$ ,则 $\angle {MFO}$ 的大小为\dotfill{(\kh{B})}
    \xx{$\ang{15}$}{$\ang{30}$}{$\ang{45}$}{不确定}
    \begin{jieda}
        如图,取 ${AB}$ 的中点 $G$ ,连接 ${MG}$ ,

        \begin{center}
        \includegraphics[width=0.2\textwidth]{bodyxb/src/bo_d3q39tk601uc738l6gig_3_663_1610_317_338_0.jpg}
        \end{center}
        
        则以 ${AB}$ 为直径的圆与准线 $l$ 切于点 $M$，根据抛物线性质, ${MG}//x$ 轴,且 ${MF} \bot  {AB}$ ,
        
        $\because \angle {AMF} = {60}^{ \circ  }$ ,
        
        $\therefore \angle {GAM} = \angle {GMA} = {30}^{ \circ  }$ ,$\therefore \angle {MFO} = \angle {GMF} = {30}^{ \circ  }$ .
    \end{jieda}
    
    \item 抛物线 $C : y = \dfrac{1}{4}{x}^{2}$ ,过焦点 $F$ 的直线 $l$ 与抛物线 $C$ 交于 $A,B$ 两点,若 $3\vv{AF} = \vv{FB}$ , 则直线 ${AB}$ 的倾斜角为\dotfill{(\kh{C})}
    \xx{${30}^{ \circ  }$}{${60}^{ \circ  }$}{${30}^{ \circ  }$ 或 ${150}^{ \circ  }$}{${60}^{ \circ  }$ 或 ${120}^{ \circ  }$}
    \begin{jieda}
        由题意可知 $F\left( {0,1}\right)$ ,不妨设 $l : y = {kx} + 1,A\left( {{x}_{1},{y}_{1}}\right) ,B\left( {{x}_{2},{y}_{2}}\right)$ ,
        
        联立直线与抛物线方程得 $\left\{  \begin{array}{l} {x}^{2} = {4y} \\  y = {kx} + 1 \end{array} \Longrightarrow  {x}^{2} - {4kx} - 4 = 0 \Longrightarrow  \left \{\begin{array}{l}
             {x}_{1} + {x}_{2} = {4k}  \\
             {x}_{1}{x}_{2} =  - 4
        \end{array}\right. \right.$ ,
        
        又 $\vv{AF} = \left( {-{x}_{1},1 - {y}_{1}}\right) ,\vv{FB} = \left( {{x}_{2},{y}_{2} - 1}\right)$ ,而 $3\vv{AF} = \vv{FB}$ ,
        
        则 $- 3{x}_{1} = {x}_{2} \Longrightarrow  3{x}_{1}^{2} = 4$ ,即 $\left\{  {\begin{array}{l} {x}_{1} = \dfrac{2\sqrt{3}}{3} \\  {x}_{2} =  - 2\sqrt{3} \end{array} \Longrightarrow  k =  - \dfrac{\sqrt{3}}{3}}\right.$ 或 $\left\{  {\begin{array}{l} {x}_{1} =  - \dfrac{2\sqrt{3}}{3} \\  {x}_{2} = 2\sqrt{3} \end{array} \Longrightarrow  k = \dfrac{\sqrt{3}}{3}}\right.$ ,
        
        所以直线的倾斜角为 ${30}^{ \circ  }$ 或 ${150}^{ \circ  }$ .
    \end{jieda}
    
    \item 已知抛物线 $C : {y}^{2} = {4x}$ 的焦点为 $F$ ,准线为 $l$ ,过点 $F$ 的直线与抛物线交于 $P\left( {x}_{1}\right.$ , $\left. {y}_{1}\right)$ , $Q\left( {{x}_{2},{y}_{2}}\right)$ 两点,点 $P$ 在 $l$ 上的射影为 ${P}_{1}$ ,则下列说法错误的是\dotfill{(\kh{D})}
    \xx{若 ${x}_{1} + {x}_{2} = 6$ ,则 $\left| {PQ}\right|  = 8$}{以 ${PQ}$ 为直径的圆与准线 $l$ 相切}{设 $M\left( {0,1}\right)$ ,则 $\left| {PM}\right|  + \left| {P{P}_{1}}\right|  \geq  \sqrt{2}$}{过点 $M\left( {0,1}\right)$ 与抛物线 $C$ 有且仅有一个公共点的直线至多有 2 条}
    \begin{jieda}
        对于选项 ${A}$ ,因为 $p = 2$ ,所以 ${x}_{1} + {x}_{2} + 2 = \left| {PQ}\right|$ ,则 $\left| {PQ}\right|  = 8$ ,故 ${A}$ 正确;
        
        对于选项 $B$,由抛物线焦点弦的性质可知, ${B}$ 正确；
        
        对于选项 ${C}$ ,因为 $F\left( {1,0}\right)$ ,所以 $\left| {PM}\right|  + \left| {P{P}_{1}}\right|  = \left| {PM}\right|  + \left| {PF}\right|  \geq  \left| {MF}\right|  = \sqrt{2}$ ,故 ${C}$ 正确;
        
        对于选项 ${D}$ ,显然直线 $x = 0,y = 1$ 与抛物线只有一个公共点,设过 $M$ 的直线方程为 $y = {kx}$  $+ 1\left( {k \neq  0}\right)$ ,
        
        联立 $\left\{  \begin{array}{l} y = {kx} + 1, \\  {y}^{2} = {4x}, \end{array}\right.$ 可得 ${k}^{2}{x}^{2} + \left( {{2k} - 4}\right) x + 1 = 0$ ,令 $\Delta  = 0$ ,得 $k = 1$ ,所以直线 $y = x + 1$ 与抛物线也只有一个公共点,此时有三条直线符合题意,故 ${D}$ 错误.

    \end{jieda}
    
    \item 已知抛物线 ${y}^{2} = {2px}\left( {p > 0}\right)$ 上三点 $A\left( {{x}_{1},{y}_{1}}\right) ,B\left( {1,2}\right) ,C\left( {{x}_{2},{y}_{2}}\right) ,F$ 为抛物线的焦点, 则下列说法错误的是\dotfill{(\kh{C})}
    \xx{抛物线的准线方程为 $x =  - 1$}{若 $\vv{FA} + \vv{FB} + \vv{FC} = \vv{0}$ ,则 $2\left| \vv{FB}\right|  = \left| \vv{FA}\right|  + \left| \vv{FC}\right|$}{若 $A,F,C$ 三点共线,则 ${y}_{1}{y}_{2} =  - 1$}{若 $\left| {AC}\right|  = 6$ ,则 ${AC}$ 的中点到 $y$ 轴距离的最小值为 2}
    \begin{jieda}
        把点 $B\left( {1,2}\right)$ 代入抛物线 ${y}^{2} = {2px}$ ,得 $p = 2$ ,所以抛物线的准线方程为 $x =  - 1$ ,故 ${A}$ 正确;
        
        因为 $A\left( {{x}_{1},{y}_{1}}\right) ,B\left( {1,2}\right) ,C\left( {{x}_{2},{y}_{2}}\right) ,F\left( {1,0}\right)$ ,所以 $\vv{FA} = \left( {{x}_{1} - 1,{y}_{1}}\right) ,\vv{FB} = \left( {0,2}\right) ,\vv{FC} = \left( {{x}_{2} - 1}\right.$ , $\left. {y}_{2}\right)$ ,又由 $\vv{FA} + \vv{FB} + \vv{FC} = \vv{0}$ ,得 ${x}_{1} + {x}_{2} = 2$ ,
        
        所以 $\left| \vv{FA}\right|  + \left| \vv{FC}\right|  = {x}_{1} + 1 + {x}_{2} + 1 = 4 = 2\left| \vv{FB}\right|$ ,故 ${B}$ 正确;
        
        因为 $A,F,C$ 三点共线,所以直线 ${AC}$ 是焦点弦,所以 ${y}_{1}{y}_{2} =  - {p}^{2} =  - 4$ ,故 ${C}$ 不正确; 
        
        
        设 ${AC}$ 的中点为 $M\left( {{x}_{0},{y}_{0}}\right)$，因为 $\left| {AF}\right|  + \left| {CF}\right|  \geq  \left| {AC}\right| ,\left| {AF}\right|  + \left| {CF}\right|  = {x}_{1} + 1 + {x}_{2} + 1 = 2{x}_{0} + 2$，所以 $2{x}_{0} + 2 \geq  6$ ,得 ${x}_{0} \geq  2$，即 ${AC}$ 的中点到 $y$ 轴距离的最小值为 2,故 ${D}$ 正确.
    \end{jieda}
\end{enumerate}

\section*{抛物线的性质}
\subsection*{抛物线的性质 (1)}

A 组：
\begin{enumerate}
    \item 若抛物线 $C : {y}^{2} = {2px}\left( {p > 0}\right)$ 的焦点在直线 $x + {2y} - 2 = 0$ 上,则 $p$ 等于\blkda{4}.
    
    \item 抛物线 ${y}^{2} = {4x}$ 上一点 $M$ 到焦点 $F$ 的距离是 10 ,则点 $M$ 到 $y$ 轴的距离是\blkda{9}.

    \item 已知 $F$ 是椭圆 ${C}_{1} : \dfrac{{x}^{2}}{4} + \dfrac{{y}^{2}}{3} = 1$ 与抛物线 ${C}_{2} : {y}^{2} = {2px}\left( {p > 0}\right)$ 的一个共同焦点, ${C}_{1}$ 与 ${C}_{2}$ 相交于 $A,B$ 两点,则线段 ${AB}$ 的长等于\blkda{$\dfrac{4\sqrt{6}}{3}$}.
    \begin{jieda}
        公共焦点为$F(1, 0)$，故抛物线方程为$y^2 = 4x$，易知$A, B$两点关于$x$轴对称，联立椭圆和抛物线方程，消去$x$可得$3y^4+64y^2-192 = 0$，解得$y^2 = \dfrac{8}{3}$，故$|AB| = \dfrac{4\sqrt{6}}{3}$
    \end{jieda}

    \item 设 $F$ 为抛物线 ${y}^{2} = {4x}$ 的焦点,若点 $P$ 在抛物线上移动,点 $Q$ 在 ${\left( x - 1\right) }^{2} + {y}^{2} = \dfrac{1}{9}$ 上移动,则 $\left| {PQ}\right|$ 的最小值为\blkda{$\dfrac{2}{3}$}.

    \item 已知 $F$ 是抛物线 $C : {y}^{2} = {4x}$ 的焦点, $P$ 是抛物线 $C$ 上一动点, $Q$ 是曲线 ${x}^{2} + {y}^{2} - {8x} - {2y} + {16} = 0$ 上一动点,则 $\left| {PF}\right|  + \left| {PQ}\right|$ 的最小值为\blkda{4}.

\end{enumerate}
B 组：
\begin{enumerate}
    \item 若 $P\left( {m,n}\right)$ 是焦点为 $F$ 的抛物线 ${y}^{2} = {2px}\left( {p > 0}\right)$ 上的一点,则 $\left| {PF}\right|  =$ \blkda{$m+\dfrac{p}{2}$}. 
    
    \item 抛物线 ${y}^{2} = {8x}$ 上的点到它的焦点的距离的最小值等于\blkda{2}.

    \item 已知 ${AB}$ 是抛物线 ${y}^{2} = {2px}$ 的焦点弦, $O$ 是抛物线的顶点. 若 $\left| {AB}\right|  = a$ ,则 $\bigtriangleup  {AOB}$ 的面积是\blkda{$\dfrac{p\sqrt{2pa}}{4}$}.
    \begin{jieda}
        设$A(x_1, y_1)$，$B(x_2, y_2)$，故$S = \dfrac{1}{2} \times \dfrac{p}{2} \times |y_1-y_2| = \dfrac{p}{4}\cdot \sqrt{(y_1+y_2)^2-4y_1y_2}$

        设直线$AB$为$x = my+\dfrac{p}{2}$，联立$AB$和抛物线方程$\left\{\begin{array}{l}
             x = my+\dfrac{p}{2}  \\
             y^2= 2px
        \end{array}\right.$，消去$x$可得$y^2= 2p\left(my+\dfrac{p}{2}\right)$，即$y^2 - 2pmy - p^2 = 0$，因为$AB$过焦点，所以必有$\Delta > 0$.

        根据韦达定理有$\left\{\begin{array}{l}
             y_1+y_2 = 2pm \\
             y_1y_2 = -p^2
        \end{array}\right.$

        又因为$a = |AB| = x_1+x_2+p = \dfrac{y_1^2}{2p} + \dfrac{y_2^2}{2p}+p$，所以$(y_1+y_2)^2 = (2pa-2p^2) -2p^2 = 2pa - 4p^2$，因此$S = \dfrac{p\sqrt{2pa}}{4}$

        \ps 若利用结论$|AB| = \dfrac{2p}{sin^2\alpha}$可知$sin\alpha = \sqrt{\dfrac{2p}{a}}$，故$tan\alpha = \pm \sqrt{\dfrac{2p}{a-2p}}$，不妨取正，因此$AB$方程为$y = \sqrt{\dfrac{2p}{a-2p}}\left(x-\dfrac{p}{2}\right)$，则$O$到$AB$的距离$h = \dfrac{\left|\dfrac{p}{2}\sqrt{\dfrac{2p}{a-2p}}\right|}{\sqrt{1+\dfrac{2p}{a-2p}}} = \dfrac{p}{2}\sqrt{\dfrac{2p}{a}}$，故$S = \dfrac{ah}{2} = \dfrac{p\sqrt{2pa}}{4}$
    \end{jieda}

    \item 抛物线的顶点在原点,对称轴为坐标轴,焦点在直线 ${3x} - {4y} - {12} = 0$ 上,则抛物线的方程是\blkda[5]{$y^2 = 16x$或$x^2 = -12y$}.
    \begin{jieda}
        \begin{dingautolist}{172}
            \item 若焦点为$(4, 0)$，则抛物线方程为$y^2 = 16x$
            \item 若焦点为$(0, -3)$，则抛物线方程为$x^2 = -12y$
        \end{dingautolist}
    \end{jieda}

    \item 若点 $P$ 到直线 $x =  - 1$ 的距离比它到点$(2,0)$的距离小 1 ,则点 $P$ 的轨迹方程为\blkda{$y^2 = 8x$}.

    \item 若动圆 $M$ 与圆 $C : {x}^{2} + {\left( y - 2\right) }^{2} = 1$ 相外切,且与 $x$ 轴相切,则圆心 $M$ 的轨迹方程为\blkda[3]{$x^2 = 6y - 3$}.
    \begin{jieda}
        $M$到$(0,2)$的距离与到$y = -1$的距离相等，故是顶点在$\left(0, \dfrac{1}{2}\right)$，$p = 3$，开口向上的抛物线，故方程为$x^2 = 6\left(y-\dfrac{1}{2}\right)$，即$x^2 = 6y - 3$

        \ps 代数推导：设点$M(x,y)$，$\sqrt{x^2+(y-2)^2} = y+1$，平方可得$x^2 = 6y-3$.
    \end{jieda}

    \item 已知 $O$ 为坐标原点, $F$ 为抛物线 ${y}^{2} = {4x}$ 的焦点, $A$ 为此抛物线上的一点. 若 $\vv{OA} \cdot  \vv{AF} =  - 4$ ,则点 $A$ 的坐标为\blkda{$(1, \pm 2)$}.
    \begin{jieda}
        设$A\left(\dfrac{y_0^2}{4}, y_0\right)$，故$\vv{OA} = \left(\dfrac{y_0^2}{4}, y_0\right)$，$\vv{AF} = \left(1-\dfrac{y_0^2}{4}, -y_0\right)$，因此$-4 = \vv{OA} \cdot  \vv{AF} = \dfrac{y_0^2}{4} - \dfrac{y_0^4}{16} - y_0^2$，解得$y_0^2 = 4$，故$A$的坐标为$(1, \pm 2)$.
    \end{jieda}

    \item 若等边三角形的三个顶点都在抛物线 ${y}^{2} = {2px}$ 上,且其中一个顶点为抛物线的顶点,则正三角形的边长为\blkda{$4\sqrt{3}p$}. (用 ${p}$ 表示)
    \begin{jieda}
        易知$x$为正三角形的中线，设$x$轴上方的顶点坐标为$\left(\dfrac{y_0^2}{2p}, y_0\right)$，则有$\sqrt{\left(\dfrac{y_0^2}{2p}\right)^2+y_0^2} = 2y_0$，即$\left(\dfrac{y_0^2}{2p}\right)^2+y_0^2 = 4y_0^2$，故$y_0^2 = 12p$，故正三角形边长为$4\sqrt{3}p$
    \end{jieda}

    \item 设抛物线 ${x}^{2} = {4y}$ 的焦点为 ${F}$ ,点 $A\text{ 、 }B\text{ 、 }C$ 在抛物线上,若 $\vv{FA} + \vv{FB} + \vv{FC} = \vv{0}$ , 则 $\left| \vv{FA}\right|  + \left| \vv{FB}\right|  + \left| \vv{FC}\right|  =$ \blkda{6}.
    \begin{jieda}
        $p = 2, F$为$\triangle ABC$重心，$y_A+y_B+y_C = 3y_F = 3$，$\left| \vv{FA}\right|  + \left| \vv{FB}\right|  + \left| \vv{FC}\right|  = y_A+y_B+y_C + \dfrac{3p}{2} = 3+3 = 6$
    \end{jieda}

    \item 已知点 ${A}\left( {1,2}\right)$ ,过点 ${P}\left( {5, - 2}\right)$ 的直线与抛物线 ${y}^{2} = {4x}$ 相交于 ${B},{C}$ 两点, 则 $\bigtriangleup  {{ABC}}$ 是\dotfill{(\kh{A})}
    \xx{直角三角形}{钝角三角形}{锐角三角形}{不能确定}
    \begin{jieda}
        \begin{dingautolist}{172}
            \item 当 ${BC}$ 斜率不存在时,方程为 $x = 5$，代入抛物线方程 ${y}^{2} = {4x}$ 得$B\left( {5,2\sqrt{5}}\right) ,C\left( {5, - 2\sqrt{5}}\right)$所以 ${AB}$ 斜率是 ${k}_{AB} = \dfrac{2\sqrt{5} - 2}{5 - 1} = \dfrac{\sqrt{5} - 1}{2}$,${AC}$ 斜率是 ${k}_{AC} = \dfrac{-2\sqrt{5} - 2}{5 - 1} = \dfrac{-\sqrt{5} - 1}{2}$，所以 ${k}_{AB} \cdot  {k}_{AC} =  - 1$
            \item 当 ${BC}$ 斜率存在时, 显然不能为 0 , 否则与抛物线只有一个公共点，所以设方程为 $x - 5 = a\left( {y + 2}\right)$，代入抛物线方程化简得${y}^{2} - {4ay} - {8a} - {20} = 0$ 
            
            设 $B\left( {{x}_{1},{y}_{1}}\right), C\left( {{x}_{2},{y}_{2}}\right)$，根据韦达定理有$\left\{\begin{array}{l}
                 {y}_{1} + {y}_{2} = {4a}  \\
                 {y}_{1}{y}_{2} =  - {8a} - {20}
            \end{array}\right.$ ,
            ${k}_{AB} \cdot  {k}_{AC} = \dfrac{{y}_{1} - 2}{{x}_{1} - 1} \cdot  \dfrac{{y}_{2} - 2}{{x}_{2} - 1} = \dfrac{{y}_{1}{y}_{2} - 2\left( {{y}_{1} + {y}_{2}}\right)  + 4}{x_{1}x_{2} - \left( {x_{1} + x_{2}}\right)  + 1} = \dfrac{{y}_{1}{y}_{2} - 2\left( {{y}_{1} + {y}_{2}}\right)  + 4}{\dfrac{y_1^2y_2^2}{16} - \dfrac{y_1^2+y_2^2}{4} + 1} = \dfrac{-8a-20-8a+4}{(2a+5)^2-(4a^2+4a+10)+1} = -1$
        \end{dingautolist}
        综上可知,三角形 ${ABC}$ 是直角三角形
    \end{jieda}

    \item 已知抛物线的顶点在原点,焦点在 $x$ 轴上,且直线 $l : {2x} - y + 1 = 0$ 被此抛物线所截得的线段的长为 $\sqrt{15}$ ,求此抛物线的方程.
    \begin{jieda}
        设所求抛物线的方程为 ${y}^{2} = {ax}\left( {a \neq  0}\right)$ ,抛物线与直线 $l$ 的交点为 $A\left( {{x}_{1},{y}_{1}}\right) ,B\left( {{x}_{2},{y}_{2}}\right)$ . 由$\left\{  \begin{array}{l} {2x} - y + 1 = 0, \\  {y}^{2} = {ax}, \end{array}\text{ 得 }4{x}^{2} + \left( {4 - a}\right) x + 1 = 0 \right.$
        
        故${\left( {x}_{1} - {x}_{2}\right) }^{2} = {\left( {x}_{1} + {x}_{2}\right) }^{2} - 4{x}_{1}{x}_{2} = {\left( \dfrac{a - 4}{4}\right) }^{2} - 1 = \dfrac{{a}^{2} - {8a}}{16}$
        
        $\left| AB\right|^{2} = 5{\left( {x}_{1} - {x}_{2}\right) }^{2} = \dfrac{5\left( {{a}^{2} - {8a}}\right) }{16} = {15}$ . 解得 $a = {12}$ 或 $a =  - 4$ .
        
        因此,所求抛物线的方程为 ${y}^{2} = {12x}$ 或 ${y}^{2} =  - {4x}$ .
    \end{jieda}

    \item 一拱桥 (拱形抛物线) 的跨度为 52 米, 拱顶离水面 6.5 米. 河面上一艘竹排上有宽 4 米、高 6 米的大木箱,此竹排能否安全从桥下通过?
    \begin{jieda}
        以拱桥顶点为原点建立平面直角坐标系，设抛物线方程为 ${x}^{2} =  - {2py}$ . 
        
        将岸边拱桥地基的坐标$(-26, - 6.5)$代入抛物线方程中,易得 $p = {52}$ ,即抛物线方程为 ${x}^{2} = {104y}$ . 
        
        当 $y =  - {0.5}$ 时, $x =  \pm  2\sqrt{13}$ . 由 $4\sqrt{13} > 4$ ,知木箱能安全通过
    \end{jieda}

    \item 已知抛物线 ${y}^{2} = {2px}$ 的焦点为点 $F$ ,过点 $F$ 的直线与抛物线交于 $A\text{ 、 }B$ 两点. 求证: $\dfrac{1}{\left| FA\right| } + \dfrac{1}{\left| FB\right| } = \dfrac{2}{p}$ .
    \begin{jieda}
        焦点 $F$ 的坐标为 $\left( {\dfrac{p}{2},0}\right)$，设直线 ${AB}$ 的方程为$x = my+\dfrac{p}{2}$.
        
        联立方程组 $\left\{  \begin{array}{l} 
            x = my+\dfrac{p}{2}, \\  
            {y}^{2} = {2px}. \end{array}\right.$ 
        消去$x$得 ${y}^{2} = 2p\left(my+\dfrac{p}{2}\right)$，整理可得$y^2 - 2pmy - p^2 = 0$，由直线过焦点易知必有$A, B$两个交点
        
        根据韦达定理有${y}_{1}{y}_{2} =  - {p}^{2},{x}_{1}{x}_{2} = \dfrac{{y}_{1}^{2}}{2p} \cdot  \dfrac{{y}_{2}^{2}}{2p} = \dfrac{{\left( {y}_{1}{y}_{2}\right) }^{2}}{4{p}^{2}} = \dfrac{{p}^{4}}{4{p}^{2}} = \dfrac{{p}^{2}}{4}$ 

        由焦半径公式知, $\left| {AF}\right|  = {x}_{1} + \dfrac{p}{2},\left| {BF}\right|  = {x}_{2} + \dfrac{p}{2}$ . 所以$\dfrac{1}{\left| AF\right| } + \dfrac{1}{\left| BF\right| } = \dfrac{1}{{x}_{1} + \dfrac{p}{2}} + \dfrac{1}{{x}_{2} + \dfrac{p}{2}} = \dfrac{{x}_{1} + {x}_{2} + p}{{x}_{1}{x}_{2} + \dfrac{p}{2}\left( {{x}_{1} + {x}_{2}}\right)  + \dfrac{{p}^{2}}{4}}.$
        
        又 ${x}_{1}{x}_{2} = \dfrac{{p}^{2}}{4},{x}_{1} + {x}_{2} = \left| {AB}\right|  - p$，代入上式，得$\dfrac{1}{\left| AF\right| } + \dfrac{1}{\left| BF\right| } = \dfrac{\left| AB\right| }{\dfrac{{p}^{2}}{4} + \dfrac{p}{2}\left( {\left| {AB}\right|  - p}\right)  + \dfrac{{p}^{2}}{4}} = \dfrac{2}{p}$.
        
        所以, $\dfrac{1}{\left| AF\right| } + \dfrac{1}{\left| BF\right| }$ 为定值 $\dfrac{2}{p}$ .
    \end{jieda}

    \item 已知抛物线 ${y}^{2} = {2x},{P}$ 是抛物线上的一点.
    \begin{enumerate}
        \item 求点 ${P}$ 和圆 ${\left( x - 3\right) }^{2} + {y}^{2} = 1$ 上的点之间的最短距离
        \item 设点 $A\left( {a,0}\right)$ ,且 $\left| {PA}\right|  = d$ ,试求 $d$ 的最小值
    \end{enumerate}
    \begin{jieda}
        \begin{enumerate}
            \item 设圆心为$C(3, 0)$，则点 ${P}$ 和圆 ${\left( x - 3\right) }^{2} + {y}^{2} = 1$ 上的点之间的最短距离为$|PC| - 1$
            
            设 $P\left( {{x}_{0},{y}_{0}}\right) \left( {{x}_{0} \geq  0}\right)$ ,$|PC| = \sqrt{(x_0-3)^2+y_0^2} = \sqrt{x_0^2-4x_0+9} \geq \sqrt{5}$（$x_0 = 2$时取等）

            因此点 ${P}$ 和圆 ${\left( x - 3\right) }^{2} + {y}^{2} = 1$ 上的点之间的最短距离为$\sqrt{5}-1$
            \item 设 $P\left( {{x}_{0},{y}_{0}}\right) \left( {{x}_{0} \geq  0}\right)$ ,则 ${y}_{0}{}^{2} = 2{x}_{0}$ ,所以$d = \left| {PA}\right|  = \sqrt{{\left( {x}_{0} - a\right) }^{2} + {y}_{0}^{2}} = \sqrt{{\left( {x}_{0} - a\right) }^{2} + 2{x}_{0}} =  \sqrt{{\left\lbrack  {x}_{0} + \left( 1 - a\right) \right\rbrack  }^{2} + {2a} - 1}$
            
            因为 ${x}_{0} \geq  0$，所以
            \begin{dingautolist}{172}
                \item 当 $a < 1$ 时, $1 - a > 0$ ,此时当 ${x}_{0} = 0$ 时, ${d}_{\min } = \sqrt{{\left( 1 - a\right) }^{2} + {2a} - 1} = |a|$
                \item 当 $a \geq  1$ 时, $1 - a \leq  0$ ,此时当 ${x}_{0} = a - 1$ 时, ${d}_{\min } = \sqrt{{2a} - 1}$ .
            \end{dingautolist}
        \end{enumerate}
    \end{jieda}

    \item 设抛物线 ${y}^{2} = {2px}\left( {p > 0}\right)$ 的焦点为 ${F}$ ,经过点 ${F}$ 的直线交抛物线于 ${A}$ 、 ${B}$ 两点，点 ${C}$ 在抛物线的准线上,且 ${{BC}}//x$ 轴. 求证: 直线 ${{AC}}$ 经过原点 ${O}$ .
    \begin{jieda}
        焦点 $F$ 的坐标为 $\left( {\dfrac{p}{2},0}\right)$，设$A(x_1, y_1), B(x_2, y_2)$，故$C\left(-\dfrac{p}{2}, y_2\right)$，故$\vv{OA} = (x_1, y_1)$，$\vv{OC} = \left(-\dfrac{p}{2}, y_2\right)$
        
        $-\dfrac{p}{2}y_1 - x_1y_2 = -\dfrac{p}{2}y_1 - \dfrac{y_1^2y_2}{2p} = -y_1\left(\dfrac{p}{2}+\dfrac{y_1y_2}{2p}\right)$
        
        设直线 ${AB}$ 的方程为$x = my+\dfrac{p}{2}$，联立方程组 $\left\{  \begin{array}{l} 
            x = my+\dfrac{p}{2}, \\  
            {y}^{2} = {2px}. \end{array}\right.$ ，消去$x$得 ${y}^{2} = 2p\left(my+\dfrac{p}{2}\right)$，整理可得$y^2 - 2pmy - p^2 = 0$，由直线过焦点易知必有$A, B$两个交点，故根据韦达定理有${y}_{1}{y}_{2} =  - {p}^{2}$
            
            因此$-\dfrac{p}{2}y_1 - x_1y_2 = -y_1\left(\dfrac{p}{2}+\dfrac{y_1y_2}{2p}\right) = 0$，即$\vv{OA} // \vv{OC}$，所以直线 ${{AC}}$ 经过原点 ${O}$ .
    \end{jieda}
\end{enumerate}

\clearpage

\subsection*{抛物线的性质 (2)}
\begin{enumerate}
    \item 若过点 $A\left( {1,0}\right)$ ,且与 $y$ 轴的夹角为 $\dfrac{\pi }{6}$ 的直线与抛物线 ${y}^{2} = {4x}$ 交于 $P\text{ 、 }Q$ 两点, 则 $\left| {PQ}\right|  =$ \blkda{$\dfrac{16}{3}$}.
    \begin{jieda}
        $|PQ| = \dfrac{4}{sin^2\dfrac{\pi}{3}} = \dfrac{16}{3}$
    \end{jieda}
    
    \item 已知直线 ${l}_{1} : {4x} - {3y} + 6 = 0$ 和 ${l}_{2} : x =  - 1$ ,抛物线 ${y}^{2} = {4x}$ 上一动点 $P$ 到 ${l}_{1}$ 和 ${l}_{2}$ 的距离之和的最小值是\blkda{2}.
    \begin{jieda}
        $P$ 到 ${l}_2$的距离等于$P$到$(1, 0)$的距离，数形结合可知目标最小值为$(1, 0)$到$l_1$的距离$d = 2$
    \end{jieda}
    
    \item 已知抛物线的准线方程为 $x + y - 2 = 0$ ,焦点是 $F\left( {5,5}\right)$ ,则抛物线的顶点坐标是\blkda{$(3, 3)$}.
    \begin{jieda}
        顶点为焦点到准线的垂线段的中点.
    \end{jieda}

    \item 过抛物线 ${y}^{2} = {4x}$ 的焦点作直线交抛物线于 $A\left( {{x}_{1},{y}_{1}}\right) \text{ 、 }B\left( {{x}_{2},{y}_{2}}\right)$ 两点,如果 ${x}_{1} + {x}_{2} = 6$ ,则 $\left| {AB}\right|$ 等于\blkda{8}.
    \begin{jieda}
        $|AB| = x_1+x_2+p = 6+2 = 8$
    \end{jieda}

    \item 已知过抛物线的焦点 $F$ 的直线与抛物线交于 $A\text{ 、 }B$ 两点. 若点 $A\text{ 、 }B$ 在抛物线的准线上的射影分别是 ${A}_{1}\text{ 、 }{B}_{1}$ ,则 $\angle {A}_{1}F{B}_{1}$ 等于\blkda{$\dfrac{\pi}{2}$}.
    \begin{jieda}
        设$A(x_1, y_1)$，$B(x_2, y_2)$，则$A_1\left(-\dfrac{p}{2}, y_1\right)$，$B_1\left(-\dfrac{p}{2}, y_2\right)$

        $\vv{FA_1} = \left(-p, y_1\right)$，$\vv{FB_1} = \left(-p, y_2\right)$，根据焦点弦的性质可知$y_1y_2 = -p^2$，故$\vv{FA_1}\cdot\vv{FB_1} = 0$
    \end{jieda}


    \item 若直线 $y = {kx} + 2$ 与抛物线 ${y}^{2} = {4x}$ 只有一个公共点,则实数 $k$ 的值为\blkda{$0$或$\dfrac{1}{2}$}.
    \begin{jieda}
        \begin{dingautolist}{172}
            \item 若$y = {kx} + 2$为切线，设切点$\left(\dfrac{y_0^2}{4}, y_0\right)$，则切线为$y_0y = 2\left(x+\dfrac{y_0^2}{4}\right)$，即$y = \dfrac{2}{y_0}x+\dfrac{y_0}{2}$，故$\dfrac{y_0}{2} = 2$，即$k = \dfrac{2}{y_0} = \dfrac{1}{2}$
            \item 若$y = {kx} + 2$不是切线，则$k = 0$.
        \end{dingautolist}
    \end{jieda}


    \item 抛物线 $y = {x}^{2}$ 上到直线 ${2x} - y = 4$ 距离最小的点的坐标是\blkda{$(1, 1)$}.
    \begin{jieda}
        \begin{enumerate}
            \item 方法1

            易知抛物线 $y = {x}^{2}$ 和直线 ${2x} - y = 4$没有公共点，故直线$y = 2x+b$和抛物线相切时，切点即为所求.

            设切点为$(x_0, x_0^2)$，则该处切线为$\dfrac{x_0^2+y}{2} = x_0x$，即$y = 2x_0x-x_0^2$，当$x_0 = 1$时，切线和已知直线平行，切点为$(1, 1)$

            \item 方法2

            设抛物线上的点为$(x_0, x_0^2)$，其到直线的距离$d = \dfrac{|2x_0-x_0^2-4|}{\sqrt{5}} = \dfrac{x_0^2-2x_0+4}{\sqrt{5}}$，当$x_0 = 1$时取得最小值，故目标点坐标为$(1,1)$
        \end{enumerate}
    \end{jieda}


    \item 若直线 $l$ 与抛物线 ${y}^{2} = {16x}$ 相交所得的弦 ${AB}$ 被点 $P\left( {3,2}\right)$ 平分,则直线 $l$ 的方程为\blkda[3]{$y = 4x - 10$}.
    \begin{jieda}
        设$A(x_1, y_1)$，$B(x_2, y_2)$，则$\left\{\begin{array}{l}
             y_1^2 = 16x_1  \\
             y_2^2 = 16x_2 
        \end{array}\right.$，两式作差可得$(y_1+y_2)(y_1-y_2) = 16(x_1-x_2)$，故$k_{AB} = \dfrac{y_1-y_2}{x_1-x_2} = \dfrac{16}{y_1+y_2} = 4$，故$l: y = 4(x-3)+2$
    \end{jieda}


    \item 设抛物线 ${y}^{2} = {8x}$ 的准线与 $x$ 轴交于点 $Q$ . 若过点 $Q$ 的直线 $l$ 与此抛物线有公共点, 则直线 $l$ 的斜率取值范围是\blkda{$[-1, 1]$}.
    \begin{jieda}
        易知相切的时候为边界情况，设切点$\left(\dfrac{y_0^2}{8}, y_0\right)$，则切线为$y_0y = 4\left(x+\dfrac{y_0^2}{8}\right)$，即$y = \dfrac{4}{y_0}x + \dfrac{y_0}{2}$，代入$(-2, 0)$，可得$y_0^2 = 16$，故切线斜率为$\pm 1$，故直线 $l$ 的斜率取值范围是$[-1, 1]$
    \end{jieda}


    \item 已知抛物线 ${y}^{2} = {4x}$ ,焦点为 $F$ ,顶点为原点 $O$ ,点 $P$ 在抛物线上移动, $Q$ 是 ${OP}$ 的中点, $M$ 是 ${FQ}$ 的中点,则点 $M$ 的轨迹方程是\blkda[3]{$y^2 = x-\dfrac{1}{2}\left(x > \dfrac{1}{2}\right)$}.
    \begin{jieda}
        设$Q(x, y)$，则满足$4{y}^{2} = 8x$，即$Q$点轨迹为$y^2 = 2x(x > 0)$.

        设$M(x, y)$，则满足$4y^2 = 2(2x-1)$，即$M$的轨迹为$y^2 = x-\dfrac{1}{2} \left(x > \dfrac{1}{2}\right)$
    \end{jieda}


    \item 若抛物线 ${y}^{2} = {2px}\left( {p > 0}\right)$ 的弦 ${PQ}$ 的中点为 $M\left( {{x}_{0},{y}_{0}}\right) \left( {{y}_{0} \neq  0}\right)$ ,则直线 ${PQ}$ 的斜率为\blkda{$\dfrac{p}{y_0}$}.
    \begin{jieda}
        设$P(x_1, y_1)$，$Q(x_2, y_2)$，则$\left\{\begin{array}{l}
             y_1^2 = 2px_1  \\
             y_2^2 = 2px_2 
        \end{array}\right.$，两式作差可得$(y_1+y_2)(y_1-y_2) = 2p(x_1-x_2)$，故$k_{PQ} = \dfrac{y_1-y_2}{x_1-x_2} = \dfrac{2p}{y_1+y_2} = \dfrac{p}{y_0}$
    \end{jieda}


    \item 若过抛物线 ${y}^{2} = {4x}$ 的焦点的弦被焦点分为长为 $m$ 和 $n$ 的两部分,则 $m$ 与 $n$ 的关系为\dotfill{(\kh{C})}
    \xx{$m + n = 4$}{${mn} = 4$}{$m + n = {mn}$}{$m + n = {2mn}$}
    \begin{jieda}
        $m = x_1+\dfrac{p}{2}$，$n = x_2+\dfrac{p}{2}$

        $m+n = x_1+x_2+p$，$mn = x_1x_2+\dfrac{p}{2}\left(x_1+x_2\right)+\dfrac{p^2}{4} = \dfrac{y_1^2y_2^2}{4p^2}+\dfrac{p}{2}\left(x_1+x_2\right) + \dfrac{p^2}{4} = \dfrac{p^2}{2}+\dfrac{p}{2}\left(x_1+x_2\right)$

        $p = 2$时有$m+n = x_1+x_2+2$，$mn = 2+x_1+x_2$
    \end{jieda}
    

    \item 如图,过抛物线 ${y}^{2} = {2px}\left( {p > 0}\right)$ 的焦点 ${F}$ 作直线交抛物线于 ${A}\text{ 、 }{B}$ 两点, ${M}$ 为准线 $I$ 上任意一点,记 $\angle {{AMF}} = \alpha ,\angle {{BMF}} =$  $\beta ,\angle {{MFO}} = \theta$ ,若 ${{AM}}\bot {{BM}}$ ,则 $\left| {\alpha  - \beta }\right|$ 与 $\theta$ 的大小关系为\dotfill{(\kh{B})}
    \xx{$\left| {a - \beta }\right|  > \theta$}{$\left| {a - \beta }\right|  = \theta$}{$\left| {a - \beta }\right|  < \theta$}{不能确定}
    
    \begin{center}
    \includegraphics[width=0.2\textwidth]{bodyxb/src/bo_d3qo4vref24c73d0hmv0_1_215_203_291_320_0.jpg}
    \end{center}

    \begin{jieda}
        当$M$位于$x$轴上方时，即当前图中所示情况，根据抛物线焦点弦的性质可知$MF \bot AB$
        
        做$BB_1 \bot l$交$l$于$B_1$，则$\triangle BB_1M \cong \triangle BFM$，故$\angle{BMB_1} = \beta$
        
        设$BM$交$x$轴于$N$，$l$与$x$轴交于$N'$，则一方面$\angle{MNN'} = \beta + \theta$，另一方面$\angle{MNN'} = \angle{MBB'} = \dfrac{\pi}{2} - \beta = \alpha$，因此有$\alpha - \beta = \theta$

        当$M$位于$x$轴下方时，则根据对称性可知有$\beta - \alpha = \theta$
    \end{jieda}


    \item 过点 $M\left( {2,4}\right)$ 作直线 $l$ ,与抛物线 ${y}^{2} = {8x}$ 只有一个公共点,这样的直线有\dotfill{(\kh{C})}
    \xx{0 条}{1 条}{2 条}{3 条}


    \item 抛物线 $y = a{x}^{2}\left( {a > 0}\right)$ 与直线 $y = {kx} + b$ 有两个交点,横坐标分别为 ${x}_{1}\text{ 、 }{x}_{2}$ ,而 ${x}_{3}$ 是直线与 $x$ 轴的交点的横坐标,则\dotfill{(\kh{D})}
    \xx{${x}_{3} = {x}_{1} + {x}_{2}$}{${x}_{3} = \dfrac{1}{{x}_{1}} + \dfrac{1}{{x}_{2}}$}{$x_{1}x_{3} = x_{2}x_{3} + x_{1}x_{2}$}{${x}_{1}{x}_{2} = {x}_{2}{x}_{3} + {x}_{1}{x}_{3}$}
    \begin{jieda}
        考虑$k>0, b>0$的情况显然$A$错，当$b$趋近于$0$时，$x_1, x_2, x_3$均趋近于$0$，显然$B$错，又因为$x_1, x_2$应当轮换对称，所以排除$C$，选$D$，下证$D$选项正确

        联立直线和抛物线，消去$y$可得$ax^2 = kx+b$，即$ax^2 - kx - b = 0$，根据韦达定理有$\left\{\begin{array}{l}
             x_1+x_2 = \dfrac{k}{a} \\
             x_1x_2 = -\dfrac{b}{a}
        \end{array}\right.$，故${x}_{1}{x}_{2} = {x}_{2}{x}_{3} + {x}_{1}{x}_{3}$即$-\dfrac{b}{a} = \dfrac{k}{a} \times \left(-\dfrac{b}{k}\right)$，成立.
    \end{jieda}

    \item 已知直线 $y = k\left( {x + 2}\right) \left( {k > 0}\right)$ 与抛物线 $C : {y}^{2} = {8x}$ 相交于 $A\text{ 、 }B$ 两点, $F$ 为 $C$ 的焦点,若 $\left| {FA}\right|  = 2\left| {FB}\right|$ ,则 $k =$\dotfill{(\kh{D})}
    \xx{$\dfrac{1}{3}$}{$\dfrac{\sqrt{2}}{3}$}{$\dfrac{2}{3}$}{$\dfrac{2\sqrt{2}}{3}$}
    \begin{jieda}
        $\left| {FA}\right|  = 2\left| {FB}\right| \Longleftrightarrow \left\{\begin{array}{l}
             x_A + 2 = 2(x_B+2) \\
             y_A = 2y_B
        \end{array}\right. \Longleftrightarrow \left\{\begin{array}{l}
             x_A = 2x_B + 2 \\
             y_A = 2y_B
        \end{array}\right. \Longleftrightarrow \left\{\begin{array}{l}
             \dfrac{y_A^2}{8} = \dfrac{y_B^2}{4} + 2 \\
             y_A = 2y_B
        \end{array}\right.$，故$y_B = 2\sqrt{2}$，$x_B = 1$，故$k = \dfrac{2\sqrt{2}}{3}$
    \end{jieda}


    \item 已知探照灯反射镜的轴截面是抛物线的一部分. 若某探照灯的灯口直径是 ${60}{{cm}}$ ,灯深是 ${40}{{cm}}$ ,则相应的抛物线的焦点到顶点的距离等于\dotfill{(\kh{D})}
    \xx{${40}{{cm}}$}{${20}{{cm}}$}{$\dfrac{45}{4}{{cm}}$}{$\dfrac{45}{8}{{cm}}$}

    \item 定长为 3 的线段 ${AB}$ 的两端点在抛物线 ${y}^{2} = x$ 上移动,记线段 ${AB}$ 的中点为 $M$ ,求点 $M$ 到 $y$ 轴的最短距离,并求此时直线 ${AB}$ 的方程.
    \begin{jieda}
        由抛物线方程 $y^2 = x$ 可知,其焦点 $F\left( {\dfrac{1}{4},0}\right)$ ,准线方程： $x =  - \dfrac{1}{4}$ ,
        从而 $\left| {AF}\right|  + \left| {BF}\right|  \geq  \left| {AB}\right|  = 3$ ,当且仅当 $A,B,F$ 三点共线时,不等式取等号
        
        设线段 ${AB}$ 的中点为 $M$ 到 $y$ 轴的距离为 $d$，由抛物线定义和梯形中位线性质可知, $2\left( {d + \dfrac{1}{4}}\right)  = \left| {AF}\right|  + \left| {BF}\right|  \geq  3$ ,即 $d \geq  \dfrac{5}{4}$ ,
        从而点 $M$ 到 $y$ 轴的最短距离为 $\dfrac{5}{4}$，此时 $A,B,F$ 三点共线
        
        故有$3 = |AB| = \dfrac{1}{sin^2\alpha}$，解得$sin\alpha = \dfrac{\sqrt{3}}{3}$，因此直线$AB$的斜率为$\pm \dfrac{\sqrt{2}}{2}$，故$AB: y = \pm \dfrac{\sqrt{2}}{2}\left(x-\dfrac{1}{2}\right)$，即${4x} - 4\sqrt{2}y - 1 = 0$ 或 ${4x} + 4\sqrt{2}y - 1 = 0$
    \end{jieda}


    \item 在平面直角坐标系 ${xOy}$ 中,直线 $l$ 与抛物线 ${y}^{2} = {2x}$ 相交于 $A\text{ 、 }B$ 两点.
    \begin{enumerate}
        \item 求证:“如果直线 $l$ 过点 $T\left( {3,0}\right)$ ,那么 $\vv{OA} \cdot  \vv{OB} = 3$ ” 是真命题
        \item 请写出第(1)题中命题的逆命题,判断它是真命题还是假命题,并说明理由
    \end{enumerate}
    \begin{jieda}
        \begin{enumerate}
            \item 证明：
            
            设过点 $T\left( {3,0}\right)$ 的直线 $l$ 交抛物线 $y^2 = {2x}$ 于点 $A\left( {{x}_{1},{y}_{1}}\right) ,B\left( {{x}_{2},{y}_{2}}\right)$ 
            
            设直线 $l$ 的方程为 $x = my+3$，联立直线和抛物线方程$\left\{  \begin{array}{l} {y}^{2} = {2x} \\  
            x = my+3 \end{array}\right.$，消去$x$可得${y}^{2} - 2my - 6 = 0$，$\Delta > 0$，根据韦达定理有${y}_{1}{y}_{2} =  - 6$

            
            所以 $\vv{OA} \cdot  \vv{OB} = {x}_{1}{x}_{2} + {y}_{1}{y}_{2} = \dfrac{y_1^2y_2^2}{4} + {y}_{1}{y}_{2} = \dfrac{36}{4} - 6 = 3$
            \item 逆命题是：“设直线与抛物线 ${y^2} = {2x}$ 相交于 $A$ 、 $B$ 两点,如果 $\vv{OA} \cdot  \vv{OB} = 3$ ,那么该直线过点 $y = \dfrac{2}{3}\left( {x + 1}\right)$ ”,该命题是假命题
            
            反例：取抛物线上的点 $A\left( {2,2}\right) ,B\left( {\dfrac{1}{2},1}\right)$ ,此时 $\vv{OA} \cdot  \vv{OB} = 3$ ,直线 ${AB}$ 的方程为 $y = \dfrac{2}{3}\left( {x + 1}\right)$ ,而 $T\left( {3,0}\right)$ 不在直线 ${AB}$ 上.
        \end{enumerate}
    \end{jieda}

    \item 过直角坐标平面 ${xOy}$ 中的抛物线 ${y}^{2} = {2px}\left( {p > 0}\right)$ 的焦点 $F$ 作一条倾斜角为 $\dfrac{\pi }{4}$ 的直线与抛物线相交于 ${A}$ 、 ${B}$ 两点.
    
    \begin{center}
    \includegraphics[width=0.2\textwidth]{bodyxb/src/bo_d3qo4vref24c73d0hmv0_2_231_976_343_194_0.jpg}
    \end{center}
    \begin{enumerate}
        \item 用 $p$ 表示 ${A}$ 、 ${B}$ 之间的距离
        \item 证明: $\angle {AOB}$ 的大小是与 $p$ 无关的定值,并求出这个值
    \end{enumerate}
    \begin{jieda}
        \begin{enumerate}
            \item 焦点 $F$ 的坐标为 $\left( {\dfrac{p}{2},0}\right)$，直线 ${AB}$ 的方程为$x = y+\dfrac{p}{2}$.
            
            联立方程组 $\left\{  \begin{array}{l} 
                x = y+\dfrac{p}{2}, \\  
                {y}^{2} = {2px}. \end{array}\right.$ 
            消去$x$得 ${y}^{2} = 2p\left(y+\dfrac{p}{2}\right)$，整理可得$y^2 - 2py - p^2 = 0$，由直线过焦点易知必有$A, B$两个交点
            
            ${y}_{1} + {y}_{2} = 2p$，因为 $x = y + \dfrac{p}{2}$ ,所以$\left| {AB}\right|  = {x}_{1} + {x}_{2} + p = (y_1+y_2) + {2p} = 4p$.
            \item $cos\angle{AOB} = \dfrac{x_1x_2+y_1y_2}{\sqrt{x_1^2+y_1^2} \cdot \sqrt{x_2^2+y_2^2}} = \dfrac{\dfrac{y_1^2y_2^2}{4p^2} + y_1y_2}{\sqrt{\dfrac{y_1^4}{4p^2}+y_1^2} \cdot \sqrt{\dfrac{y_2^4}{4p^2}+y_2^2}} = \dfrac{-\dfrac{3}{4}p^2}{p^2 \cdot \sqrt{\dfrac{y_1^2}{4p^2}+1} \cdot \sqrt{\dfrac{y_2^2}{4p^2}+1}}  = \dfrac{-\dfrac{3}{4}}{\sqrt{\dfrac{y_1^2y_2^2}{16p^4} + \dfrac{y_1^2+y_2^2}{4p^2}+1}} = \dfrac{-\dfrac{3}{4}}{\sqrt{\dfrac{1}{16}+\dfrac{4p^2+2p^2}{4p^2}+1}} = -\dfrac{3\sqrt{41}}{41}$

            故$\angle{AOB} = \pi - arccos\left(\dfrac{3\sqrt{41}}{41}\right)$
        \end{enumerate}
    \end{jieda}
    
    \item 设 $A,B$ 是抛物线 ${y}^{2} = {2px}\left( {p > 0}\right)$ 上的两点,且 ${OA} \bot  {OB}$ .
    \begin{enumerate}
        \item 求证：$A$ 、 $B$ 两点的横坐标之积,纵坐标之积为定值
        \item 求证：直线 ${AB}$ 过定点
        \item 求弦 ${AB}$ 中点 $P$ 的轨迹方程
        \item 求点 $O$ 在 ${AB}$ 上的射影 $M$ 的轨迹方程
    \end{enumerate}
    \begin{jieda}
        \begin{enumerate}
            \item 设 $A\left( {{x}_{1},{y}_{1}}\right) ,B\left( {{x}_{2},{y}_{2}}\right)$ ,则 $\left\{\begin{array}{l}
                 {y}_{1}^{2} = {2p}{x}_{1}  \\
                 {y}_{2}^{2} = {2p}{x}_{2}
            \end{array}\right.$ . 所以 ${y}_{1}^{2}{y}_{2}^{2} = 4{p}^{2}{x}_{1}{x}_{2}$ .
            
            又因为 ${OA} \bot  {OB}$ ,所以 ${x}_{1}{x}_{2} + {y}_{1}{y}_{2} = 0$ . 由此,可解得 ${x}_{1}{x}_{2} = 4{p}^{2},{y}_{1}{y}_{2} =  - 4{p}^{2}$ .
            \item 易知直线$AB$的斜率不为0，$\dfrac{{x}_{2} - {x}_{1}}{{y}_{2} - {y}_{1}} = \dfrac{\dfrac{{y}_{2}^{2}}{2p} - \dfrac{{y}_{1}^{2}}{2p}}{{y}_{2} - {y}_{1}} = \dfrac{{y}_{1} + {y}_{2}}{2p}$ .
            
            此时直线 ${AB}$ 的方程为 $x - \dfrac{{y}_{1}^{2}}{2p} = \dfrac{{y}_{1} + {y}_{2}}{2p}\left( y - y_1\right)$ . 即 $y\left( {{y}_{1} + {y}_{2}}\right)  - {y}_{1}{y}_{2} = {2px}$ . 

            由上一问可知${y}_{1}{y}_{2} =  - 4{p}^{2}$，故直线方程为$y\left( {{y}_{1} + {y}_{2}}\right)  +4p^2 = {2px}$，即$x = \dfrac{y_1+y_2}{2p} y + 2p$，过定点$(2p, 0)$
            \item 设中点$P(x_0, y_0)$，则直线$AB: x = \dfrac{y_0}{p} y + 2p$，故$x_0 = \dfrac{y_0^2}{p}+2p$，因此$P$的轨迹为$x = \dfrac{y^2}{p}+2p$
            \item $OM$的方程为$y = -\dfrac{y_0}{p}x$，
            
            \begin{enumerate}
                \item 方法1

                $M$满足$OM$和$AB$直线方程，消去$y_0$即可得到$M$的轨迹方程.

                $\left\{\begin{array}{l}
                 x = \dfrac{y_0}{p} y + 2p  \\
                 y = -\dfrac{y_0}{p}x 
                \end{array}\right. \Longleftrightarrow \left\{\begin{array}{l}
                 y_0 = \dfrac{p(x-2p)}{y}  \\
                 y_0 = -\dfrac{py}{x} 
                \end{array}\right.$
                消去$y_0$得到$-\dfrac{py}{x} = \dfrac{p(x-2p)}{y}$，故$M$ 的轨迹方程为$x^2 - 2px + y^2 = 0 (x \neq 0)$
                \item 方法2

                联立$OM$和$AB$：$\left\{\begin{array}{l}
                 x = \dfrac{y_0}{p} y + 2p  \\
                 y = -\dfrac{y_0}{p}x 
                \end{array}\right.$，消去$y$可得$x = -\dfrac{y_0^2}{p^2}x + 2p$，故$x_M = \dfrac{2p^3}{p^2+y_0^2}$，将$y_0 = \dfrac{py_M}{-x_M}$代入有$x_M = \dfrac{2p^3}{p^2+\dfrac{p^2y_M^2}{x_M^2}} = \dfrac{2p}{1+\dfrac{y_M^2}{x_M^2}}$，去分母有$x_M + \dfrac{y_M^2}{x_M} = 2p$，进一步变形有$x_M^2 + y_M^2 = 2px_M (x_M \neq 0)$
    
                故$M$ 的轨迹方程为$x^2 - 2px + y^2 = 0 (x \neq 0)$
            \end{enumerate}
        \end{enumerate}
    \end{jieda}
\end{enumerate}
\section*{求轨迹的方程}
\subsection*{求轨迹的方程 (1)}

A 组：
\begin{enumerate}
    \item “曲线 $C$ 上的点的坐标都是方程 $f\left( {x,y}\right)  = 0$ 的解”是“方程 $f\left( {x,y}\right)  = 0$ 是曲线 $C$ 的方程”的\blkda[2]{必要不充分}条件.

    \item 下列命题中真命题的序号是\blkda{\ding{173}}.
    \begin{dingautolist}{172}
        \item 方程 ${\left( x - y\right) }^{2} + {\left( xy - 1\right) }^{2} = 0$ 的曲线是一条直线和一条双曲线
        \item 两坐标轴所成的角在第二、第四象限的平分线方程是 $x + y = 0$
        \item 若 $A\left( {5,0}\right) ,B\left( {0, - 5}\right)$ ,则线段 ${AB}$ 的方程是 $x - y - 5 = 0$
    \end{dingautolist}

    
    \item 若命题 “曲线 $C$ 上的点的坐标都是方程 $f\left( {x,y}\right)  = 0$ 的解” 是正确的,则以下命题中正确的是\dotfill{(\kh{D})}
    \xx{方程 $f\left( {x,y}\right)  = 0$ 是曲线 $C$ 的方程}{坐标满足方程 $f\left( {x,y}\right)  = 0$ 的点都在曲线 $C$ 上}{曲线 $C$ 是方程 $f\left( {x,y}\right)  = 0$ 的曲线}{方程 $f\left( {x,y}\right)  = 0$ 的曲线不一定是曲线 $C$}
    
    \item 已知 $A\left( {-1,0}\right)$ , $B\left( {1,0}\right)$ 分别为 $\bigtriangleup  {ABC}$ 的两个顶点,动点 $C$ 满足 $\vv{CA} \cdot  \vv{CB} = 0$ ,则动点 $C$ 的轨迹方程为\blkda[3]{$x^2+y^2 = 1(y \neq 0)$}.

    \item 下面各组方程是否表示同一曲线,为什么?
    \begin{dingautolist}{172}
        \item ${x}^{2} + {y}^{2} = 0$ 与 ${xy} = 0$
        \item $\left| y\right|  = \left| x\right|$ 与 ${y}^{2} = {x}^{2}$
        \item $y = \lg {x}^{2}$ 与 $y = 2\lg x$
        \item $y = \sin \left( {\arcsin x}\right)$ 与 $y = x$
    \end{dingautolist}
    \begin{jieda}
        \begin{dingautolist}{172}
            \item 不是，前者是坐标原点，后者是两坐标轴
            \item 是
            \item 不是，定义域不同，前者两支，后者一支
            \item 不是，定义域不同
        \end{dingautolist}
    \end{jieda}
\end{enumerate}


B 组：
\begin{enumerate}
    \item 已知点 $A\left( {-1,0}\right)$ , $B\left( {2,0}\right)$ ,动点 $M$ 满足 $2\left| {MA}\right|  = \left| {MB}\right|$ . 则动点 $M$ 的轨迹方程为\blkda[3]{${\left( x + 2\right) }^{2} + {y}^{2} = 4$}
    
    \item 已知命题 $A$ : 点 $P$ 在曲线 $y = \left| x\right|$ 上,命题 $B$ : 点 $P$ 到两坐标轴的距离相等,命题 $A$ 是命题 $B$ 的\blkda[2]{充分不必要}条件.

    \item 已知 $A\text{ 、 }B$ 两点的坐标分别是 $\left( {-3,5}\right) \text{ 、 }\left( {3,5}\right)$ . 若动点 $M$ 满足 $\left| {MA}\right|  + \left| {MB}\right|  = 6$ ,则点 $M$ 的轨迹方程是\blkda[3]{$y = 5, x \in [-3, 3]$}.

    \item 若 $x$ 轴上方的一动点 $P$ 到 $x$ 轴的距离与它到原点的距离的比是 $\dfrac{\sqrt{7}}{4}$ ,则点 $P$ 的轨迹方程是\blkda[3]{$7x^2=9y^2, y > 0$}.
    \begin{jieda}
        设$P(x, y), y > 0$，则根据题意有$\dfrac{y}{\sqrt{x^2+y^2}} = \dfrac{\sqrt{7}}{4}$，两边平方整理得$7x^2=9y^2$
    \end{jieda}

    \item 方程 $\left( {\left| x\right|  - \left| y\right|  - 1}\right) \left( {{x}^{2} - 9}\right)  = 0$ 表示的曲线所围成的封闭图形的面积 $S =$ \blkda{8}.
    \begin{jieda}
        $\left( {\left| x\right|  - \left| y\right|  - 1}\right) \left( {{x}^{2} - 9}\right)  = 0$ 表示$|x| - |y| = 1$和$x^2 = 9$的并集

        如图所示，计算可得面积为$8$.
        \begin{center}
            \includegraphics[width=0.5\linewidth]{bodyxb//src/guijifangcheng1B5.png}
        \end{center}
    \end{jieda}

    \item 曲线 ${y}^{2} =  - {4x}$ 关于直线 $l : x - y - 1 = 0$ 对称的曲线方程是\blkda[4]{$x^2-2x+4y+5=0$}.
    \begin{jieda}
        设对称曲线上的点为$(x, y)$，则关于直线 $l : x - y - 1 = 0$ 对称的点为$(y+1, x-1)$，因此$(x-1)^2 = -4(y+1)$，整理得$x^2-2x+4y+5=0$
    \end{jieda}

    \item 已知一条曲线在 $x$ 轴上方,这条曲线上的每一点到点 $A\left( {0,2}\right)$ 的距离减去它到 $x$ 轴的距离之差都是 2 ,则这条曲线的方程是\blkda[3]{$x^2= 8y, y \neq 0$}.

    \item 若曲线 $C$ 的方程不是 $F\left( {x,y}\right)  = 0$ ,则曲线 $C$ 上\dotfill{(\kh{D})}
    \xx{任意一点的坐标都不满足方程 $F\left( {x,y}\right)  = 0$}{存在无数个点,它们的坐标都不满足方程 $F\left( {x,y}\right)  = 0$}{至少存在一点,其坐标不满足方程 $F\left( {x,y}\right)  = 0$}{任意一点的坐标有可能都满足方程 $F\left( {x,y}\right)  = 0$}

    \item 方程 $x{y}^{2} - {x}^{2}y = {8x}$ 的曲线$\overline{\quad}$对称\dotfill{(\kh{C})}
    \xx{关于 $y$ 轴}{关于 $x + y = 0$}{关于原点}{关于 $x - y = 0$}
    \begin{jieda}
        将$(-x, -y)$代换$(x, y)$方程依然成立.
    \end{jieda}
    
    \item 方程 ${x}^{2} - {xy} + {y}^{2} + x - y - 1 = 0$ 的曲线具有\dotfill{(\kh{A})}
    \xx{对称轴 $x + y = 0$}{对称轴 $x - y = 0$}{对称轴 $y = 0$}{对称中心 $\left( {0,0}\right)$}
    \begin{jieda}
        将$(-y, -x)$代换$(x, y)$方程依然成立.
    \end{jieda}
    
    \item 过抛物线 $y = {x}^{2}$ 的顶点 $O$ 作互相垂直的两条直线交抛物线于点 $A$ 和 $B$ .
    \begin{enumerate}
        \item 求 ${AB}$ 的中点 $P$ 的轨迹方程
        \item 求顶点 $O$ 在直线 ${AB}$ 上的射影 $M$ 的轨迹方程
    \end{enumerate}
    \begin{jieda}
        \begin{enumerate}
            \item 设$A(x_1, y_2), B(x_2, y_2)$，则根据$OA \bot OB$可知$x_1x_2+y_1y_2 = 0$，即$x_1x_2+x_1^2x_2^2 = 0$，故$x_1x_2 = -1$

            $\left\{\begin{array}{l}
                 y_1 = x_1^2 \\
                 y_2 = x_2^2
            \end{array}\right.$，故$y_1+y_2 = x_1^2+x_2^2 = (x_1+x_2)^2-2x_1x_2 = (x_1+x_2)^2+2$，设$P(x, y)$，则$2y_P = (2x_P)^2+2$，即$P$的轨迹是$y = 2x^2+1$
            \item 设$AB: y = kx+b$，与抛物线方程联立后根据韦达定理有$x_1x_2 = -b$，又因为$x_1x_2 = -1$，所以$AB$恒过$(0, 1)$，因此$AB: y = kx+1$，$OM: y = -\dfrac{1}{k}x$，故$k = \dfrac{y_M - 1}{x_M}$且$k = -\dfrac{x_M}{y_M}$，因此$M$的轨迹是$\dfrac{y - 1}{x} = -\dfrac{x}{y}$，即$x^2+y^2-y = 0(y \neq 0)$
        \end{enumerate}
    \end{jieda}
    
    \item 已知 $\bigtriangleup  {ABC}$ 的两个顶点坐标分别是 $B\left( {-5,0}\right)$ 和 $C\left( {5,0}\right)$ ,且 $\angle {ABC} + \angle {ACB} = {135}^{ \circ  }$ ,顶点 $A$ 在 $x$ 轴上方,求顶点 $A$ 的轨迹方程.
    \begin{jieda}
        根据三角形内角和为$\ang{180}$可知$\angle{BAC} = \ang{45}$，因此$A$的轨迹是一段以$(0, 5)$为圆心，$5\sqrt{2}$为半径的圆弧，即$x^2+(y-5)^2 = 50(y>0)$
    \end{jieda}

    \item 已知曲线 ${C} : \dfrac{{x}^{2}}{25} + \dfrac{{y}^{2}}{16} = 1$ 及点 $F\left( {-3,0}\right)$ ,若点 $Q$ 为曲线 ${C}$ 上一点,且点 $P$ 满足 $\vv{FQ} = 2\vv{QP}$ ,求点 $P$ 的轨迹方程.
    \begin{jieda}
        设$P(x, y)$，则根据$\vv{FQ} = 2\vv{QP}$可知$Q\left(\dfrac{2x-3}{3}, \dfrac{2y}{3}\right)$，又因为点 $Q$ 为曲线 ${C}$ 上一点，所以$\dfrac{\left(\dfrac{2x-3}{3}\right)^{2}}{25} + \dfrac{\left(\dfrac{2y}{3}\right)^{2}}{16} = 1$，整理可得$\dfrac{{\left( 2x - 3\right) }^{2}}{225} + \dfrac{{y}^{2}}{36} = 1$
    \end{jieda}

    \item 已知定点 $A\left( {2,0}\right)$ 和曲线 ${C} : {x}^{2} + {y}^{2} = 1$ (表示单位圆) 上的动点 $Q$ , $\angle {AOQ}$ 的角平分线交 ${AQ}$ 于点 $P$ ,求点 $P$ 的轨迹方程.
    \begin{jieda}
        根据角分线定理可知$\dfrac{|AP|}{|QP|} = \dfrac{|AO|}{|QO|} = 2$，设$P(x, y)$，则$Q\left(\dfrac{3}{2}x-1, \dfrac{3y}{2}\right)$，又因为点 $Q$ 为曲线 ${C}$ 上一点，所以$\left(\dfrac{3}{2}x-1\right)^2+\left(\dfrac{3y}{2}\right)^2=1$，即$\left(x-\dfrac{2}{3}\right)^2+y^2 = \dfrac{4}{9}$

        当$Q$为$(1, 0)$时不符合情景，因此点 $P$ 的轨迹方程为$\left(x-\dfrac{2}{3}\right)^2+y^2 = \dfrac{4}{9} \left(x \neq \dfrac{4}{3}\right)$
    \end{jieda}
\end{enumerate}

C组：
\begin{enumerate}
    \item 古希腊数学家阿波罗尼奥斯在研究圆锥曲线时发现了椭圆的光学性质:从椭圆的一个焦点射出的光线, 经椭圆反射, 其反射光线必经过椭圆的另一焦点. 设椭圆方程 $\dfrac{{x}^{2}}{{a}^{2}} + \dfrac{{y}^{2}}{{b}^{2}} = 1\left( {a > b > 0}\right) ,{F}_{1}$ 、 ${F}_{2}$ 为其左、右焦点,若从右焦点 ${F}_{2}$ 发出的光线经椭圆上的点 ${A}$ 和点 $B$ 反射后,满足 ${AB} \bot  {AD},\cos \angle {ABC} = \dfrac{4}{5}$ ,则该椭圆的离心率为\blkda{$\dfrac{\sqrt{2}}{2}$}.
    \begin{center}
    \includegraphics[width=0.3\textwidth]{bodyxb/src/bo_d43fvmc601uc738mn8og_3_214_355_365_318_0.jpg}
    \end{center}
    
    \begin{jieda}
        由题意, 可作图如下:
        
        \begin{center}
        \includegraphics[width=0.3\textwidth]{bodyxb/src/bo_d43g3ak601uc738mn9ng_0_241_1253_362_315_0.jpg}
        \end{center}
        
        则 $\cos \angle {AB}{F}_{1} = \dfrac{4}{5} = \dfrac{\left| AB\right| }{\left| B{F}_{1}\right| },\sin \angle {AB}{F}_{1} = \sqrt{1 - {\cos }^{2}\angle {AB}{F}_{1}} = \dfrac{3}{5} = \dfrac{\left| A{F}_{1}\right| }{\left| B{F}_{1}\right| }$ ,即 $\left| {AB}\right|  : \left| {A{F}_{1}}\right|  : \left| {B{F}_{1}}\right|  = 4 : 3 : 5$ , 可设 $\left| {AB}\right|  = {4k}$ , $\left| {A{F}_{1}}\right|  = {3k}$ , $\left| {B{F}_{1}}\right|  = {5k}$ ,
        
        由 $\left| {AB}\right|  + {\left| AF\right| }_{1} + \left| {B{F}_{1}}\right|  = \left| {A{F}_{2}}\right|  + \left| {B{F}_{2}}\right|  + \left| {A{F}_{1}}\right|  + \left| {B{F}_{1}}\right|  = {4a}$ ,则 ${4k} + {3k} + {5k} = {4a}$ ,即 ${3k} = a$ ,
        
        $\left| {A{F}_{2}}\right|  = {2a} - \left| {A{F}_{1}}\right|  = {3k}$ ,在 Rt $\bigtriangleup A{F}_{1}{F}_{2}$ 中, $\left| {{F}_{1}{F}_{2}}\right|  = \sqrt{{\left| A{F}_{1}\right| }^{2} + {\left| A{F}_{2}\right| }^{2}} = 3\sqrt{2}k = {2c}$ ,
        
        则 $e = \dfrac{2c}{2a} = \dfrac{3\sqrt{2}k}{6k} = \dfrac{\sqrt{2}}{2}$ .
    \end{jieda}
    
    \item 已知双曲线 $G$ 的方程为 $\dfrac{{x}^{2}}{16} - \dfrac{{y}^{2}}{9} = 1$ ,其左、右焦点分别是 ${F}_{1}$ , ${F}_{2}$ ,点 $P$ 坐标为 $\left( {4,2}\right)$ , 双曲线 $G$ 上的 $Q\left( {{x}_{0},{y}_{0}}\right) \left( {{x}_{0} > 0,{y}_{0} > 0}\right)$ 满足 $\dfrac{\vv{Q{F}_{1}} \cdot  \vv{P{F}_{1}}}{\left| \vv{Q{F}_{1}}\right| } = \dfrac{\vv{{F}_{2}{F}_{1}} \cdot  \vv{P{F}_{1}}}{\left| \vv{{F}_{2}{F}_{1}}\right| }$ ,则 $\left( {{S}_{\bigtriangleup {F}_{1}{PQ}} - {S}_{\bigtriangleup {F}_{2}{PQ}}}\right) \left( {\left| {Q{F}_{2}}\right|  - \left| {Q{F}_{1}}\right| }\right)  =$ \blkda{$-64$}.
    \begin{jieda}
        如图,设 $\bigtriangleup  {Q{F}_{1}{F}_{2}}$ 的内切圆与三边分别相切于 $D,E,G$ ,
        
        可得 $\left| {QD}\right|  = \left| {QG}\right| ,\left| {{F}_{1}D}\right|  = \left| {{F}_{1}E}\right| ,\left| {{F}_{2}E}\right|  = \left| {{F}_{2}G}\right|$ ,
        
        又由双曲线定义可得 $\left| {Q{F}_{1}}\right|  - \left| {Q{F}_{2}}\right|  = {2a} = 8$ ,
        
        则 $\left| {QD}\right|  + \left| {DF}\right|  - \left( {\left| {QG}\right|  + \left| {G{F}_{2}}\right| }\right)  = \left| {D{F}_{1}}\right|  - \left| {G{F}_{2}}\right|  = E{F}_{1} + E{F}_{2} = 2$ ,
        
        又 $\left| {E{F}_{1}}\right|  + \left| {E{F}_{2}}\right|  = {2c}$ ,解得 $\left| {E{F}_{1}}\right|  = a + c$ ,
        
        则点 $E$ 横坐标为 $a$ ,即内切圆圆心横坐标为 $a$ .
        
        又 $\dfrac{\vv{Q{F}_{1}} \cdot  \vv{P{F}_{1}}}{\left| \vv{Q{F}_{1}}\right| } = \dfrac{\vv{{F}_{2}{F}_{1}} \cdot  \vv{P{F}_{1}}}{\left| \vv{{F}_{2}{F}_{1}}\right| }$ ,可得 $\dfrac{\left| \vv{Q{F}_{1}}\right|  \cdot  \left| \vv{P{F}_{1}}\right| \cos \angle P{F}_{1}Q}{\left| Q{F}_{1}\right| } = \dfrac{\left| \vv{{F}_{2}{F}_{1}}\right|  \cdot  \left| \vv{P{F}_{1}}\right| \cos \angle P{F}_{1}{F}_{2}}{\left| \vv{{F}_{2}{F}_{1}}\right| }$ ,
        
        化简得 $\cos \angle P{F}_{1}Q = \cos \angle P{F}_{1}{F}_{2}$ ,即 $\angle P{F}_{1}Q = \angle P{F}_{1}{F}_{2}$ ,
        
        即 $P{F}_{1}$ 是 $\angle Q{F}_{1}{F}_{2}$ 的平分线,由于 $P\left( {4,2}\right) ,a = 4$ ,
        
        所以点 $P$ 是 $\bigtriangleup  {Q}_{1}^{r}{P}_{2}^{r}$ 的内心,且半径 $r$ 为 2,
        
        则 ${S}_{\bigtriangleup {F}_{1}{PQ}} - {S}_{\bigtriangleup {F}_{2}{PQ}} = \dfrac{1}{2}r\left( {\left| {Q{F}_{1}}\right|  - \left| {Q{F}_{2}}\right| }\right)  = \dfrac{1}{2} \times  2 \times  8 = 8$ ,
        
        又 $\left| {Q{F}_{2}}\right|  - \left| {Q{F}_{1}}\right|  =  - {2a} =  - 8$ ,
        
        所以 $\left( {{S}_{\vee {F}_{1}{PQ}} - {S}_{\vee {F}_{2}{PQ}}}\right) \left( {\left| {Q{F}_{2}}\right|  - \left| {Q{F}_{1}}\right| }\right)  = 8 \times  \left( {-8}\right)  =  - {64}$ .
        
        
        \begin{center}
        \includegraphics[width=0.4\textwidth]{bodyxb/src/bo_d43g3ak601uc738mn9ng_2_240_195_532_392_0.jpg}
        \end{center}
    \end{jieda}

\end{enumerate}

\clearpage
\subsection*{求轨迹的方程 (2)}
\begin{enumerate}
    \item 已知动点 $M\left( {x,y}\right)$ 满足 $\sqrt{{\left( x + 2\right) }^{2} + {y}^{2}} - \sqrt{{\left( x - 2\right) }^{2} + {y}^{2}} = 4$ ,则动点 $M$ 的轨迹方程是\blkda[3]{$y = 0\left( {x \geq  2}\right)$}.
    
    \item 在平面直角坐标系 ${xOy}$ 中,动点 $P$ 关于 $x$ 轴对称的点为 $Q$ ,且 $\vv{OP} \cdot  \vv{OQ} = 2$ ,则点 $P$ 的轨迹方程为\blkda[3]{${x}^{2} - {y}^{2} = 2$}.
    
    \item 数学中有很多形状优美,寓意美好的曲线,曲线 $C : {x}^{2} + {y}^{2} - 2\left| x\right|  - 2\left| y\right|  = 0$ 就是其中之一,则曲线 $C$ 所围成的封闭图形的面积是\blkda{$8 + 4\pi$}.
    \begin{jieda}
        方程 ${x}^{2} + {y}^{2} - 2\left| x\right|  - 2\left| y\right|  = 0$ ,对 $x,y$ 分类讨论

        \begin{dingautolist}{172}
            \item $x \geq  0,y \geq  0$ 时,化为: ${\left( x - 1\right) }^{2} + {\left( y - 1\right) }^{2} = 2$
            \item $x \geq  0,y \leq  0$ 时,化为: ${\left( x - 1\right) }^{2} + {\left( y + 1\right) }^{2} = 2$
            \item $x \leq  0,y \geq  0$ 时,化为: ${\left( x + 1\right) }^{2} + {\left( y - 1\right) }^{2} = 2$
            \item $x \leq  0,y \leq  0$ 时,化为: ${\left( x + 1\right) }^{2} + {\left( y + 1\right) }^{2} = 2$
        \end{dingautolist}
        \begin{center}
        \includegraphics[width=0.2\textwidth]{bodyxb/src/bo_d43hsi3ef24c73d4tekg_1_241_638_349_331_0.jpg}
        \end{center}

        所表示的曲线所围成的图形面积 $S = {\left( 2\sqrt{2}\right) }^{2} + {2\pi } \times  {\left( \sqrt{2}\right) }^{2} = 8 + {4\pi }$ .
    \end{jieda}
    
    \item 已知曲线 ${C}_{1} : \dfrac{{x}^{2}}{m} + {y}^{2} = 1$ 与曲线 ${C}_{2} : \left| y\right|  = {2x} + 2$ ,且曲线 ${C}_{1}$ 和 ${C}_{2}$ 恰有两个不同的交点,则实数 $m$ 的取值范围是\blkda[4]{$\left( {-\infty , - \dfrac{1}{4}}\right\rbrack   \cup  \left\{  \dfrac{3}{4}\right\}   \cup  \left( {1, + \infty }\right)$}.
    \begin{jieda}
        由题意,曲线 ${C}_{1}$ 过定点 $\left( {0, \pm  1}\right)$ ,曲线 ${C}_{2} : \left| y\right|  = {2x} + 2 \geq  0 \Longrightarrow  x \geq   - 1$ , 故 $\left| y\right|  = {2x} + 2$ 的图象为 $y = {2x} + 2\left( {x \geq   - 1}\right)$ 的图象及其关于 $x$ 轴对称的部分,如图所示:
        
        \begin{center}
        \includegraphics[width=0.2\textwidth]{bodyxb/src/bo_d43hsi3ef24c73d4tekg_4_238_201_349_338_0.jpg}
        \end{center}
        
        当 $m < 0$ 时,曲线 ${C}_{1}$ 为双曲线,则曲线 ${C}_{1}$ 和 ${C}_{2}$ 在 $y$ 轴左侧必有两个交点
        
        又渐近线为 $y =  \pm  \dfrac{1}{\sqrt{-m}}x$ ,故当 $2 \geq  \dfrac{1}{\sqrt{-m}}$ 即 $m \leq   - \dfrac{1}{4}$ 时,曲线 ${C}_{1}$ 和 ${C}_{2}$ 在 $y$ 轴右侧无交点，满足题意;
        
        当 $m > 0$ 时,曲线 ${C}_{1}$ 为圆或椭圆，联立 $\left\{  \begin{array}{l} \dfrac{{x}^{2}}{m} + {y}^{2} = 1, \\  y = {2x} + 2 \end{array}\right.$ 消 $y$ 得 $\left( {\dfrac{1}{m} + 4}\right) {x}^{2} + {8x} + 3 = 0$ ,
        
        当曲线 ${C}_{1}$ 与 $y = {2x} + 2$ 相切时, $\Delta  = {64} - {12}\left( {\dfrac{1}{{m}} + 4}\right)  = 0 \Longrightarrow  {m} = \dfrac{3}{4}$

        \begin{dingautolist}{172}
            \item 故当 ${m} = \dfrac{3}{4}$ 时,曲线 ${C}_{1}$ 和 ${C}_{2}$ 有有两个不同的交点
            \item 当 $0 < {m} < \dfrac{3}{4}$ 时,曲线 ${C}_{1}$ 和 ${C}_{2}$ 有零个不同的交点
            \item 当 $\dfrac{3}{4} < m < 1$ 时,曲线 ${C}_{1}$ 和 ${C}_{2}$ 有四个不同的交点
            \item 当 $m = 1$ 时,曲线 ${C}_{1}$ 和 ${C}_{2}$ 有三个不同的交点
            \item 当 $m > 1$ 时,曲线 ${C}_{1}$ 和 ${C}_{2}$ 有两个不同的交点
        \end{dingautolist}
        
        综上,曲线 ${C}_{1}$ 和 ${C}_{2}$ 恰有两个不同的交点,实数 $m$ 的取值范围为 $\left( {-\infty , - \dfrac{1}{4}}\right\rbrack   \cup  \left\{  \dfrac{3}{4}\right\}   \cup  \left( {1, + \infty }\right)$ .
    \end{jieda}
    
    \item 数学中有许多形状优美、寓意美好的曲线,曲线 $C : {x}^{2} + {y}^{2} = 4 + \left| x\right| y$ 就是其中之一. 曲线 $C$ 对应的图象如图所示
    
    \begin{center}
    \includegraphics[width=0.2\textwidth]{bodyxb/src/bo_d43fvo4601uc738mn8pg_0_240_1019_290_267_0.jpg}
    \end{center}
    下列结论:
    \begin{dingautolist}{172}
        \item 直线 ${AB}$ 的方程为: $x + y + 2 = 0$
        \item 曲线 $C$ 与圆 ${x}^{2} + {y}^{2} = 8$ 有 2 个交点
        \item 曲线 $C$ 阶围成的 “心形” 区域的面积大于 12
        \item 曲线 $C$ 恰好经过 4 个整点(即横、纵坐标均为整数的点).
    \end{dingautolist}
    其中正确的是:\blkda{\ding{173}\ding{174}}. (填写所有正确结论的编号)
    \begin{jieda}
        \begin{dingautolist}{172}
            \item 令 $x = 0$ ,则 $y =  \pm  2$ ; 令 $y = 0$ ,则 $x =  \pm  2$，所以点 $A\left( {2,0}\right) ,B\left( {0,2}\right)$ ,所以直线 ${AB}$ 的方程为: $\dfrac{x}{2} + \dfrac{y}{2} = 1$ 即 $x + y - 2 = 0$，故错误
            \item 由 $\left\{  \begin{array}{l} {x}^{2} + {y}^{2} = 4 + \left| x\right| y \\  {x}^{2} + {y}^{2} = 8 \end{array}\right.$ 可得 $\left\{  \begin{array}{l} x = 2 \\  y = 2 \end{array}\right.$ 或 $\left\{  \begin{array}{l} x =  - 2 \\  y = 2 \end{array}\right.$

            所以曲线 $C$ 与圆 ${x}^{2} + {y}^{2} = 8$ 有 2 个交点 $\left( {2,2}\right)$ , $\left( {-2,2}\right)$ ,故正确
            \item 在曲线上取点 $D\left( {2,2}\right)$ , $E\left( {-2,2}\right)$ , $F\left( {-2,0}\right)$ , $G\left( {0, - 2}\right)$ ,顺次连接各点,如图,
            
            \begin{center}
            \includegraphics[width=0.4\textwidth]{bodyxb/src/bo_d43hsi3ef24c73d4tekg_2_238_884_541_537_0.jpg}
            \end{center}
            
            则 ${S}_{ADEFG} = 2 \times  4 + \dfrac{1}{2} \times  4 \times  2 = {12}$ ,
            
            所以曲线 $C$ 所围成的“心形”区域的面积大于 12，故正确
            \item 曲线经过的整点有: $\left( {\pm 2,0}\right) ,\left( {0, \pm  2}\right) ,\left( {\pm 2,2}\right)$ ,有 6 个，故错误
        \end{dingautolist}
    \end{jieda}
    
    \item 曲线 $C$ 是平面内与两个定点 ${F}_{1}\left( {-1,0}\right)$ 和 ${F}_{2}\left( {1,0}\right)$ 距离之积等于 4 的点的轨迹,已知 $O$ 为坐标原点,点 $P$ 为曲线 $C$ 上的动点,关于曲线 $C$ 有以下 3 个结论:
    \begin{dingautolist}{172}
        \item 曲线 $C$ 是封闭曲线且不过原点
        \item 曲线 $C$ 关于 $x$ 轴对称,且关于 $y$ 轴对称
        \item 曲线 $C$ 上的动点 $P\left( {x,y}\right)$ 的坐标满足 $\left| x\right|  \leq  \sqrt{5}$ , $\left| y\right|  \leq  \sqrt{3}$
    \end{dingautolist}则所有正确的结论为\blkda{\ding{172}\ding{173}\ding{174}}.(填写所有正确结论的编号)
    \begin{jieda}
        $\because$ 动点 $M\left( {x,y}\right)$ 到平面两个定点 ${F}_{1}\left( {-1,0}\right)$ 和 ${F}_{2}\left( {1,0}\right)$ 距离之积等于 4,
        
        $\therefore \left\lbrack  {{\left( x + 1\right) }^{2} + {y}^{2}}\right\rbrack   \cdot  \left\lbrack  {{\left( x - 1\right) }^{2} + {y}^{2}}\right\rbrack   = {16}$ ,
        
        $\therefore \left( {{x}^{2} + {y}^{2} + 1 + {2x}}\right)  \cdot  \left( {{x}^{2} + {y}^{2} + 1 - {2x}}\right)  = {16}$ ,
        
        $\therefore {\left( {x}^{2} + {y}^{2} + 1\right) }^{2} = 4{x}^{2} + {16} \Longrightarrow  {y}^{2} = \sqrt{4{x}^{2} + {16}} - {x}^{2} - 1$
        \begin{dingautolist}{172}
            \item 令$t = \sqrt{4x^2+16},t \geq 4$，设$g(t) = t-\dfrac{t^2-16}{4}-1$，易知$g(t)$在$[4, +\infty)$严减，故根据同增异减可知$f(x) = \sqrt{4{x}^{2} + {16}} - {x}^{2} - 1$在$[0, +\infty)$严减，故曲线$C$在第一象限的部分是一段严格减函数的图象，根据\ding{173}中的对称性可知曲线为封闭图形，且不过原点，故正确
            \item 把方程中 $x$ 用 $- x$ 代换, $y$ 用 $- y$ 代换,方程不变,曲线 $C$ 关于 $x$ 轴对称,且关于 $y$ 轴对称，故正确
            \item 根据\ding{172}中分析的单调性可知$x = 0$时，$y$取得最值，$y = 0$时，$x$取得最值，代值验证可知命题正确.
        \end{dingautolist}
    \end{jieda}

    
    \item 方程 $\sqrt{{\left( x - 2\right) }^{2} + {y}^{2}} + \sqrt{{\left( x + 2\right) }^{2} + {y}^{2}} = {12}$ ,化简的结果是\dotfill{(\kh{B})}
    \xx{$\dfrac{{x}^{2}}{36} + \dfrac{{y}^{2}}{4} = 1$}{$\dfrac{{x}^{2}}{36} + \dfrac{{y}^{2}}{32} = 1$}{$\dfrac{{x}^{2}}{36} + \dfrac{{y}^{2}}{16} = 1$}{$\dfrac{{y}^{2}}{36} + \dfrac{{x}^{2}}{16} = 1$}
    \begin{jieda}
        以$(2, 0), (-2,0)$为焦点，长轴长为$12$的椭圆.
    \end{jieda}
    
    \item 下列命题中为真命题的是\dotfill{(\kh{B})}
    \xx{到 $x$ 轴的距离等于 $\sqrt{2}$ 的点所形成的直线是直线 $y = \sqrt{2}$}{两相交直线 $y = \dfrac{\sqrt{3}}{3}x$ 与 $y = \sqrt{3}x$ 的夹角的平分线的方程为 $y = x$}{若 $\bigtriangleup  {ABC}$ 的三顶点坐标分别为 $A\left( {1,1}\right)$ 、 $B\left( {3,1}\right)$ 、 $C\left( {1,3}\right)$ ,则 ${BC}$ 边上的中线的方程为 $y = x$}{与两定点 $A\left( {-1,0}\right) \text{ 、 }B\left( {1,0}\right)$ 连成的直线的夹角为 ${90}^{ \circ  }$ 的动点的轨迹方程为 ${x}^{2} + {y}^{2} = 1$}
    
    \item 到两坐标轴距离之和为 1 的点的轨迹是下图中的 \dotfill{(\kh{D})}
    \xx{\includegraphics[width=0.3\textwidth]{bodyxb/src/bo_d43fvo4601uc738mn8pg_1_239_1180_229_239_0.jpg}}{\includegraphics[width=0.3\textwidth]{bodyxb/src/bo_d43fvo4601uc738mn8pg_1_480_1179_234_239_0.jpg}}{\includegraphics[width=0.3\textwidth]{bodyxb/src/bo_d43fvo4601uc738mn8pg_1_729_1180_233_238_0.jpg}}{\includegraphics[width=0.3\textwidth]{bodyxb/src/bo_d43fvo4601uc738mn8pg_1_977_1182_235_236_0.jpg}}
    
    \item 到两坐标轴距离之差的绝对值等于 2 的点的轨迹是\dotfill{(\kh{D})}
    \xx{2 条直线}{4 条直线}{4 条射线}{8 条射线}
    
    \item 曲线 ${C} : F\left( {x,y}\right)  = 0$ 关于直线 $x - y - 2 = 0$ 对称的曲线 $C'$的方程是 \dotfill{(\kh{A})}
    \xx{$f\left( {y + 2,x - 2}\right)  = 0$}{$f\left( {y + 2,x}\right)  = 0$}{$f\left( {x - 2,y}\right)  = 0$}{$f\left( {y - 2,x + 2}\right)  = 0$}
    
    \item 若 $a\text{ 、 }b$ 为实数,且关于 $x$ 的一元二次方程 ${x}^{2} - {ax} + b = 0$ 的两个实数根为 $\sin \theta$ 和 $\cos \theta \left( {-\dfrac{\pi }{4} \leq  \theta  \leq  \dfrac{\pi }{4}}\right)$ ,则点 $P\left( {a,b}\right)$ 的轨迹方程是 \dotfill{(\kh{A})}
    \xx{${a}^{2} = 2\left( {b + \dfrac{1}{2}}\right) \left( {0 \leq  a \leq  \sqrt{2}}\right)$}{${a}^{2} = 2\left( {b + \dfrac{1}{2}}\right) \left( {-\sqrt{2} \leq  a \leq  \sqrt{2}}\right)$}{${a}^{2} = 2\left( {b + \dfrac{1}{2}}\right) \left( {-2 \leq  a \leq  2}\right)$}{${a}^{2} = 2\left( {b + \dfrac{1}{2}}\right) \left( {0 \leq  a \leq  2}\right)$}
    \begin{jieda}
        $\left\{\begin{array}{l}
            sin\theta + cos\theta = a \\
            sin\theta cos\theta = b 
        \end{array}\right.$，故$a^2 = 1+2b$. 因为$\theta \in \left[-\dfrac{\pi}{4}, \dfrac{\pi}{4}\right]$，所以$a \in [0, \sqrt{2}]$，$b \in \left[-\dfrac{1}{2}, \dfrac{1}{2}\right]$
    \end{jieda}


    \item 已知直线 $y = {kx}\left( {k > 0}\right)$ 交抛物线 $y = {x}^{2} - {2x} + 2$ 于不同的两点 $M$ 、 $P$ . 在线段 ${MP}$ 上是否存在动点 $Q$ ,使 $\dfrac{1}{\left| OM\right| }\text{ 、 }\dfrac{1}{\left| OQ\right| }\text{ 、 }\dfrac{1}{\left| OP\right| }$ 成等差数列? 若存在,求出点 $Q$ 的轨迹方程; 若不存在, 说明理由.
    \begin{jieda}
        \begin{enumerate}
            \item 方法1 

            设 $M\left( {{x}_{1},{y}_{1}}\right) ,P\left( {{x}_{2},{y}_{2}}\right) ,Q\left( {x,y}\right)$ .
            联立 $\left\{  \begin{array}{l} y = {kx} \\  y = {x}^{2} - {2x} + 2 \end{array}\right.$ ,化为：
            ${x}^{2} - \left( {2 + k}\right) x + 2 = 0$ ,
            由 $\Delta = {\left( 2 + k\right) }^{2} - 8 > 0$ , $k > 0$ ,解得：
            $k > 2\sqrt{2} - 2$ .
            
            根据韦达定理有$\left\{ \begin{array}{l}
                 {x}_{1} + {x}_{2} = 2 + k \\
                 {x}_{1}{x}_{2} = 2
            \end{array}\right.$.
            
            假设在 ${MP}$ 上存在动点 $Q$ ,使 $\dfrac{1}{\left| OM\right| }$ 、 $\dfrac{1}{\left| OQ\right| }$ 、
            $\dfrac{1}{\left| OP\right| }$ 成等差数列,
            则 $\dfrac{1}{\left| OM\right| } + \dfrac{1}{\left| OP\right| } = \dfrac{2}{\left| OQ\right| }$ ,即
            $\dfrac{\left| OQ\right| }{\left| OM\right| } + \dfrac{\left| OQ\right| }{\left| OP\right| } = 2$
            
            $\therefore \dfrac{x}{{x}_{1}} + \dfrac{x}{{x}_{2}} = 2$ ,化为：
            $x = \dfrac{2{x}_{1}{x}_{2}}{{x}_{1} + {x}_{2}} = \dfrac{2 \times  2}{2 + k} = \dfrac{4}{2 + k}.$
            $\therefore y = {kx} = \dfrac{4k}{2 + k}$ ,
            消去 $k$ 可得： ${2x} + y - 4 = 0$ .
            因此在 ${MP}$ 上存在动点 $Q$ ,使 $\dfrac{1}{\left| OM\right| }$ 、 $\dfrac{1}{\left| OQ\right| }$ 、
            $\dfrac{1}{\left| OP\right| }$ 成等差数列,点 $Q$ 的轨迹方程为：${2x} + y - 4 = 0(x \in(0,\sqrt{2})).$

            \item 方法2

            设直线参数方程为$\left\{\begin{array}{l}
                 x = tcos\theta  \\
                 y = tsin\theta
            \end{array}\right.$，$ \theta \in \left(0, \dfrac{\pi}{2}\right)$，代入抛物线的直角坐标方程，$tsin\theta = (tcos\theta)^2-2tcos\theta+2$，即$cos^2\theta \cdot t^2 - (sin\theta + 2cos\theta) t + 2 = 0$

            $\Delta > 0 \Longleftrightarrow (sin\theta + 2cos\theta)^2-8cos^2\theta > 0 \Longleftrightarrow (2+tan\theta)^2>8 \Longleftrightarrow tan\theta > 2\sqrt{2}-2$

            $|OM| = t_1, |OP| = t_2$，故$\dfrac{1}{\left| OM\right| }\text{ 、 }\dfrac{1}{\left| OQ\right| }\text{ 、 }\dfrac{1}{\left| OP\right| }$ 成等差数列即$\dfrac{1}{|OM|} + \dfrac{1}{|OP|} = \dfrac{1}{|OQ|}$，即$\dfrac{t_1+t_2}{t_1t_2} = \dfrac{2}{|OQ|}$，即$|OQ| = \dfrac{2t_1t_2}{t_1+t_2} = \dfrac{4}{sin\theta+2cos\theta}$，因此$\left\{\begin{array}{l}
                 x_Q = \dfrac{4cos\theta}{sin\theta+2cos\theta} = \dfrac{4}{2+tan\theta}  \\
                 y_Q = \dfrac{4sin\theta}{sin\theta+2cos\theta} = \dfrac{4tan\theta}{2+tan\theta} 
            \end{array}\right.$，故$2x_Q+y_Q = 4, x_Q \in (0, \sqrt{2})$.

            故点 $Q$ 的轨迹方程为：${2x} + y - 4 = 0(x \in(0,\sqrt{2})).$
        \end{enumerate}
        
    \end{jieda} 
    
    
    \item 如图,点 $A\text{ 、 }B$ 分别是椭圆 $\dfrac{{x}^{2}}{36} + \dfrac{{y}^{2}}{20} = 1$ 长轴的左、右端点,点 $F$ 是椭圆的右焦点,点 $P$ 在椭圆上,且位于 $x$ 轴上方, ${PA} \bot  {PF}$ .
    \begin{center}
    \includegraphics[width=0.3\textwidth]{bodyxb/src/bo_d43fvo4601uc738mn8pg_2_241_1102_360_278_0.jpg}
    \end{center}
    \begin{enumerate}
        \item 求点 $P$ 的坐标
        \item 设 $M$ 是椭圆长轴 ${AB}$ 上的一点, $M$ 到直线 ${AP}$ 的距离等于 $\left| {MB}\right|$ ,求椭圆上的点到点 $M$ 的距离 $d$ 的最小值.
    \end{enumerate}
    \begin{jieda}
        \begin{enumerate}
            \item 由题知$A\left( {-6,0}\right) ,B\left( {6,0}\right) ,F\left( {4,0}\right)$
            
            $\because {PA} \bot  {PF},\therefore \vv{AP} \cdot  \vv{FP} = 0$，即 ${x}^{2} + {2x} - {24} + {y}^{2} = 0$，和$\dfrac{{x}^{2}}{36} + \dfrac{{y}^{2}}{20} = 1$两式联立可得:$x = \dfrac{3}{2},y = \dfrac{5\sqrt{3}}{2}$，故点 $P$ 的坐标为 $\left( {\dfrac{3}{2},\dfrac{5\sqrt{3}}{2}}\right)$

            \item ${l}_{AP} : x - \sqrt{3}y + 6 = 0$ 
            、
            不妨设 $M\left( {m,0}\right) ,m \in  \left\lbrack  {-6,6}\right\rbrack$，由题意可得$\dfrac{\left| m + 6\right| }{\sqrt{1 + 3}} = 6 - m$，即 $\left| {m + 6}\right|  = {12} - {2m}$，因为 $m \in  \left\lbrack  {-6,6}\right\rbrack$，解得 $m = 2$
            
            设 $Q\left( {x,y}\right)$ 为椭圆任一点，则 $d = {QM} = \sqrt{{\left( x - 2\right) }^{2} + {y}^{2}} = \sqrt{{x}^{2} - {4x} + 4 + {20} - \dfrac{5{x}^{2}}{9}} = \sqrt{\dfrac{4{x}^{2}}{9} - {4x} + {24}} = \sqrt{\dfrac{4}{9}{\left( x - \dfrac{9}{2}\right) }^{2} + {15}}$ ,
            
            故当 $x = \dfrac{9}{2}$ 时 $d$ 有最小值为 $\sqrt{15}$ .
        \end{enumerate}
    \end{jieda}
    
\end{enumerate}

C 组：
\begin{enumerate}
    \item 已知曲线 $C : {x}^{2} + {y}^{2} = {2a}\left| x\right|  + {2y}$ ,下列说法正确的有\blkda{\ding{172}\ding{174}\ding{175}}.
    \begin{dingautolist}{172}
        \item 曲线 $C$ 关于 $y$ 轴对称
        \item 存在 $a$ ,使得曲线 $C$ 与坐标轴的交点个数为 3
        \item 曲线 $C$ 围成的区域面积 $S\left( a\right)$ 是关于 $a$ 的增函数
        \item 当 $a \neq  0$ 时,直线 $l : y = {4ax}$ 与曲线 $C$ 有且仅有 2 个交点
    \end{dingautolist}
    \begin{jieda}

    \begin{dingautolist}{172}
        \item 用$-x$代换$x$，方程依然成立，故正确
        \item 当 $x = 0$ 时, ${y}^{2} = {2y}$ ,所以曲线 $C$ 过点 $\left( {0,0}\right) ,\left( {0,2}\right)$
        
        当 $y = 0$ 时, ${x}^{2} = {2a}\left| x\right|$ ,即: ${\left| x\right| }^{2} = {2a}\left| x\right|$ ,所以曲线 $C$ 过点 $\left( {0,0}\right) ,\left( {\pm {2a},0}\right)$

        因此当 $a = 0$ 时,曲线 $C$ 与坐标轴有两个交点,当 $a \neq  0$ 时,曲线 $C$ 与坐标轴有四个交点，故错误.
        \item 当 $a = 0$ 时, $S = \pi$，如图所示:
        
        \begin{center}
        \includegraphics[width=0.2\textwidth]{bodyxb/src/bo_d43hsi3ef24c73d4tekg_6_242_1459_305_292_0.jpg}
        \end{center}

        当 $a \neq  0$ 时,曲线 $C$ 表示的两个圆的方程分别为:
        
        ${C}_{1} : {\left( x + a\right) }^{2} + {\left( y - 1\right) }^{2} = {a}^{2} + 1$ ,圆心为 $\left( {-a,1}\right)$ ,半径 $r = \sqrt{{a}^{2} + 1}$ ,
        
        ${C}_{2} : {\left( x - a\right) }^{2} + {\left( y - 1\right) }^{2} = {a}^{2} + 1$ ,圆心为 $\left( {a,1}\right)$ ,半径 $r = \sqrt{{a}^{2} + 1}$

        当 $a < 0$ 时,如图所示:

        \begin{center}
        \includegraphics[width=0.3\textwidth]{bodyxb/src/bo_d43hsi3ef24c73d4tekg_7_240_201_363_291_0.jpg}
        \end{center}
        
        由图可知,曲线 $C$ 恒过 $\left( {0,0}\right) ,\left( {0,2}\right)$ ,且两个圆的圆心纵坐标均为 1，所以曲线 $C$ 围成的区域面积 $S\left( a\right)$ 随 $a$ 的增大而增大

        当 $a > 0$ 时,如图所示:
        
        \begin{center}
        \includegraphics[width=0.3\textwidth]{bodyxb/src/bo_d43hsi3ef24c73d4tekg_7_240_744_362_290_0.jpg}
        \end{center}
        
        由图可知,曲线 $C$ 同样恒过 $\left( {0,0}\right) ,\left( {0,2}\right)$ ,且两个圆的圆心纵坐标还是 1，所以曲线 $C$ 围成的区域面积 $S\left( a\right)$ 随 $a$ 的增大而增大

        且 $a = 0$ 时,区域面积 $S\left( a\right)$ 介于 $a < 0$ 和 $a > 0$ 之间，所以曲线 $C$ 围成的区域面积 $S\left( a\right)$ 是关于 $a$ 的增函数，故正确

        \item 若 $a < 0$ ,联立 $\left\{  \begin{array}{l} y = {4ax} \\  {\left( x + a\right) }^{2} + {\left( y - 1\right) }^{2} = {a}^{2} + 1 \end{array}\right.$，解得: $x = 0$ 或 $x = \dfrac{6a}{1 + {16}{a}^{2}}$，即 $y = {4ax}$ 与曲线 $C$ 有且仅有 2 个交点;
        
        若 $a > 0$ ,联立 $\left\{  \begin{array}{l} y = {4ax} \\  {\left( x - a\right) }^{2} + {\left( y - 1\right) }^{2} = {a}^{2} + 1 \end{array}\right.$，解得: $x = 0$ 或 $x = \dfrac{10a}{1 + {16}{a}^{2}}$ ,
        
        即 $y = {4ax}$ 与曲线 $C$ 有且仅有 2 个交点，故正确
    \end{dingautolist}

    \end{jieda}
    
    \item 已知定圆 $M : {x}^{2} + {y}^{2} = 1$ 点 $A$ 是圆 $M$ 所在平面内一定点,点 $P$ 是圆 $M$ 上的动点, 若线段 ${PA}$ 的中垂线交直线 ${PM}$ 于点 $Q$ ,则点 $Q$ 的轨迹可能是:(1)椭圆；(2)双曲线； (3)抛物线；(4)圆；(5)直线；(6)一个点.其中所有可能的结果有\dotfill{(\kh{C})}
    \xx{2 个}{3 个}{4 个}{5 个}
    \begin{jieda}
    \begin{dingautolist}{172}
        \item 椭圆 (点 $A$ 在圆 $M$ 内, 不包括圆心)，如图:
        
        \begin{center}
        \includegraphics[width=0.3\textwidth]{bodyxb/src/bo_d43hsi3ef24c73d4tekg_8_243_746_356_332_0.jpg}
        \end{center}
        
        $\because \left| {AQ}\right|  = \left| {PQ}\right| ,\therefore \left| {AQ}\right|  + \left| {MQ}\right|  = \left| {PQ}\right|  + \left| {MQ}\right|  = \left| {MP}\right|  = 1$ ,即 $Q$ 点的轨迹是以 $A,M$ 为焦点的椭圆. 
        
        \item 双曲线(点 $A$ 在圆 $M$ 外)，如图:
        
        \begin{center}
        \includegraphics[width=0.3\textwidth]{bodyxb/src/bo_d43hsi3ef24c73d4tekg_8_242_1307_357_334_0.jpg}
        \end{center}
        
        $\because \left| {AQ}\right|  = \left| {PQ}\right| ,\therefore \left| {MQ}\right|  - \left| {AQ}\right|  = \left| {MQ}\right|  - \left| {PQ}\right|  = \left| {MP}\right|  = 1$ ,即 $Q$ 点的轨迹是以 $A,M$ 为焦点的双曲线.
        
        \item 圆(点 $A$ 恰为圆心)，如图:
        
        \begin{center}
        \includegraphics[width=0.3\textwidth]{bodyxb/src/bo_d43hsi3ef24c73d4tekg_9_241_197_359_337_0.jpg}
        \end{center}
        
        当 ${A}$ 与 $M$ 重合时, $Q$ 在 ${MP}$ 中点,所以 $Q$ 点的轨迹是以 $M$ 为圆心,半径为 $\dfrac{1}{2}$ 的圆.
        
        \item 一个点(点 $A$ 在圆 $M$ 上)，如图:
        
        \begin{center}
        \includegraphics[width=0.3\textwidth]{bodyxb/src/bo_d43hsi3ef24c73d4tekg_9_243_700_356_331_0.jpg}
        \end{center}
        
        当 ${A}$ 在圆上时, ${AP}$ 为弦,所以 $Q$ 一定与 $M$ 重合,所以 $Q$ 点的轨迹是点 $M$ .
    \end{dingautolist}
    \end{jieda}
    
    \item 已知 $P\text{ 、 }Q$ 是曲线 $\Gamma$ 上两点,若对于平面内任意不在曲线 $\Gamma$ 上的 $M$ 点,使得曲线 $\Gamma$ 上存在点 $P$ 和点 $Q$ 使得 ${MP} \bot  {MQ}$ ,则称曲线 $\Gamma$ 是 “ ${YY}$ 曲线”. 现有如下两个命题:  \ding{172} 任意椭圆都是 “ ${YY}$ 曲线”； \ding{173} 任意双曲线都是 “ ${YY}$ 曲线”. 则\dotfill{(\kh{D})}
    \xx{\ding{172} 成立 \ding{173} 成立}{\ding{172} 成立, \ding{173} 不成立}{\ding{172} 不成立, \ding{173} 成立}{\ding{172} 不成立, \ding{173} 不成立}
    \begin{jieda}
        对椭圆：在椭圆的蒙日圆$x^2+y^2 = a^2+b^2$外的点$M$必然不能找到对应的$P,Q$.
        
        对双曲线：若$a>b$，则在该双曲线的蒙日圆$x^2+y^2 = a^2-b^2$内的点$M$必然不能找到对应的$P,Q$.
    \end{jieda}
\end{enumerate}
\section*{简单的参数方程}

\begin{enumerate}
    \item 已知曲线 $C : \left\{  {\begin{array}{l} x = \cos \theta , \\  y = 2\sin \theta  \end{array}\;\text{ (参数 }\theta  \in  \mathbf{R}\text{ ) 经过点 }\left( {m,\dfrac{1}{2}}\right) }\right.$ ,则实数 $m =$ \blkda{$\pm \dfrac{\sqrt{15}}{4}$}.
    
    \item 已知定点 $A\left( {1,0}\right) ,F$ 是曲线 $\left\{  \begin{array}{l} x = 2\cos \theta , \\  y = 1 + \cos {2\theta } \end{array}\right.$ (参数 $\left. {\theta  \in  \mathbf{R}}\right)$ 的焦点,则 $\left| {AF}\right|  =$ \blkda{$\dfrac{\sqrt{5}}{2}$}.

    \item 已知直线 $\left\{  \begin{array}{l} x =  - 1 - {3t}, \\  y = 2 + {4t} \end{array}\right.$ 分别与双曲线 ${\left( y - 2\right) }^{2} - {x}^{2} = 1$ 相交于 $A\text{ 、 }B$ 两点,若点 $P$ 坐标为 $\left( {-1,2}\right)$ ,则 $\left| {PA}\right|  \cdot  \left| {PB}\right|  =$ \blkda{$\dfrac{50}{7}$}, $\left| {AB}\right|  =$ \blkda{$\dfrac{10\sqrt{23}}{7}$},弦 ${AB}$ 的中点 $M$ 与点 $P$ 的距离 $\left| {MP}\right|  =$ \blkda{$\dfrac{15}{7}$}.

    \item 已知实数 $x,y$ 满足方程 ${x}^{2} + {y}^{2} - {6x} - {4y} + {12} = 0$ ,则 $x + {2y}$ 的最大值是\blkda{$7+\sqrt{5}$}.

    \item 已知曲线 $C : {x}^{2} + {y}^{2} - {2y} = 0$ 上任意点 $\left( {x,y}\right)$ 的坐标都使不等式 $x + y + m \geq  0$ 恒成立, 则实数 $m$ 的取值范围是\blkda[2]{$[\sqrt{2}-1, +\infty)$}.

    \item 若点 $P\left( {x,y}\right)$ 在椭圆 $\dfrac{{x}^{2}}{4} + \dfrac{{y}^{2}}{9} = 1$ 上,则 ${2x} - y + 3$ 的最大值为\blkda{8},最小值为\blkda{-2}.

    \item 已知直线 $l$ 的参数方程为 $\left\{  \begin{array}{l} x = 3 + {4t}, \\  y =  - 4 + {3t} \end{array}\right.$ ( $t$ 为参数),则下列说法中错误的是\dotfill{(\kh{C})}
    \xx{直线 $l$ 的斜率是 $\dfrac{3}{4}$}{直线 $l$ 过点 $\left( {3, - 4}\right)$}{当 $t = 1$ 时,直线 $l$ 上所对应的点到 $\left( {3, - 4}\right)$ 的距离等于 1}{直线 $l$ 不经过第二象限}
    
    \begin{center}
    \includegraphics[width=0.3\textwidth]{bodyxb/src/bo_d43fvr3ef24c73d4s7ig_0_1135_1589_377_310_0.jpg}
    \end{center}
    
    \item 如图,长为 6 的线段 ${AB}$ 的两端 $A$ 、 $B$ 分别在 $x$ 轴、 $y$ 轴上滑动, $P$ 为线段 ${AB}$ 上的定点,且 ${AP} : {PB} = 2 : 1$ ,设 $\angle {OAP} = \theta$ 为参数,则点 $P$ 的轨迹的参数方程为\blkda[3]{$\left\{  \begin{array}{l} x = 2\cos \theta \\  y = 4\sin \theta  \end{array}\right.$}.
    \begin{jieda}
        考虑$P$在第一象限的情况，其他象限可以根据对称性得到
        
        设$P(x_0, y_0)$，则$AB: \left\{\begin{array}{l}
             x = x_0+tcos(\pi-\theta) \\
             y = y_0+tsin(\pi-\theta)
        \end{array}\right.$，代入到$xy = 0$有$(x_0-tcos\theta)(y_0+tsin\theta) = 0$，即$(cos\theta sin\theta)t^2-(x_0sin\theta-y_0cos\theta)t-x_0y_0 = 0$

        根据题意可知方程的解为$t_A = -4$，$t_B = 2$，故根据韦达定理有$\left\{\begin{array}{l}
             -2 = \dfrac{x_0sin\theta-y_0cos\theta}{sin\theta cos\theta} \\
             -8 = -\dfrac{x_0y_0}{sin\theta cos\theta} 
        \end{array}\right.$，解得$\left\{  \begin{array}{l} x_0 = 2\cos \theta \\  y_0 = 4\sin \theta  \end{array}\right.$，根据对称性可知轨迹的参数方程为$\left\{  \begin{array}{l} x = 2\cos \theta \\  y = 4\sin \theta  \end{array}\right.$.
    \end{jieda}
    
    \item 参数方程 $\left\{  \begin{array}{l} x = {\cos }^{2}\theta \\  y = \cos {2\theta } \end{array}\right.$ ( $\theta$ 为参数) 所表示的曲线是\dotfill{(\kh{C})}
    \xx{圆}{直线}{线段}{抛物线}
    
    \item 参数方程 $\left\{  {\begin{array}{l} x = t\cos a \\  y = t\sin a \end{array}\left( {t \neq  0,a \in  \lbrack 0,{2\pi }}\right) }\right.$ )表示的曲线是\dotfill{(\kh{D})}
    \xx{直线}{圆}{可能是直线也可能是圆}{以上都不正确}
    
    \item 已知曲线 $C : \left\{  {\begin{array}{l} x = 2\cos \theta , \\  y = 2\sin \theta  \end{array}\;\text{ ( }\theta }\right.$ 为参数)和直线 $l : \left\{  \begin{array}{l} x = t, \\  y = t + b \end{array}\right.$ ( $t$ 为参数, $b$ 为实数),若曲线 $C$ 上恰有 3 个点到直线 $l$ 的距离等于 1,则 $b =$ \dotfill{(\kh{D})}
    \xx{$\sqrt{2}$}{$- \sqrt{2}$}{0}{$\pm  \sqrt{2}$}
    
    \item 对于参数方程 $\left\{  {\begin{array}{l} x = {at} + \lambda \cos \theta \\  y = {bt} + \lambda \sin \theta  \end{array}\left( {a,b,\lambda \text{ 均不为零, }0 \leq  \theta  \leq  {2\pi }}\right) }\right.$ ,指出分别在下列情况下对应的曲线:
    \begin{enumerate*}
        \item $t$ 为参数
        \item $\lambda$ 为参数
        \item $\theta$ 为参数
    \end{enumerate*}
    \begin{jieda}
        \begin{enumerate}
            \item 直线, ${bx} - {ay} = {\lambda b}\cos \theta  - {\lambda a}\sin \theta$
            \item 直线, $\left( {x - {at}}\right) \sin \theta  - \left( {y - {bt}}\right) \cos \theta  = 0$
            \item 圆, ${\left( x - at\right) }^{2} + {\left( y - bt\right) }^{2} = {\lambda }^{2}$
        \end{enumerate}
    \end{jieda}
    
    \item 已知直线 $l$ 过点 $P\left( {0,\dfrac{1}{2}}\right)$ ,且与抛物线 $C : y = \dfrac{1}{2}{x}^{2}$ 相交于 $A$ 、 $B$ 两点,求线段 ${AB}$ 的中点的轨迹方程.
    \begin{jieda}
        \begin{enumerate}
            \item 方法1
            
            设$A(2t_1, 2t_1^2)$，$B(2t_2, 2t_2^2)$，则中点$M(t_1+t_2, t_1^2+t_2^2)$
    
            直线$AB: y = \dfrac{(2t_2^2-2t_1^2)}{2t_2-2t_1}x+\dfrac{1}{2}$，即$y = (t_1+t_2)x+\dfrac{1}{2}$，故$y_M = x_M \cdot x_M + \dfrac{1}{2}$，因此中点轨迹为$y = x^2+\dfrac{1}{2}$.

            \item 方法2

            设直线的参数方程为$\left\{\begin{array}{l}
                 x = tcos\alpha \\
                 y = \dfrac{1}{2} + t sin\alpha
            \end{array}\right.$，代入到抛物线的直角坐标方程中整理得$cos^2\alpha \cdot t^2 - 2sin\alpha \cdot t - 1 = 0$，$\Delta > 0$

            故$\left\{\begin{array}{l}
                 x_M = \dfrac{t_1+t_2}{2} \cdot cos\alpha = \dfrac{sin\alpha}{cos\alpha} \\
                 y_M = \dfrac{1}{2} + \dfrac{t_1+t_2}{2} \cdot sin\alpha = \dfrac{1}{2} + \dfrac{sin^2\alpha}{cos^2\alpha}
            \end{array}\right.$，故$y_M = x_M^2+\dfrac{1}{2}$，因此中点轨迹为$y = x^2+\dfrac{1}{2}$
        \end{enumerate}
        
    \end{jieda}
    
    \item 过点 $M\left( {-2,4}\right)$ 引倾斜角为 $\dfrac{3\pi }{4}$ 的直线交抛物线 ${y}^{2} = {2px}\left( {p > 0}\right)$ 于 $P\text{ 、 }Q$ 两点且 $\left| {MP}\right| \text{ 、 }\left| {PQ}\right| \text{ 、 }\left| {MQ}\right|$ 依次成等比数列,求 $p$ 的值.
    \begin{jieda}
        \begin{enumerate}
            \item 方法1

            直线$PQ$的参数方程：$\left\{\begin{array}{l}
                 x = -\dfrac{\sqrt{2}t}{2} - 2 \\
                 y = \dfrac{\sqrt{2}t}{2} + 4
            \end{array}\right.$

            代入抛物线方程有$\left(\dfrac{\sqrt{2}t}{2} + 4\right)^2 = 2p\left(-\dfrac{\sqrt{2}t}{2} - 2\right)$，整理得$2{t}^{2} + 4\sqrt{2}\left( {4 + p}\right) t + {16}\left( {4 + p}\right)  = 0$

            根据韦达定理有$\left\{\begin{array}{l}
                 t_1+t_2 = -2\sqrt{2}(4+p)  \\
                 t_1t_2 = 8(4+p)
            \end{array}\right.$

            $|MP| = -t_1, |MQ| = -t_2, |PQ| = |t_2-t_1|$，故$t_1t_2 = (t_2-t_1)^2$，即$t_1t_2 = (t_1+t_2)^2-4t_1t_2$，即$40(4+p) = 8(4+p)^2$，解得$p = 1$
            
            经检验直线和抛物线有两个交点.
            
            \item 方法2

            设$P(2pt_1^2, 2pt_1)$，$Q(2pt_2^2, 2pt_2)$
    
            根据题意有$PQ: y = -x+2$，故$\left\{\begin{array}{l}
                 2pt_1 = -2pt_1^2+2  \\
                 2pt_2 = -2pt_2^2+2
            \end{array}\right.$
    
            两式做差可得$2p(t_1-t_2) = -2p(t_1^2-t_2^2)$，即$t_1+t_2=-1$
            
            两式做和可得$2p(t_1+t_2) = -2p(t_1^2+t_2^2)$，整理得$t_1t_2=-\dfrac{1}{p}$
    
            $\left| {MP}\right| \text{ 、 }\left| {PQ}\right| \text{ 、 }\left| {MQ}\right|$ 依次成等比数列即其纵坐标差值成等比数列，即$|2pt_1-4|\cdot|2pt_2-4| = |2pt_1-2pt_2|^2$，即$|4p^2t_1t_2-8p(t_1+t_2)+16| = 4p^2[(t_1+t_2)^2-4t_1t_2]$，即$4p+16 = 4p^2+16p$，解得$p = 1$
            
            经检验直线和抛物线有两个交点.
        \end{enumerate}
    \end{jieda}
    
    \item 过定点 $P\left( {0, - \sqrt{5}}\right)$ 的直线 $l$ 分别与椭圆 $\dfrac{{x}^{2}}{4} + {y}^{2} = 1$ 交于 $A\text{ 、 }B$ 两点 (包括 $A\text{ 、 }B$ 两点重合).
    \begin{enumerate}
        \item 求直线 $l$ 的斜率 $k$ 与倾斜角 $\alpha$ 的取值范围
        \item 求 $\left| {PA}\right|  \cdot  \left| {PB}\right|$ 的最大值与最小值
    \end{enumerate}
    \begin{jieda}
        \begin{enumerate}
            \item 设$l: y = kx-\sqrt{5}$，联立椭圆和直线，消去$y$可得$\left(\dfrac{1}{4}+k^2\right)x^2-2\sqrt{5}kx+4=0$，$\Delta = 4k^2-4 \geq 0$，故$k^2 \geq 1$，即$k \in (-\infty, -1] \cup [1, +\infty)$，$\alpha \in \left[\dfrac{\pi}{4}, \dfrac{3\pi}{4}\right]$
            \item 直线的参数方程：$\left\{\begin{array}{l}
                 x = tcos\alpha \\
                 y = tsin\alpha-\sqrt{5}
            \end{array}\right.$，代入椭圆方程有$\dfrac{t^2cos^2\alpha}{4}+(tsin\alpha-\sqrt{5})^2 = 1$，化简得$\dfrac{1 + 3sin^{2}\alpha}{4}{t}^{2} - 2\sqrt{5}sin\alpha{t} + 4 = 0$

            $|PA| \cdot |PB| = t_1t_2 = \dfrac{16}{1+3sin^2\alpha}$，因为$\alpha \in \left[\dfrac{\pi}{4}, \dfrac{3\pi}{4}\right]$，所以$sin\alpha \in \left[\dfrac{\sqrt{2}}{2}, 1\right]$，故$|PA| \cdot |PB| \in \left[4, \dfrac{32}{5}\right]$，故最大值为$\dfrac{32}{5}$（$\alpha = \dfrac{\pi}{2}$），最小值为$4$（$\alpha = \dfrac{\pi}{4}$或$\dfrac{3\pi}{4}$）.
        \end{enumerate}
    \end{jieda}
\end{enumerate}

\section*{极坐标系与极坐标方程}
\begin{enumerate}
    \item 把以下坐标系中的点化为直角坐标
    
    \begin{enumerate*}
        \item $A\left( {3,\dfrac{\pi }{2}}\right)$\blkda{$\left( {0,3}\right)$}
        \item $B\left( {4,\pi }\right)$ \blkda{$(-4,0)$}
        \item $C\left( {2,\dfrac{\pi }{6}}\right)$ \blkda{$\left( {\sqrt{3},1}\right)$}
        \item $D\left( {-2,\dfrac{\pi }{3}}\right)$ \blkda{$\left( {-1, - \sqrt{3}}\right)$}
    \end{enumerate*}
    
    \item 把以下直角坐标系中的点化为极坐标 $\left( {\rho  > 0,0 \leq  \theta  < {2\pi }}\right)$
    
    \begin{enumerate*}
        \item $A\left( {-1,1}\right)$ \blkda{$\left( {\sqrt{2},\dfrac{3}{4}\pi }\right)$}
        \item $B\left( {3, - \sqrt{3}}\right)$ \blkda{$\left( {2\sqrt{3},\dfrac{11}{6}\pi }\right)$}
        \item $C\left( {0, - 1}\right)$ \blkda{$\left( {1,\dfrac{3}{2}\pi }\right)$}
    \end{enumerate*}
    
    \item 若点 ${A}$ 在极坐标系中的坐标为 $\left( {\rho ,\theta }\right)$ ,则它在直角坐标系中的坐标为\blkda[2]{$(\rho cos\theta, \rho sin\theta)$}.

    \item 极坐标系中一点 $A$ 的极坐标为 $\left( {3, - 4}\right)$ ,则点 $A$ 到极点的距离为\blkda{3}.

    \item 极坐标系中,点 $\left( {2,\dfrac{\pi }{6}}\right)$ 和 $\left( {2, - \dfrac{\pi }{6}}\right)$ 之间的距离为\blkda{2}.

    \item 极坐标系中,点 $\left( {-1,\dfrac{\pi }{3}}\right)$ 与点 $\left( {1,\theta }\right) ,\theta  \in  \lbrack 0,{2\pi })$ 表示同一个点,则 $\theta  =$ \blkda{$\dfrac{4\pi}{3}$}.

    \item 已知椭圆的极坐标方程为 $\rho  = \dfrac{3}{2 - \cos \theta }$ ,则椭圆的焦距是\blkda{2}；长轴长是\blkda{4}.

    \item 极坐标系中,点 $\left( {\rho ,\theta }\right)$ 与点 $\left( {-\rho ,\pi  + \theta }\right)$ 的关系是
    \dotfill{(\kh{C})}
    \xx{关于极轴对称}{关于极点对称}{重合}{以上都不对}
    
    \item 极坐标系中 $A\left( {-3,\dfrac{\pi }{6}}\right) ,B\left( {3,\dfrac{\pi }{2}}\right) ,{O}$ 为极点,求 $\bigtriangleup {AOB}$ 的面积.
    \begin{jieda}
        $A\left(3, \dfrac{7\pi}{6}\right)$，则$S = \dfrac{1}{2} \times 3 \times 3 \times sin\dfrac{2\pi}{3} = \dfrac{9\sqrt{3}}{4}$
    \end{jieda}

    \item 分别画出下列极坐标方程和直角坐标方程的图形:
    \begin{enumerate}
        \item 极坐标方程 $\rho  = 2$ 和直角坐标方程 $x = 2$
        \item 极坐标方 $\theta  = \dfrac{\pi }{4}$ 和直角坐标方程 $x = \dfrac{\pi }{4}$
    \end{enumerate}
    
    \item 在平面直角坐标系 ${xOy}$ 中,曲线 $C$ 的参数方程为 $\left\{  {\begin{array}{l} x = 4 + 2\cos \varphi \\  y = 2\sin \varphi  \end{array}\left( {\varphi  \in  \lbrack 0,{2\pi })}\right) }\right.$ . 在以坐标原点 $O$ 为极点, $x$ 轴的正半轴为极轴的极坐标系中,直线 $l$ 的方程为 $\rho \cos \left( {\theta  + \dfrac{\pi }{4}}\right)  = \sqrt{2}.$
    \begin{enumerate}
        \item 求曲线 $C$ 的普通方程和直线 $l$ 的直角坐标方程
        \item 已知点 $P\left( {0, - 2}\right)$ ,直线 $l$ 与曲线 $C$ 分别交于 $A$ , $B$ 两点,点 $M$ 是 ${AB}$ 的中点,求 ${PM}$ 的长
    \end{enumerate}
    \begin{jieda}
        \begin{enumerate}
            \item $C : {\left( x - 4\right) }^{2} + {y}^{2} = 4,l : y - x + 2 = 0$
            \item $3\sqrt{2}$
        \end{enumerate}
    \end{jieda}
    
\end{enumerate}